% Use the base layout without conference or publication branding.
\PassOptionsToPackage{capitalise,noabbrev}{cleveref}
\PassOptionsToPackage{hypertexnames=false}{hyperref}
\documentclass[12pt,cleveref]{jmlr}

\usepackage{times}
\usepackage{mathtools}
\usepackage{microtype}
\usepackage{tikz}
\usepackage{wrapfig}
\setlength{\marginparwidth}{1in}
\setlength{\marginparsep}{0.15in}
\usepackage{todonotes}
\hypersetup{linkcolor=blue!55!black,citecolor=blue!55!black,urlcolor=blue!55!black}

% Keep named proofs while using the class's proof typography and QED mark.
\renewenvironment{proof}[1][\proofname]
    {\par\noindent{\bfseries\upshape #1\ }}
    {\jmlrQED}
\usepackage{thm-restate}
\usepackage{autonum} % load last

% Restore theorem link targets omitted by the class on recent LaTeX kernels.
\makeatletter
\let\ps@jmlrtps\ps@plain
\patchcmd{\@jmlrmaketitle}
    {\jmlrpreauthor \ignorespaces\@author \jmlrpostauthor}
    {\ifx\@author\@empty\vskip\aftermaketitskip\else\jmlrpreauthor \ignorespaces\@author \jmlrpostauthor\fi}
    {}{\errmessage{Title layout patch failed}}
\ifdefined\MakeLinkTarget
    \patchcmd{\@begintheorem}{\hskip\labelsep}{\MakeLinkTarget{\@currentcounter}\hskip\labelsep}{}{\errmessage{Theorem hyperlink patch failed}}
    \patchcmd{\@opargbegintheorem}{\hskip\labelsep}{\MakeLinkTarget{\@currentcounter}\hskip\labelsep}{}{\errmessage{Named theorem hyperlink patch failed}}
\fi
\makeatother

\crefname{equation}{Eq.}{Eqs.}
\Crefname{equation}{Eq.}{Eqs.}
\crefname{equality}{Eq.}{Eqs.}
\Crefname{equality}{Eq.}{Eqs.}
\crefname{inequality}{Ineq.}{Ineqs.}
\Crefname{inequality}{Ineq.}{Ineqs.}
\creflabelformat{equality}{\textup{(#2#1#3)}}
\creflabelformat{inequality}{\textup{(#2#1#3)}}
\crefrangelabelformat{equality}{\textup{(#3#1#4)--(#5#2#6)}}
\crefrangelabelformat{inequality}{\textup{(#3#1#4)--(#5#2#6)}}
\newcommand{\eqlabel}[1]{\label[equality]{#1}}
\newcommand{\ineqlabel}[1]{\label[inequality]{#1}}
\newcommand{\proofpart}[2]{%
    \par\addvspace{1.2\baselineskip}
    \noindent
    \textcolor{blue!55!black}{\rule[-0.2ex]{2pt}{1.25em}}%
    \hspace{0.6em}%
    {\sffamily\bfseries\color{blue!55!black}Part #1}%
    \hspace{0.8em}%
    {\sffamily\bfseries #2}%
    \par\nobreak\smallskip
    \noindent\ignorespaces
}
\newcommand{\proofstepinline}[1]{%
    \par\addvspace{0.6\baselineskip}
    \noindent
    {\normalfont\normalsize\bfseries\color{blue!55!black}Step #1.}\enspace\ignorespaces
}
\newcommand{\reason}[1]{\texttt{(#1)}}
\newcommand{\explan}[2]{\stackrel{\reason{#2}}{#1}}
\newcommand{\R}{\mathbb R}
\newcommand{\E}{\mathbb E}
\newcommand{\Pp}{\mathbb P}
\newcommand{\K}{\mathcal K}
\newcommand{\cP}{\mathcal P}
\newcommand{\F}{\mathcal F}
\newcommand{\N}{\mathcal N}
\newcommand{\T}{\mathcal T}
\newcommand{\tr}{\operatorname{tr}}
\newcommand{\Lip}{\operatorname{Lip}}
\newcommand{\op}{\mathrm{op}}
\newcommand{\argmin}{\operatornamewithlimits{arg\,min}}
\newcommand{\1}{\mathbf 1}
\newcommand{\eps}{\varepsilon}
\renewcommand{\epsilon}{\varepsilon}
\newcommand{\ip}[2]{\left\langle #1,#2\right\rangle}
\newcommand{\norm}[1]{\left\lVert #1\right\rVert}
\newcommand{\normm}[2]{\left\lVert #1\right\rVert_{#2}}
\newcommand{\alireza}[1]{\todo[color=blue!20,size=\tiny]{Alireza: #1}}
\newcommand{\kevin}[1]{\todo[color=green!20,size=\tiny]{Kevin: #1}}
\newcommand{\lillian}[1]{\todo[color=purple!20,size=\tiny]{Lillian: #1}}

\title[Direct Quadratic Estimation for Bandit Convex Optimization]
{Direct Quadratic Estimation for Stochastic Bandit Convex Optimization}
\author{}

\begin{document}
\maketitle

\begin{abstract}
    We consider stationary stochastic bandit convex optimization with a bounded
    convex loss and subgaussian observation noise.  Gaussian online Newton step
    methods are usually analyzed through an optimistic Gaussian surrogate whose
    derivatives are estimated using a likelihood-ratio correction.  We show that
    neither device is needed in the stationary problem.  The algorithm estimates
    the first two Gaussian Hermite coefficients of a queryable convex extension
    directly from the observed loss and applies online Newton step to the resulting
    (damped) quadratic estimates, without projecting the mean.
    For a rounded convex body $\K$, the regret is
    \begin{align}
        O\!\left(d^2 M\sqrt n\,
        \log^8\!\frac{nd}{\delta}\right)
    \end{align}
    with probability at least $1-\delta$, where $M$ is $\tilde{O}(d^{-1/2})$
    when $\frac1d \K$ is in L\"owner position,
    which recovers the stationary rate of \citet{fokkema2024online}
    of $\widetilde O(d^{3/2}\sqrt n)$.
\end{abstract}

\section{Introduction}

The learner in bandit convex optimization chooses a point in a convex body
and observes only a noisy value of the loss at that point.  A useful way to
extract second-order information is to sample from a Gaussian and estimate the
first two coefficients in the corresponding Hermite expansion.  The first
coefficient behaves like a gradient and the second like an average Hessian, so
the natural update is online Newton step.

There is one difficulty with this direct idea.  The online Newton step
potential contains the curvature of the quadratic at the comparator, while
convexity controls regret through a Bregman divergence of the original loss.
These quantities need not be comparable globally.  The usual Gaussian ONS
analysis resolves this by first replacing the loss with an optimistic
Gaussian surrogate and then estimating its derivatives with a change of
measure; see \citet{lattimore2023second} and Chapters~10--12 of
\citet{lattimore2026bandit}.  Here we take a
different route for the stationary stochastic problem.  We compare the direct
Hermite curvature with the Bregman divergence on an ellipsoidal region around
the Gaussian mean.  The comparison loses only a small part of the usual
log-determinant trace budget.

The resulting algorithm has one Gaussian sample and one function evaluation
per round, with no projection of the mean.
The proof still has to address two features of bandit feedback:
the one-sample Hessian estimate is indefinite, and the queryable extension
amplifies noise outside the constraint set.
We handle the first by comparing the random covariance with its
predictable counterpart, and the second by bounding the extension value at the means and closing
a feedback-moment budget through extension regret.

The paper is organized as follows.  \Cref{sec:setting} introduces the problem,
the extension and the algorithm.
The stability and regret analysis, including the trace, moment and
concentration estimates, appears in Appendix~\ref{sec:stability}, with the
supporting proofs in the remaining appendices.

\section{Problem formulation and algorithm}
\label{sec:setting}

For $r>0$, let $\mathbb{B}_r=\{x\in\R^d:\norm{x}\leq r\}$
and $\mathbb{S}_r=\{x\in\R^d:\norm{x}=r\}$ be the Euclidean ball and sphere of radius $r$.
Let $\K\subset\R^d$ be a convex body.  An affine preprocessing step puts the body
in the position $\mathbb{B}_1\subseteq\K\subseteq \mathbb{B}_{2d}$.
The loss $f:\K\to[0,1]$ is convex and does not change with time.  In round
$t$, the learner chooses $A_t\in\K$ and observes $f(A_t)+\xi_t$
where $\xi_t$ is the noise in round $t$.
Let $\zeta_t$ be the randomness of the algorithm in round $t$ which it uses to make decisions, and let
$\F_t=\sigma(\zeta_1,\xi_1,\ldots,\zeta_t,\xi_t)$ be the filtration at the end of round $t$,
and $\E_t[\cdot]=\E[\cdot\mid\F_{t-1}]$ and $\Pp_t(\cdot)=\Pp(\cdot\mid\F_{t-1})$ be the conditional expectation and probability given the history.
The noise is conditionally subgaussian, i.e., for all $t$,
\begin{align}
    \E_t[\xi_t|\zeta_t]=0,
    \qquad
    \E_t[\exp(\xi_t^2)|\zeta_t]\leq2.
    \label{eq:noise}
\end{align}
The quantity of interest is the realized (\textit{pseudo-})regret
\begin{align}
    \operatorname{Reg}_n
    =\sum_{t=1}^n\left(f(A_t)-\inf_{x\in\K}f(x)\right).
    \eqlabel{eq:regret}
\end{align}

Let $\pi(x)=\inf\{r>0:x\in r\K\}$ be the Minkowski functional of $\K$ and let
$U$ be uniform on the Euclidean unit sphere.
The only geometric quantity in the result is
\begin{align}
    M=\max\left\{\frac1{\sqrt d},\E[\pi(U)]\right\}.
    \eqlabel{eq:M}
\end{align}
The rounding  $\mathbb{B}_1\subseteq\K$ implies that $\pi$ is $1$-Lipschitz.
Write $\gamma_{\mu,\Sigma}$ for the density of the multivariate normal distribution
$\N(\mu,\Sigma)$, for $\mu\in\R^d$ and $\Sigma\in\R^{d\times d}$ symmetric positive semidefinite.
We write $Dh$ and $D^2h$ for the first and second distributional derivatives of $h$;
see Appendix~\ref{app:distributional-derivatives} for their definitions.

The loss function $f$ has the restriction of being only defined on the convex body $\K$.
To overcome this limitation, we construct a queryable convex extension of $f$ to the entire space $\R^d$.
For some $\eps\in(0,1/2]$ and write $\K_\eps=(1-\eps)\K$ and
\begin{align}
    \pi_\eps(x)
     & =\max\left\{1,\frac{\pi(x)}{1-\eps}\right\}
    =\begin{cases}
         1               & \text{if }x\in\K_\eps, \\
         \pi(x)/(1-\eps) & \text{otherwise}\,,
     \end{cases}
    \quad\text{and}\quad v(x)=\pi_\eps(x)-1\,.
    \eqlabel{eq:extension-notation}
\end{align}
Then, we use the queryable extension from
\citet[Eq.~(10.6)]{lattimore2026bandit} which is defined for all $x\in\R^d$ by
\begin{align}
    e(x)
    =\pi_\eps(x)f\!\left(\frac{x}{\pi_\eps(x)}\right)
    +\frac{2}{\eps}v(x)\,,
    \quad\text{and satisfies }\,\,
    \Lip(e)\leq\frac{8}{\eps}\,.
    \label{eq:extension-lipschitz}
\end{align}
Moreover, the extension $e$ is convex and agrees with $f$ on $\K_\eps$.
These properties are established in \cref{app:extension}.

% \subsection{Gaussian $L^2$ projection and Hermite basis}
% \label{sec:hermite-projection}
% The Gaussian space $L^2(\N(\mu,\Sigma))$ consists of the functions that are
% square-integrable under the law of $X \sim \N(\mu, \Sigma)$, i.e.
% $L^2(\N(\mu,\Sigma))
%     =
%     \left\{
%     r:\R^d\to\R:
%     \E[r(X)^2]<\infty
%     \right\}$.
% Since $e$ is Lipschitz, we have
% \begin{align}
%     \E\left[
%         e(X)^2
%         \right]
%     \leq
%     \E\left[
%         (e(0)+\Lip(e)\norm{X})^2
%         \right]
%     \leq
%     2e(0)^2 + 2\Lip(e)^2\E[\norm{X}^2] < \infty,
% \end{align}
% which implies that $e\in L^2(\N(\mu,\Sigma))$.
% Next, define
% \begin{align}
%     \mathcal P_{\mu,\Sigma}e(x)
%     =
%     \E[e(X)]
%     +\ip{\bar{g}}{x-\mu}
%     +\frac12
%     \left(
%     \|x-\mu\|^2_{\bar{H}}
%     -\tr(\bar{H}\Sigma)
%     \right).
%     \label{eq:hermite-best-approximation}
% \end{align}
\section{The algorithm}
The algorithm maintains a Gaussian sampling distribution and uses one observation per round to update its mean and covariance.
In round $t$, it samples $X_t\sim\N(\mu_t,\Sigma_t)$, plays the radial contraction
$A_t=X_t/\pi_\eps(X_t)\in\K$, and observes $f(A_t)+\xi_t$.
The extension defined above lets us turn this observation into feedback for $e$:
\begin{align}
    Y_t
     & =\pi_\eps(X_t)\bigl(f(A_t)+\xi_t\bigr)
    +\frac2\eps\bigl(\pi_\eps(X_t)-1\bigr)\nonumber \\
     & =e(X_t)+\pi_\eps(X_t)\xi_t,
    \eqlabel{eq:Y}
\end{align}
so $\E_t[Y_t\mid X_t]=e(X_t)$.
From the pair $(X_t,Y_t)$, we compute $Z_t =\Sigma_t^{-1/2}(X_t-\mu_t)$,
\begin{align}
    g_t =Y_t\Sigma_t^{-1/2}Z_t\qquad \text{and}\qquad
    H_t =Y_t\Sigma_t^{-1/2}(Z_tZ_t^\top-I_d)\Sigma_t^{-1/2}.
    \eqlabel{eq:sampled-hermite-coefficients}
\end{align}
The coefficients $g_t$ and $H_t$ define a stochastic quadratic approximation to a Gaussian smoothing of $e$, as explained in \cref{sec:quadratic-approximation}.
With learning rate $\eta>0$ and curvature scaling $\lambda\in(0,1)$, we use these coefficients to update the Gaussian parameters by unprojected online Newton steps (ONS), as specified in \cref{alg:direct-hermite}.
The optimization interpretation of these updates is developed in \cref{sec:population-ons}.
The covariance update requires $\Sigma_t^{-1}+\lambda\eta H_t\succ0$; Appendix~\ref{sec:stability} controls this condition in the presence of estimation error.

\begingroup
\RestyleAlgo{ruled}
\DontPrintSemicolon
\SetAlgoNoLine
\SetAlgoLongEnd
\SetAlFnt{\small}
\SetKwInput{KwIn}{Input}
\begin{algorithm2e}[!t]
    \caption{Direct Hermite--ONS algorithm}
    \label{alg:direct-hermite}
    \KwIn{$\eta,\lambda,\sigma^2,\eps$}
    $\mu_1=0$ and $\Sigma_1=\sigma^2I_d$\;
    \BlankLine
    \For{$t=1,\ldots,n$}{
        Sample $X_t\sim\N(\mu_t,\Sigma_t)$\;
        Play $A_t=X_t/\pi_\eps(X_t)$ and observe $f(A_t)+\xi_t$\;
        Compute $Y_t$ from the observed feedback using \cref{eq:Y}\;
        Compute $Z_t$, $g_t$ and $H_t$ using \cref{eq:sampled-hermite-coefficients}\;
        Set $\Sigma_{t+1}^{-1}=\Sigma_t^{-1}+\lambda\eta H_t$\;
        Set $\mu_{t+1}=\mu_t-\eta\Sigma_{t+1}g_t$\;
    }
\end{algorithm2e}
\endgroup
\FloatBarrier

\subsection{Main result}

\begin{theorem}
    \label{thm:main}
    For $\delta\in(0,1/2)$, let
    \begin{align}
        \ell    & =128\log\frac{128nd}{\delta},
                &
        \sigma  & =\frac1{M\sqrt{2d}},
                &
        \lambda & =\frac{1}{1000d^{3/2}M\ell^4},
                &
        \eta    & =\frac{8dM}{\sqrt n},
                &
        \eps_0  & =\frac{100d^2M\ell^8}{\sqrt n}.
        \eqlabel{eq:parameters-one}
        \eqlabel{eq:parameters-two}
    \end{align}
    If $\eps_0\leq1/2$, run \cref{alg:direct-hermite} with $\eps=\eps_0$.
    Otherwise play any fixed point in $\K$.  In the setting above,
    \begin{align}
        \operatorname{Reg}_n
        \leq 200d^2M\ell^8\sqrt n
        \ineqlabel{eq:main-regret}
    \end{align}
    with probability at least $1-\delta$.
\end{theorem}

\begin{corollary}
    \label{cor:loewner}
    If $\K/d$ is in L\"owner position, then
    \begin{align}
        \operatorname{Reg}_n
        =O\!\left(
        d^{3/2}\sqrt{\log(ed)}\,\ell^8\sqrt n
        \right)
        =\widetilde O(d^{3/2}\sqrt n)
    \end{align}
    with probability at least $1-\delta$.
\end{corollary}

\begin{proof}
    In L\"owner position,
    $\E[\pi(U)]=O(\sqrt{\log(ed)/d})$; see
    \citet[Table~10.1(a)]{lattimore2026bandit}.  Substitute this bound in
    \cref{eq:main-regret}.
\end{proof}

\section{Idealized Behavior}
We will now describe the idealized behavior of the algorithm, assuming access to exact expectations rather than noisy observations as we discuss later.
The purpose of this chapter is to provide intuition about both the behavior of the algorithm and the structure of the regret analysis.
The formal proof is presented in Appendix~\ref{sec:stability}, where we rigorously handle the noisy observations and establish the regret bounds.

First, we argue that we should think of $e$ as the primary loss function and control the regret with respect to $e$.
The key observation is that, for every comparator $x\in\K$,
\begin{align}
    f(A_t)-f(x)
    \leq e(X_t)-e((1-\eps)x)+\eps.
\end{align}
Thus controlling regret for $e$ against the contracted comparator controls
the original regret, with an additional cost of at most $n\eps \leq 100 d^2 M \ell^8 \sqrt{n}$ when $\eps=\eps_0$.\footnote{
    For $\eps_0>1/2$ the regret bound is immediate as explained in the main proof.
}
We therefore work directly with $e$ as the loss and $Y$ as its noisy
observation, satisfying $\E[Y\mid X]=e(X)$.

\subsection{Single Point Second Order Approximation}
\label{sec:quadratic-approximation}
Learning $e$ from noisy observations $(X_t,Y_t)$ is challenging, especially
because the space of convex functions is infinite dimensional. This is
quite different from linear bandits, where we only need to estimate a
finite-dimensional loss vector, which can be easily estimated, for example through least squares.
What is possible to estimate is a local quadratic approximation to a smoothed version of $e$ around the mean $\mu_t$ of the Gaussian distribution as we explain next.

For a covariance matrix $\Sigma_t$, define the \textit{Gaussian-smoothing} of $e$ around $\mu_t$ as
\begin{align}
    F_{\Sigma_t}(y)
    =\E[e(y+Z)],\qquad Z\sim\N(0,\Sigma_t).
\end{align}
Note that $F_{\Sigma_t}(\mu_t)$ is the expected loss the algorithm actually incurs when it samples $X_t \sim \N(\mu_t, \Sigma_t)$.
Therefore, the regret in round $t$ can be expressed as\footnote{Recall that $x_\star = \argmin_{x \in \mathcal{K}_\eps}e(x)$.}
\begin{align}
    \E_t[e(X_t)]-e(x_\star)
    = F_{\Sigma_t}(\mu_t) - e(x_\star).
\end{align}
It's also worth mentioning that $F_{\Sigma_t}$ is infinitely differentiable and by Jensen's inequality upper bounds the original function $e$ uniformly, i.e., $F_{\Sigma_t}(y) \ge e(y)$ for all $y \in \R^d$.

In round $t$, the algorithm samples $X_t \sim \N(\mu_t, \Sigma_t)$, and it is easy to see that the loss $Y_t$ it observes
is an unbiased estimate of the Gaussian-smoothed loss $F_{\Sigma_t}(\mu_t)$, i.e.,
\begin{align}
    \E_t[Y_t] = \E_t[\E_t[Y_t\mid X_t]] = F_{\Sigma_t}(\mu_t)\,.
\end{align}
However, the algorithm can do more than just estimate the Gaussian-smoothed loss at the mean; it can also estimate
the first and the second derivatives of $F_{\Sigma_t}$ at $\mu_t$.
To see this, let's start with the gradient first as
\begin{align}
    \nabla F_{\Sigma_t}(\mu_t) & = \E_t[\nabla e(X_t)] = \int_{\R^d} \nabla e(x) \gamma_{\mu_t,\Sigma_t}(x) \, dx
    = -\int_{\R^d} e(x) \nabla \gamma_{\mu_t, \Sigma_t}(x) \, dx                                                  \\
                               & = \E_t\!\left[\Sigma_t^{-1}(X_t-\mu_t)e(X_t)\right]
    = \E_t\!\left[\Sigma_t^{-1}(X_t-\mu_t)Y_t\right]
    \,,
\end{align}
where the third inequality follows from integration by parts and the fact that the Gaussian density vanishes at infinity, so we can drop
the boundary terms, and the final inequality follows from $E_t[Y_t\mid X_t] = e(X_t)$.
Note that this provides $\Sigma_t^{-1}(X_t-\mu_t)e(X_t)$ as an unbiased estimator of the gradient $\nabla F_{\Sigma_t}(\mu_t)$.
Similarly, one can obtain an unbiased estimator of the Hessian $\nabla^2 F_{\Sigma_t}(\mu_t)$ as
\begin{align}
    \nabla^2 F_{\Sigma_t}(\mu_t) = \E_t\!\left[\left(\Sigma_t^{-1}(X_t-\mu_t)(X_t-\mu_t)^\top \Sigma_t^{-1} - \Sigma_t^{-1}\right) Y_t\right]\,.
\end{align}
These are precisely the estimators $g_t$ and $H_t$ in \cref{eq:sampled-hermite-coefficients}.
Let $\lambda>0$ be a scaling parameter for the quadratic term, whose purpose will be clarified later, and define
\begin{align}
    q_t^\lambda(x) = Y_t + \ip{g_t}{x-\mu_t} + \frac\lambda2 (x-\mu_t)^\top H_t (x-\mu_t)\,,
\end{align}
which is an unbiased estimator of the second-order Taylor expansion of $F_{\Sigma_t}$ around $\mu_t$ as
\begin{align}
    \bar{q}_t^\lambda(x) : = F_{\Sigma_t}(\mu_t) + \ip{\nabla F_{\Sigma_t}(\mu_t)}{x-\mu_t} + \frac\lambda2 \|(x-\mu_t)\|_{\nabla^2 F_{\Sigma_t}(\mu_t)}^2 \, = \E_t[q_t^\lambda(x)].
\end{align}

\begin{figure}[!t]
    \centering
    \begingroup
    \definecolor{onsloss}{HTML}{46515B}
    \definecolor{onssmooth}{HTML}{147D92}
    \definecolor{onsquad}{HTML}{C46A29}
    % Exact example: e(x)=|x|+0.35x, slopes -0.65 and 1.35, Sigma_t=9/4, mu_t=0.7.
    % The wider Gaussian keeps the smoothing visibly above the loss.
    \def\onsSmoothCoordinates{%
        (-2.200,1.52482059)
        (-2.150,1.49967445)
        (-2.100,1.47500443)
        (-2.050,1.45083354)
        (-2.000,1.42718535)
        (-1.950,1.40408389)
        (-1.900,1.38155369)
        (-1.850,1.35961973)
        (-1.800,1.33830735)
        (-1.750,1.31764229)
        (-1.700,1.29765058)
        (-1.650,1.27835853)
        (-1.600,1.25979267)
        (-1.550,1.24197969)
        (-1.500,1.22494641)
        (-1.450,1.20871970)
        (-1.400,1.19332643)
        (-1.350,1.17879341)
        (-1.300,1.16514733)
        (-1.250,1.15241467)
        (-1.200,1.14062170)
        (-1.150,1.12979434)
        (-1.100,1.11995812)
        (-1.050,1.11113813)
        (-1.000,1.10335894)
        (-0.950,1.09664452)
        (-0.900,1.09101820)
        (-0.850,1.08650255)
        (-0.800,1.08311939)
        (-0.750,1.08088967)
        (-0.700,1.07983342)
        (-0.650,1.07996969)
        (-0.600,1.08131651)
        (-0.550,1.08389080)
        (-0.500,1.08770834)
        (-0.450,1.09278373)
        (-0.400,1.09913030)
        (-0.350,1.10676011)
        (-0.300,1.11568391)
        (-0.250,1.12591106)
        (-0.200,1.13744957)
        (-0.150,1.15030599)
        (-0.100,1.16448547)
        (-0.050,1.17999168)
        (0.000,1.19682684)
        (0.050,1.21499168)
        (0.100,1.23448547)
        (0.150,1.25530599)
        (0.200,1.27744957)
        (0.250,1.30091106)
        (0.300,1.32568391)
        (0.350,1.35176011)
        (0.400,1.37913030)
        (0.450,1.40778373)
        (0.500,1.43770834)
        (0.550,1.46889080)
        (0.600,1.50131651)
        (0.650,1.53496969)
        (0.700,1.56983342)
        (0.750,1.60588967)
        (0.800,1.64311939)
        (0.850,1.68150255)
        (0.900,1.72101820)
        (0.950,1.76164452)
        (1.000,1.80335894)
        (1.050,1.84613813)
        (1.100,1.88995812)
        (1.150,1.93479434)
        (1.200,1.98062170)
        (1.250,2.02741467)
        (1.300,2.07514733)
        (1.350,2.12379341)
        (1.400,2.17332643)
        (1.450,2.22371970)
        (1.500,2.27494641)
        (1.550,2.32697969)
        (1.600,2.37979267)
        (1.650,2.43335853)
        (1.700,2.48765058)
        (1.750,2.54264229)
        (1.800,2.59830735)
        (1.850,2.65461973)
        (1.900,2.71155369)
        (1.950,2.76908389)
        (2.000,2.82718535)
        (2.050,2.88583354)
        (2.100,2.94500443)
        (2.150,3.00467445)
        (2.200,3.06482059)
    }
    % Left panel.
    \begin{minipage}[t]{0.48\linewidth}
        \centering
        \begin{tikzpicture}[
                x=1.34cm,y=1.00cm,
                line cap=round,line join=round,
                every node/.style={font=\footnotesize}]
            % A compact legend: distinct widths and dashes remain legible at contact.
            \draw[onsloss,line width=0.9pt] (-2.12,3.65) -- (-1.72,3.65);
            \node[anchor=west,text=onsloss] at (-1.65,3.65) {$e$};
            \draw[onssmooth,line width=1.5pt] (-0.48,3.65) -- (-0.08,3.65);
            \node[anchor=west,text=onssmooth] at (-0.01,3.65) {$F_{\Sigma_t}$};
            \draw[onsquad,line width=0.9pt,dash pattern=on 3pt off 2.5pt]
            (1.25,3.65) -- (1.65,3.65);
            \node[anchor=west,text=onsquad] at (1.72,3.65) {$\bar q_t^1$};

            \draw[black!35,line width=0.5pt,->] (-2.43,0) -- (2.48,0)
            node[right,text=black!70] {$x$};
            \draw[black!35,line width=0.5pt] (0,0.035) -- (0,-0.035);
            \node[below=5pt,text=black!70] at (0,0) {$0$};
            \draw[black!35,line width=0.5pt] (0.7,0.035) -- (0.7,-0.035);
            \node[below=5pt,text=black!70] at (0.7,0) {$\mu_t$};

            % Jensen's gap, with one copy of the smoothing data shared by fill and line.
            \path[fill=onssmooth!7] plot[smooth] coordinates {\onsSmoothCoordinates}
            -- (2.2,2.97) -- (0,0) -- (-2.2,1.43) -- cycle;
            \draw[onsloss,line width=0.9pt] (-2.2,1.43) -- (0,0) -- (2.2,2.97);
            \draw[onssmooth,line width=1.5pt]
            plot[smooth] coordinates {\onsSmoothCoordinates};
            \draw[onsquad,line width=0.9pt,dash pattern=on 3pt off 2.5pt,
                domain=-2.2:2.2,samples=101]
            plot (\x,{1.5698334202+0.7092616182*(\x-0.7)+0.2385222861*(\x-0.7)^2});

            \draw[black!35,line width=0.5pt,densely dashed]
            (0.7,0.09) -- (0.7,1.5028334202);
            \draw[black!40,line width=0.5pt] (0.65,1.6898334202) -- (0.2,2.31);
            \node[align=center,text=black!75,font=\scriptsize] at (-0.1,2.64)
            {Same value, slope,\\and curvature at $\mu_t$};
            \fill[white] (0.7,1.5698334202) circle[radius=2.7pt];
            \fill[black!70] (0.7,1.5698334202) circle[radius=1.6pt];
        \end{tikzpicture}
    \end{minipage}\hfill
    % Right panel: shrinking and recentering the focus region in two dimensions.
    \begin{minipage}[t]{0.48\linewidth}
        \centering
        \begin{tikzpicture}[
                x=1cm,y=1cm,
                line cap=round,line join=round,
                every node/.style={font=\footnotesize}]
            \draw[onssmooth,line width=1pt,dash pattern=on 4pt off 2.5pt]
            (-2.60,3.65) -- (-2.08,3.65);
            \node[anchor=west,text=onssmooth] at (-1.98,3.65) {Round $t$};
            \draw[onsquad,line width=1.3pt] (0.10,3.65) -- (0.62,3.65);
            \node[anchor=west,text=onsquad] at (0.72,3.65) {Round $t+1$};

            \draw[black!25,line width=0.5pt,->] (-3.38,0) -- (2.48,0)
            node[right,text=black!60] {$x_1$};
            \draw[black!25,line width=0.5pt,->] (-3.38,0) -- (-3.38,3.20)
            node[above,text=black!60] {$x_2$};

            % Both ellipses have principal axes rotated by 20 degrees.
            % Squared semiaxes decrease from (6.5025,1.5625) to (2.1025,0.5184),
            % so their shape matrices satisfy strict Loewner shrinkage.
            \path[fill=onssmooth!6,rotate around={20:(-0.7,1.8)}]
            (-0.7,1.8) ellipse[x radius=2.55cm,y radius=1.25cm];
            \path[fill=onsquad!11,rotate around={20:(0.85,2.2)}]
            (0.85,2.2) ellipse[x radius=1.45cm,y radius=0.72cm];
            % Draw both boundaries after the fills so neither outline is obscured.
            \draw[onssmooth,line width=1pt,dash pattern=on 4pt off 2.5pt,
                rotate around={20:(-0.7,1.8)}]
            (-0.7,1.8) ellipse[x radius=2.55cm,y radius=1.25cm];
            \draw[onsquad,line width=1.3pt,rotate around={20:(0.85,2.2)}]
            (0.85,2.2) ellipse[x radius=1.45cm,y radius=0.72cm];

            \draw[black!65,line width=0.8pt,->,shorten <=4pt,shorten >=5pt]
            (-0.7,1.8) -- (0.85,2.2);
            % node[midway,above=5pt,sloped,font=\scriptsize] {mean update};
            \fill[onssmooth] (-0.7,1.8) circle[radius=2pt];
            \node[below=5pt,text=onssmooth] at (-0.7,1.8) {$\mu_t$};
            \fill[onsquad] (0.85,2.2) circle[radius=2pt];
            \node[below=5pt,text=onsquad] at (0.85,2.2) {$\mu_{t+1}$};

            % x_star lies inside both ellipses, but outside the smaller ellipse
            % if it were kept at the old mean: normalized squared radius 2.208.
            \fill[white] (1.35,2.25) circle[radius=3pt];
            \fill[black!80] (1.35,2.25) circle[radius=1.9pt];
            \node[above=6pt,text=black!80] at (1.35,2.25) {$x_\star$};
            \node[text=black!70] at (-0.35,-0.30) {$\Sigma_{t+1}\preceq\Sigma_t$};
        \end{tikzpicture}
    \end{minipage}
    \endgroup
    \caption{Left: The loss function $e(x)=|x|+0.35x$, its Gaussian smoothing with
        $\Sigma_t=9/4$, and the second-order Taylor polynomial $\bar q$ at $\mu_t=0.7$.
        Right: successive focus regions in two dimensions. The covariance shrinks,
        while the mean moves so that $x_\star$ remains in both regions.
    }
    \label{fig:gaussian-quadratic}
\end{figure}
\FloatBarrier

\subsection{The focus region}
At this point, it is unclear if this quadratic approximation is a good surrogate for regret minimization.
The following theorem provides a connection between the regret of the quadratic surrogate and the regret of the original loss.\alireza{Consider adding the proof of this in the main text.}
\begin{lemma}[Regret switching]
    \label{lem:regret-switch-main}
    Let $\bar{H}_t = \E_t[H_t]$.
    Then for any $x_0 \in \R^d$, we have
    \begin{align}
        F_{\Sigma_t}(\mu_t) - e(x_0)
        =
        \bigl(\bar{q}_t^\lambda(\mu_t) - \bar{q}_t^\lambda(x_0)\bigr)
        +\tr(\bar{H}_t\Sigma_t)
        +\frac\lambda2\normm{x_0-\mu_t}{\bar{H}_t}^2
        -\E_t[D_e(x_0\,\|\,X_t)].
        \eqlabel{eq:q-regret-switch}
    \end{align}
\end{lemma}
Let us emphasize that this lemma provides a direct link between the regret of the quadratic surrogate and the regret of the original loss,
as the term on the LHS $F_{\Sigma_t}(\mu_t) - e(x_0)$ represents the conditional expected regret against the original loss function $e$ and
the term on the RHS $\bigl(\bar{q}_t^\lambda(\mu_t) - \bar{q}_t^\lambda(x_0)\bigr)$ represents the regret of the quadratic surrogate.
However, there are other terms on the RHS.
We deal with the trace term later and for now we focus on the remaining terms.
The next corollary shows as long as the comparator point $x_\star$ lies within a certain ellipsoid centered around the current mean $\mu_t$, which is referred to as the \textit{focus region}, the remaining terms on the RHS of \eqref{eq:q-regret-switch} are well controlled.
\begin{corollary}[Curvature comparison]
    \label{cor:bracket-ell}
    % Let $\ell$, $\lambda$ and $\sigma$ be as in \cref{thm:main}, and let $\eps>0$.
    Suppose that $0\prec\Sigma_t\preceq2\sigma^2I_d$.
    Then for every $x_0\in\R^d$ satisfying
    \begin{align}
        \frac12\normm{\mu_t-x_0}{\Sigma_t^{-1}}^2
        \leq\frac{1}{128^2\lambda^2\ell},
        \ineqlabel{eq:bracket-ell-focus}
    \end{align}
    it holds that
    \begin{align}
        \frac\lambda2\normm{\mu_t-x_0}{\bar{H}_t}^2
        -\E_t[D_e(x_0\,\|\,X_t)]
        \leq\tr(\bar{H}_t\Sigma_t)
        +\frac{n^{-8}}{\eps}.
        \ineqlabel{eq:bracket-ell}
    \end{align}
\end{corollary}
Corollary~\ref{cor:bracket-ell} is obtained from the general comparison in \cref{prop:bracket} in Appendix~\ref{app:bracket} using the parameters in \cref{thm:main}.

The implication here is that as long as the comparator point $x_\star$ lies withing the focus region, we have the following upper bound on the conditional expected regret
\begin{align}
    \sum_{t=1}^n \E_t[e(X_t) - e(x_\star)]
    \leq
    \sum_{t=1}^n (\bar{q}_t^\lambda(\mu_t) - \bar{q}_t^\lambda(x_\star))
    + 2 \sum_{t=1}^n \tr(\bar{H}_t\Sigma_t)
    +\frac{n^{-7}}{\eps}.
\end{align}
It is not trivial to ensure that the comparator point $x_\star$ always lies within the focus region.
More precisely, if we again consider the idealized scenario where $q_t$ perfectly matches $\bar{q}_t$,
and so $H_t = \bar{H}_t$,  the update to the covariance matrix $\Sigma_{t+1} = \Sigma_t + \lambda \eta \bar{H}_t$,
would shrink the size of the focus region in each round as $\bar{H}_t = \nabla^2 F_{\Sigma_t}(\mu_t)$ is a psd matrix by the convexity
of $F_{\Sigma_t}$.
Therefore, the tracking of the comparator point $x_\star$ within the focus region becomes challenging as the focus region shrinks over time.
Note that the choice of $\lambda$ directly controls the shrinkage rate of the focus region.
We will show later that the Online Newton Step update rule can maintain this invariant throughout the rounds, where this would be crucial to choose $\lambda < 1$ appropriately.

\subsection{Online Newton Step}
\label{sec:population-ons}
The mean update in \cref{alg:direct-hermite} minimizes the quadratic approximation together with a penalty for moving away from the current mean.
Whenever $\Sigma_t^{-1}+\lambda\eta H_t\succ0$, completing the square gives
\begin{align}
    \mu_{t+1}
    =\argmin_{\mu\in\R^d}
    \left\{
    \frac12\normm{\mu-\mu_t}{\Sigma_t^{-1}}^2
    +\eta q_t^\lambda(\mu)
    \right\}.
    \eqlabel{eq:projected-ons-update}
\end{align}

If we forget about the estimation error and suppose that $q_t$ perfectly matches $\bar{q}_t$ in an idealized setting,
then the Online Newton Step (ONS) algorithm is corresponding to choosing a distribution $p_{t+1}$ in round $t+1$ which is
the \textit{contineouous exponential weights} applied to the history of quadratic losses $(q_i)_{i\in[t]}$:
\begin{align}
    p_{t+1}(x) \propto \exp\Bigl(-\eta \sum_{i=1}^t q_i(x) - \|x\|^2_{\Sigma_1} \Bigr)\,.
\end{align}
This is a Gaussian distribution since the exponent is a quadratic function of $x$.
Next, we show that the Online Newton Step update rule controls both: the regret of the quadratic surrogate losses and the tracking of the comparator point within the shrinking focus region.
Just from the definition of $\bar{q}_t^\lambda$ and the Online Newton Step update rule, we have
\begin{align}
    \frac12\normm{\mu_{t+1}-x_\star}{\Sigma_{t+1}^{-1}}^2
     & =\frac12\normm{\mu_t-x_\star-\eta\Sigma_{t+1}g_t}
    {\Sigma_{t+1}^{-1}}^2                                                   \\
     & =\frac12\normm{\mu_t-x_\star}{\Sigma_{t+1}^{-1}}^2
    +\eta\ip{g_t}{x_\star-\mu_t}
    +\frac{\eta^2}{2}\normm{g_t}{\Sigma_{t+1}}^2                            \\
     & =\frac12\normm{\mu_t-x_\star}{\Sigma_t^{-1}}^2
    +\frac{\lambda\eta}{2}(\mu_t-x_\star)^\top H_t (\mu_t-x_\star)\nonumber \\
     & \quad +\eta\ip{g_t}{x_\star-\mu_t}
    +\frac{\eta^2}{2}\normm{g_t}{\Sigma_{t+1}}^2                            \\
     & =\frac12\normm{\mu_t-x_\star}{\Sigma_t^{-1}}^2
    +\eta\bigl(q_t^\lambda(x_\star)-q_t^\lambda(\mu_t)\bigr)
    +\frac{\eta^2}{2}\normm{g_t}{\Sigma_{t+1}}^2.
\end{align}
By telescoping the above identity over $t$, we obtain
\begin{align}
    \sum_{t=1}^n
    \bigl(\bar{q}_t^\lambda(\mu_t) - \bar{q}_t^\lambda(x_\star)\bigr)
     & =
    \frac{\norm{x_\star - \mu_1}_{\Sigma_1^{-1}}^2 - \norm{x_\star - \mu_{n+1}}_{\Sigma_{n+1}^{-1}}^2}{2\eta}
    +\frac\eta2\sum_{t=1}^n\normm{g_t}{\Sigma_{t+1}}^2.
    \eqlabel{eq:population-ons-regret}
\end{align}
Next, we show why the sum including the gradient norms $\sum_{t=1}^n\normm{g_t}{\Sigma_{t+1}}^2$ is controlled, under this idealized setting.
\subsection{The Decrease in the expected loss}
\begin{wrapfigure}{r}{0.40\textwidth}
    \vspace{-6pt}
    \centering
    \definecolor{onspotentialold}{HTML}{147D92}
    \definecolor{onspotentialnew}{HTML}{C46A29}
    \definecolor{onspotentialloss}{HTML}{46515B}
    % Exact Gaussian smoothings of e(y)=|y|+0.35y:
    % Sigma_t=9/4 and Sigma_{t+1}=9/25. Reverse the second path for the fill.
    \def\onsWideCoordinates{%
        (-1.8,1.3383074)
        (-1.7,1.2976506)
        (-1.6,1.2597927)
        (-1.5,1.2249464)
        (-1.4,1.1933264)
        (-1.3,1.1651473)
        (-1.2,1.1406217)
        (-1.1,1.1199581)
        (-1.0,1.1033589)
        (-0.9,1.0910182)
        (-0.8,1.0831194)
        (-0.7,1.0798334)
        (-0.6,1.0813165)
        (-0.5,1.0877083)
        (-0.4,1.0991303)
        (-0.3,1.1156839)
        (-0.2,1.1374496)
        (-0.1,1.1644855)
        (0.0,1.1968268)
        (0.1,1.2344855)
        (0.2,1.2774496)
        (0.3,1.3256839)
        (0.4,1.3791303)
        (0.5,1.4377083)
        (0.6,1.5013165)
        (0.7,1.5698334)
        (0.8,1.6431194)
        (0.9,1.7210182)
        (1.0,1.8033589)
        (1.1,1.8899581)
        (1.2,1.9806217)
        (1.3,2.0751473)
        (1.4,2.1733264)
        (1.5,2.2749464)
        (1.6,2.3797927)
        (1.7,2.4876506)
        (1.8,2.5983074)
    }
    \def\onsNarrowCoordinates{%
        (1.8,2.4304586)
        (1.7,2.2958162)
        (1.6,2.1614180)
        (1.5,2.0274050)
        (1.4,1.8939833)
        (1.3,1.7614448)
        (1.2,1.6301888)
        (1.1,1.5007451)
        (1.0,1.3737919)
        (0.9,1.2501682)
        (0.8,1.1308741)
        (0.7,1.0170569)
        (0.6,0.9099786)
        (0.5,0.8109659)
        (0.4,0.7213436)
        (0.3,0.6423559)
        (0.2,0.5750833)
        (0.1,0.5203644)
        (0.0,0.4787307)
        (-0.1,0.4503644)
        (-0.2,0.4350833)
        (-0.3,0.4323559)
        (-0.4,0.4413436)
        (-0.5,0.4609659)
        (-0.6,0.4899786)
        (-0.7,0.5270569)
        (-0.8,0.5708741)
        (-0.9,0.6201682)
        (-1.0,0.6737919)
        (-1.1,0.7307451)
        (-1.2,0.7901888)
        (-1.3,0.8514448)
        (-1.4,0.9139833)
        (-1.5,0.9774050)
        (-1.6,1.0414180)
        (-1.7,1.1058162)
        (-1.8,1.1704586)
    }
    \begin{tikzpicture}[
            x=1.30cm,y=0.88cm,
            line cap=round,line join=round,
            every node/.style={font=\footnotesize}]
        \draw[onspotentialold,line width=1.2pt] (-1.65,3.10) -- (-1.30,3.10);
        \node[anchor=west,text=onspotentialold] at (-1.23,3.10) {$F_{\Sigma_t}$};
        \draw[onspotentialnew,line width=1.2pt] (0.14,3.10) -- (0.49,3.10);
        \node[anchor=west,text=onspotentialnew] at (0.56,3.10) {$F_{\Sigma_{t+1}}$};

        \draw[black!30,line width=0.5pt,->] (-1.98,0) -- (1.97,0)
        node[right,text=black!65] {$y$};
        \draw[black!30,line width=0.5pt] (0,0.035) -- (0,-0.035);
        \node[anchor=north west,xshift=3pt,yshift=-3pt,text=onspotentialloss]
        at (0,0) {$x_\star$};

        \path[fill=onspotentialold!9]
        plot[smooth] coordinates {\onsWideCoordinates}
        -- plot[smooth] coordinates {\onsNarrowCoordinates} -- cycle;
        \draw[onspotentialloss,line width=0.8pt,dash pattern=on 2.5pt off 2pt]
        (-1.8,1.17) -- (0,0) -- (1.8,2.43);
        \fill[onspotentialloss] (0,0) circle[radius=1.6pt];
        \node[text=onspotentialloss] at (-1.55,0.59) {$e$};
        \draw[onspotentialold,line width=1.2pt]
        plot[smooth] coordinates {\onsWideCoordinates};
        \draw[onspotentialnew,line width=1.2pt]
        plot[smooth] coordinates {\onsNarrowCoordinates};
        % Compare the losses at the same points, independently of the mean update.
        \draw[black!45,line width=0.65pt,->,shorten <=2pt,shorten >=3pt]
        (0,1.1968268) -- (0,0.4787307);
        \draw[black!45,line width=0.65pt,->,shorten <=2pt,shorten >=3pt]
        (1,1.8033589) -- (1,1.3737919);
        % Exact illustrative ONS step: mu_t=-0.35 -> mu_{t+1}=-1,
        % lambda=0.4131848102, eta=10.9095313965, x_star=0.
        % Distance increases from 0.35 to 1, while expected regret
        % decreases from 1.1067601104 to 0.6737918620 (e(x_star)=0).
        \draw[onspotentialold!45,line width=0.5pt,densely dotted]
        (-0.35,0.04) -- (-0.35,1.1067601);
        \draw[onspotentialnew!55,line width=0.5pt,densely dotted]
        (-1,0.04) -- (-1,0.6737919);
        % The arrow connects the two rounds' expected losses.
        \draw[black!60,line width=0.7pt,->,shorten <=4pt,shorten >=4pt]
        (-0.35,1.1067601) -- (-1,0.6737919);
        \fill[white] (-0.35,1.1067601) circle[radius=2.8pt];
        \fill[onspotentialold] (-0.35,1.1067601) circle[radius=1.8pt];
        \fill[white] (-1,0.6737919) circle[radius=2.8pt];
        \fill[onspotentialnew] (-1,0.6737919) circle[radius=1.8pt];
        \draw[onspotentialold,line width=0.6pt] (-0.35,0.035) -- (-0.35,-0.035);
        \draw[onspotentialnew,line width=0.6pt] (-1,0.035) -- (-1,-0.035);
        \node[below=5pt,text=onspotentialold] at (-0.35,0) {$\mu_t$};
        \node[below=5pt,text=onspotentialnew] at (-1,0) {$\mu_{t+1}$};
    \end{tikzpicture}
    \caption{The expected regret decreases even if $\mu_t$ gets further from $x_\star$.}
    \label{fig:covariance-smoothing-decrease}
    \vspace{-4pt}
\end{wrapfigure}
To control the sum of gradient norms, we use the expected loss $F_{\Sigma_t}(\mu_t)$ as a second potential.
We continue in the idealized setting, where $g_t=\bar g_t=\nabla F_{\Sigma_t}(\mu_t)$ and $H_t=\bar H_t\succeq0$.
% We will show that each update decreases this potential by at least $\eta\normm{g_t}{\Sigma_{t+1}}^2/2$.
First, shrinking the covariance can only decrease the expected loss at a fixed mean.
Indeed, the update $\Sigma_{t+1}^{-1}=\Sigma_t^{-1}+\lambda\eta\bar H_t$ gives
$\Sigma_{t+1}\preceq\Sigma_t\preceq\sigma^2I_d$.
% A Gaussian with covariance $\Sigma_t$ can be obtained by adding independent, centered Gaussian noise to one with covariance $\Sigma_{t+1}$.
Convexity of $e$ and Jensen's inequality therefore give
$F_{\Sigma_{t+1}}(y)\leq F_{\Sigma_t}(y)$ for every $y\in\R^d$.

It remains to check that moving the mean also decreases the potential.
The mean moves in a descent direction, but the step must be small enough to avoid overshooting.
Gaussian integration by parts and the Lipschitz property of $e$ give the uniform curvature bound
\begin{align}
    \nabla^2F_{\Sigma_t}(y)
    \preceq\sigma\Lip(e)\,\Sigma_t^{-1}
    \qquad\text{for every }y\in\R^d,
    \label{eq:population-loss-curvature}
\end{align}
where we used $\Sigma_t\preceq\sigma^2I_d$.
For the parameter choices in \cref{thm:main}, the step size satisfies
\begin{align}
    \eta\sigma\Lip(e)
    \leq\frac{8\eta\sigma}{\eps}
    =\frac{16\sigma}{25d\ell^8}
    \leq1.
\end{align}
Applying the curvature bound to the mean update $\mu_{t+1}=\mu_t-\eta\Sigma_{t+1}g_t$, we obtain
\begin{align}
    F_{\Sigma_{t+1}}(\mu_{t+1})
     & \leq F_{\Sigma_t}(\mu_{t+1})
    \leq F_{\Sigma_t}(\mu_t)
    -\eta\normm{g_t}{\Sigma_{t+1}}^2
    +\frac{\eta^2\sigma\Lip(e)}2
    g_t^\top\Sigma_{t+1}\Sigma_t^{-1}\Sigma_{t+1}g_t\nonumber \\
     & \leq F_{\Sigma_t}(\mu_t)
    -\frac\eta2\normm{g_t}{\Sigma_{t+1}}^2,
    \ineqlabel{eq:population-loss-descent}
\end{align}
where the last inequality uses
$\Sigma_{t+1}\Sigma_t^{-1}\Sigma_{t+1}\preceq\Sigma_{t+1}$ and the step-size bound above.
Telescoping now bounds the gradient term in the ONS regret inequality by the total decrease in expected loss:
\begin{align}
    \frac\eta2\sum_{t=1}^n\normm{g_t}{\Sigma_{t+1}}^2
     & \leq F_{\Sigma_1}(\mu_1)
    -F_{\Sigma_{n+1}}(\mu_{n+1}).
    \ineqlabel{eq:population-gradient-budget}
\end{align}
It's also worth mentioning that \cref{eq:population-loss-descent} also shows that the expected regret $\E_{t+1}[e(X_t)]=F_{\Sigma_{t+1}}(\mu_{t+1})$ in round $t+1$ is never greater than that in round $t$, i.e., $\E_{t+1}[e(X_{t+1})] \leq \E_{t}[e(X_{t})]$, which implies that the expected loss is non-increasing over time.

\subsection{The Trace Term}
It remains to control $\sum_{t=1}^n\tr(\bar H_t\Sigma_t)$, which appears in the regret-switching bound.
For the quadratic surrogate, this trace measures the extra expected loss caused by sampling around the mean:
\begin{align}
    \E_t[\bar q_t^\lambda(X_t)]-\bar q_t^\lambda(\mu_t)
    =\frac\lambda2\tr(\bar H_t\Sigma_t).
\end{align}
Large curvature in a direction with large sampling variance makes this term large.
The covariance update responds by reducing the variance in that direction.
We use this reduction to control the accumulated trace cost through the potential $\log\det(\Sigma_t^{-1})$.
In the idealized setting, the covariance update gives
\begin{align}
    \log\frac{\det(\Sigma_{t+1}^{-1})}{\det(\Sigma_t^{-1})}
    =\log\det\!\left(I_d+\lambda\eta\Sigma_t^{1/2}\bar H_t\Sigma_t^{1/2}\right).
\end{align}
The curvature bound from \cref{eq:population-loss-curvature} in the preceding subsection ensures that
\begin{align}
    0\preceq\lambda\eta\Sigma_t^{1/2}\bar H_t\Sigma_t^{1/2}
    \preceq\lambda\eta\sigma\Lip(e)I_d
    \preceq I_d,
\end{align}
and thus the relative change is small in every direction, i.e.
the eigenvalues of $\lambda\eta\Sigma_t^{1/2}\bar H_t\Sigma_t^{1/2}$ are all in $[0,1]$.
Applying $s\leq2\log(1+s)$ to these eigenvalues, yields
\begin{align}
    \tr(\lambda\eta \bar H_t\Sigma_t)
    \leq2
    \log\det\!\left(I_d+\lambda\eta\Sigma_t^{1/2}\bar H_t\Sigma_t^{1/2}\right)
    =2
    \left(\log\det(\Sigma_{t+1}^{-1})-\log\det(\Sigma_t^{-1})\right).
\end{align}
Telescoping bounds the total trace cost by the logarithm of the total covariance shrinkage:
\begin{align}
    \sum_{t=1}^n\tr(\bar H_t\Sigma_t)
     & \leq\frac{2}{\lambda\eta}
    \log\frac{\det(\Sigma_1)}{\det(\Sigma_{n+1})}
    \leq\frac{4d}{\lambda\eta}\log(1+n).
    \ineqlabel{eq:population-trace-budget}
\end{align}
For the last inequality, \cref{lem:predictable-spectrum} gives
$\|\Sigma_{n+1}^{-1}\|_{\op}^{1/2}\leq\sigma^{-1}+n$, since the predictable covariance agrees with $\Sigma_t$ in the idealized setting.
Together with $\Sigma_1=\sigma^2I_d$ and $\sigma\leq1$, this bounds the determinant ratio by $(1+n)^{2d}$.
% Appendix~\ref{app:trace} establishes the corresponding trace bound in the presence of estimation error.

\subsection{The Regret Bound}
We can now combine the preceding bounds to control regret in the idealized setting.
For now, suppose that $x_\star$ remains in the focus region in every round, so that the regret-switching bound applies throughout.
The ONS identity controls the quadratic regret by the initial distance to $x_\star$ and the sum of gradient norms.
The decrease in expected loss controls the gradient sum, while the log-determinant potential controls the trace sum.
Combining \cref{eq:population-ons-regret,eq:population-gradient-budget,eq:population-trace-budget} and dropping the nonnegative terminal-distance term gives
\begin{align}
    \sum_{t=1}^n\E_t[e(X_t)-e(x_\star)]
     & \leq\frac{\normm{x_\star-\mu_1}{\Sigma_1^{-1}}^2}{2\eta}
    +F_{\Sigma_1}(\mu_1)
    % -F_{\Sigma_{n+1}}(\mu_{n+1})
    +\frac{8d\ell}{\lambda\eta}
    +1.
    \ineqlabel{eq:population-regret-bound}
\end{align}
The initialization $\mu_1=0$, $\Sigma_1=\sigma^2I_d$ and the containment $x_\star\in\K_\eps\subseteq\mathbb B_{2d}$ give
\begin{align}
    \frac{\normm{x_\star-\mu_1}{\Sigma_1^{-1}}^2}{2\eta}
    \leq\frac{2d^2}{\eta\sigma^2}.
\end{align}
For the expected-loss potential, the initial sample $X_1\sim\N(0,\sigma^2I_d)$ is spread around the origin, where $e(0)=f(0)\leq1$.
Lipschitz property of $e$ bounds the extra loss from sampling away from the origin by $(8/\eps)\norm{X_1}$, so
\begin{align}
    F_{\Sigma_1}(0)
     & =\E[e(X_1)]
    \leq e(0)+\frac8\eps\E[\norm{X_1}]
    \leq1+\frac8\eps\sqrt{\E[\norm{X_1}^2]}
    =1+\frac{8\sigma\sqrt d}{\eps},
\end{align}
where we used Cauchy--Schwarz and $\E[\norm{X_1}^2]=\tr(\Sigma_1)=d\sigma^2$.
Returning to the original loss adds at most $n\eps$.
Combining \cref{eq:population-regret-bound} with the two initialization bounds gives
\begin{align}
    \E[\operatorname{Reg}_n]
          & \leq n\eps+\sum_{t=1}^n\E[e(X_t)-e(x_\star)]
    \leq n\eps+\frac{2d^2}{\eta\sigma^2}
    +\frac{8\sigma\sqrt d}{\eps}
    +\frac{8d\ell}{\lambda\eta}+2                        \\
    \quad & =O\!\left(d^2M\ell^8\sqrt n\right)
    =\widetilde O\!\left(d^2M\sqrt n\right)\,,
    \ineqlabel{eq:population-original-regret}
\end{align}
which proves the intended regret bound.
% The only remaining task is to ensure that the focus region is maintained and that the estimation error from noisy observations is properly controlled, as discussed in Appendix~\ref{sec:stability}.
% where the last line substitutes the parameters from \cref{thm:main}, with $\eps=\eps_0\leq1/2$.
% This bound is conditional on maintaining the focus region; Appendix~\ref{sec:stability} handles this tracking requirement together with the estimation error from noisy observations.

\newpage
\begin{thebibliography}{9}
    \raggedright

    \bibitem[Bercu et al.(2015)Bercu, Delyon, and Rio]{bercu2015concentration}
    Bernard Bercu, Bernard Delyon, and Emmanuel Rio.
    \newblock \emph{Concentration Inequalities for Sums and Martingales}.
    \newblock SpringerBriefs in Mathematics. Springer, Cham, 2015.
    \newblock \href{https://doi.org/10.1007/978-3-319-22099-4}
    {doi:10.1007/978-3-319-22099-4}.

    \bibitem[Evans and Gariepy(2015)]{evansgariepy2015measure}
    Lawrence C. Evans and Ronald F. Gariepy.
    \newblock \emph{Measure Theory and Fine Properties of Functions}.
    \newblock Revised edition, CRC Press, 2015.

    \bibitem[Fokkema et al.(2024)Fokkema, van der Hoeven, Lattimore, and Mayo]{fokkema2024online}
    Hidde Fokkema, Dirk van der Hoeven, Tor Lattimore, and Jack J. Mayo.
    \newblock Online Newton method for bandit convex optimisation: Extended abstract.
    \newblock In \emph{Proceedings of the 37th Conference on Learning Theory},
    volume 247 of \emph{Proceedings of Machine Learning Research}, pages
    1713--1714, 2024.
    \newblock Full version: \href{https://arxiv.org/abs/2406.06506}{arXiv:2406.06506}.

    \bibitem[Lattimore(2026)]{lattimore2026bandit}
    Tor Lattimore.
    \newblock \emph{Bandit Convex Optimisation}.
    \newblock Cambridge University Press, 2026.
    \newblock \href{https://doi.org/10.1017/9781009607551}{doi:10.1017/9781009607551}.

    \bibitem[Lattimore and Gy\"orgy(2023)]{lattimore2023second}
    Tor Lattimore and Andr\'as Gy\"orgy.
    \newblock A second-order method for stochastic bandit convex optimisation.
    \newblock In \emph{Proceedings of the 36th Conference on Learning Theory},
    volume 195 of \emph{Proceedings of Machine Learning Research}, pages
    2067--2094, 2023.

    \bibitem[Laurent and Massart(2000)]{laurent2000adaptive}
    B\'eatrice Laurent and Pascal Massart.
    \newblock Adaptive estimation of a quadratic functional by model selection.
    \newblock \emph{The Annals of Statistics}, 28(5):1302--1338, 2000.
    \newblock \href{https://doi.org/10.1214/aos/1015957395}
    {doi:10.1214/aos/1015957395}.

    \bibitem[Vershynin(2018)]{vershynin2018high}
    Roman Vershynin.
    \newblock \emph{High-Dimensional Probability: An Introduction with
        Applications in Data Science}.
    \newblock Cambridge University Press, 2018.


\end{thebibliography}

\newpage
\appendix
\section{Stability and regret}
\label{sec:stability}
\label{sec:stochastic-control}

Let $x_\star \in\argmin_{x\in\K_\eps}f(x)$ be a minimizer of $f$ on the contracted set $\K_\eps$.
First, if $\eps_0>1/2$, then, since $f$ takes values in $[0,1]$,
\begin{align}
    \operatorname{Reg}_n\leq n<2n\eps_0=200d^2M\ell^8\sqrt n,
\end{align}
so \cref{thm:main} holds deterministically.
So in all that follows, we assume that $\eps=\eps_0\leq1/2$ and that the parameters $\sigma,\lambda,\eta$ are as in \cref{thm:main}.
\proofpart{I}{Reduction to stopped regret}
\proofstepinline{1}
We first bound realized regret in terms of extension regret, as in \citet{lattimore2026bandit}[Section~10.6]. Define the extension regret by
\begin{align}
    \operatorname{eReg}_{s:u}(x)
    \triangleq
    \sum_{t=s}^{u}\bigl(\E_t[e(X_t)]-e(x)\bigr).
\end{align}
The following lemma shows that bounding the regret with respect to the extension is sufficient to bound the realized regret.
\begin{lemma}[{Adapted from \citet[Section~10.6, Eq.~(10.14)]{lattimore2026bandit}}]
    \label{lem:actual-to-extension}
    Let $\mathcal E_{\mathrm{reg}}$ be the event that
    \begin{align}
        \operatorname{Reg}_n
        \leq
        \operatorname{eReg}_{1:n}(x_\star)
        +n\eps+\sqrt{2n\ell}.
        \ineqlabel{eq:actual-to-extension}
    \end{align}
    Then $\Pp(\mathcal E_{\mathrm{reg}})\geq1-\delta/4$.
\end{lemma}
\proofstepinline{2}
To apply \cref{cor:bracket-ell}, we need to ensure that $\Sigma_t$ remains positive definite and that $x_\star$ remains in the focus ellipsoid around $\mu_t$.
To this end, recall that $\bar{H}_t = \E_t[H_t] = \nabla^2 F_{\Sigma_t}(\mu_t)\succeq 0$ and define the predictable counterpart of $\Sigma_t^{-1}$ as
\begin{align}
    \bar{\Sigma}_1^{-1}=\sigma^{-2}I_d=\Sigma_1^{-1},
    \qquad
    \bar{\Sigma}_{t+1}^{-1}
    =\bar{\Sigma}_t^{-1}+\lambda\eta \bar{H}_t.
    \eqlabel{eq:predictable-covariance}
\end{align}
Next, define $\tau_{\mathrm{cov}}$ as the first time at which the empirical covariance update $\Sigma_{t+1}$ deviates from its predictable counterpart $\bar{\Sigma}_{t+1}$ by more than a constant factor:
\begin{align}
    \tau_{\mathrm{cov}}
    =
    \min\left\{
    t \in [n]:
    \frac12 \bar{\Sigma}_{t+1}^{-1}
    \preceq
    \Sigma_{t+1}^{-1}
    \preceq
    \frac32 \bar{\Sigma}_{t+1}^{-1}
    \text{ fails}
    \right\}\wedge (n + 1).
    \eqlabel{eq:covariance-stopping-time}
\end{align}
Both $\Sigma_{t+1}^{-1} = \Sigma_t^{-1} + \lambda\eta H_t$ and $\bar{\Sigma}_{t+1}^{-1}$ are $\F_t$-measurable, so $\tau_{\mathrm{cov}}$ is a stopping time with respect to $(\F_t)$.
For $t \leq \tau_{\mathrm{cov}}$, the empirical covariance $\Sigma_t$ is positive definite, so we may define the focus radius around $\mu_t$ by
\begin{align}
    R_t = \left\|\mu_t - x_\star\right\|_{\Sigma_t^{-1}}\,, \quad t \leq \tau_{\mathrm{cov}}.
\end{align}
Given a parameter $0 < F \leq \frac{1}{128^2\lambda^2 \ell}$ to be fixed later, let
\begin{align}
    \tau_{\mathrm{focus}}
    =
    \min\left\{
    t \in [n\wedge \tau_{\mathrm{cov}}]:
    R_t^2 > F
    \right\}
    \wedge (n + 1),
\end{align}
which, together with the covariance bound, allows us to apply \cref{cor:bracket-ell} with $x_0=x_\star$ for every $t<\tau_{\mathrm{cov}}\wedge\tau_{\mathrm{focus}}$.
Also define the two additional stopping times
\begin{align}
    \tau_{\mathrm{shell}}
    =\min\{t\in[n \wedge \tau_{\mathrm{cov}}]:v(\mu_{t})>\ell\}\wedge (n+1)\,,
    \eqlabel{eq:energy-shell-stopping-time}
\end{align}
and
\begin{align}
    \tau_{\mathrm{moment}}
    =\min\left\{
    t \in [n \wedge \tau_{\mathrm{cov}}]: \sum_{i=1}^{t} \E_i[Y_i^2] > S
    \right\}\wedge (n+1)\,,
\end{align}
for some parameter $S > 0$ to be fixed later.
Later in the proof, we will use these two additional stopping times to keep the variability of the feedback $Y_t$ under control, allowing us to apply concentration inequalities.
Finally, let $\tau = \tau_{\mathrm{cov}} \wedge \tau_{\mathrm{focus}} \wedge \tau_{\mathrm{shell}} \wedge \tau_{\mathrm{moment}}$,
which represents the first time any of the individual stopping conditions is triggered.

\proofstepinline{3}
Define
\begin{align}
    q_t(x)
     & \triangleq
    \ip{g_t}{x-\mu_t}
    +\frac{\lambda}{2}(x-\mu_t)^\top H_t(x-\mu_t),
    \\
    \bar q_t(x)
     & \triangleq
    \ip{\nabla F_{\Sigma_t}(\mu_t)}{x-\mu_t}
    +\frac{\lambda}{2}(x-\mu_t)^\top
    \nabla^2 F_{\Sigma_t}(\mu_t)(x-\mu_t).
\end{align}
Compared to $q_t^\lambda$ and $\bar q_t^\lambda$, these definitions
drop the constant terms $Y_t$ and $F_{\Sigma_t}(\mu_t)$, respectively.
This is due to the fact that these constants cancel when we take
the difference between the values at $\mu_t$ and $x_\star$ to compute regret.

For every $t<\tau$, the definitions of $\tau_{\mathrm{focus}}$ and
$\tau_{\mathrm{cov}}$ give
\begin{align}
    \frac12 R_t^2 \leq \frac{1}{128^2\lambda^2\ell},
    \qquad\text{and}\qquad
    0 \prec \Sigma_t \preceq 2\bar{\Sigma}_t \preceq 2\sigma^2 I_d.
\end{align}
Thus both conditions of \cref{cor:bracket-ell} hold with $x_0=x_\star$.
Therefore, by \cref{eq:q-regret-switch,cor:bracket-ell} we obtain
\begin{align}
    \operatorname{eReg}_{1:\tau -1}(x_{\star})
    \leq
    % \overline{\operatorname{qReg}}_{1:\tau_1}(x_{\circ})
    \sum_{t=1}^{\tau -1}
    \left(\bar{q}_t(\mu_t)-\bar{q}_t(x_{\star})\right)
    +
    2\sum_{t=1}^{\tau -1}
    \tr(\bar{H}_t\Sigma_t)
    +\frac{(\tau -1) n^{-8}}{\eps}\,.
    \ineqlabel{eq:regret-switch-focus}
\end{align}
The term $(\tau-1)n^{-8}/\eps$ is bounded by $1$, since
$\tau-1\leq n$ and $\eps\geq n^{-1/2}$, and we control the trace sum
using a standard log-determinant argument in Step~7.

Bounding the quadratic regret $
    \overline{\operatorname{qReg}}_{1:\tau -1}(x_{\star})
    \triangleq
    \sum_{t=1}^{\tau -1}
    \left(\bar{q}_t(\mu_t)-\bar{q}_t(x_{\star})\right)
$ in \cref{eq:regret-switch-focus} is quite delicate.
The main difficulty is that the quadratic $\bar{q}_t$ is not observed directly, but only through the sampled stochastic quadratic $q_t$.
Writing $D_t=\bar{q}_t(x_{\star})-q_t(x_{\star})$, which satisfies $\E_t[D_t]=0$,
the absolute difference between the observed and expected quadratic regret is
\begin{align}
    \left|
    \overline{\operatorname{qReg}}_{1:\tau-1}(x_{\star})-\operatorname{qReg}_{1:\tau-1}(x_{\star})
    \right|
    =
    \left|
    \sum_{t=1}^{\tau -1}
    \bar{q}_t(\mu_t) - \bar{q}_t(x_{\star})
    -
    q_t(\mu_t) + q_t(x_{\star})
    \right|
    =
    \left|
    \sum_{t=1}^{\tau-1}
    D_t
    \right|\,,
\end{align}
where $\operatorname{qReg}_{1:\tau-1}(x_{\star})$ is the observed quadratic regret,
and the last equality follows from the fact that $q_t(\mu_t)=\bar{q}_t(\mu_t)=0$.
Next, the following lemma provides an identity for the ONS update that will be useful in analyzing the observed quadratic regret.
\begin{lemma}
    \label{lem:ons-identity}
    For $t < \tau$, we have the following identity:
    \begin{align}
        \frac12
        \left(\normm{\mu_{t+1}-x_\star}{\Sigma_{t+1}^{-1}}^2
        -\normm{\mu_t-x_\star}{\Sigma_t^{-1}}^2
        \right)
         & =
        \eta\bigl(q_t(x_\star)-q_t(\mu_t)\bigr)
        +\frac{\eta^2}{2}\normm{g_t}{\Sigma_{t+1}}^2.
    \end{align}
\end{lemma}
\begin{proof}
    Since $t < \tau_{\mathrm{cov}}$ we have $\Sigma_{t+1}, \Sigma_t \succ 0$.
    Therefore, we obtain
    \begin{align}
        \frac12\normm{\mu_{t+1}-x_\star}{\Sigma_{t+1}^{-1}}^2
         & =\frac12\normm{\mu_t-x_\star-\eta\Sigma_{t+1}g_t}
        {\Sigma_{t+1}^{-1}}^2                                                   \\
         & =\frac12\normm{\mu_t-x_\star}{\Sigma_{t+1}^{-1}}^2
        +\eta\ip{g_t}{x_\star-\mu_t}
        +\frac{\eta^2}{2}\normm{g_t}{\Sigma_{t+1}}^2                            \\
         & =\frac12\normm{\mu_t-x_\star}{\Sigma_t^{-1}}^2
        +\frac{\lambda\eta}{2}(\mu_t-x_\star)^\top H_t (\mu_t-x_\star)\nonumber \\
         & \quad +\eta\ip{g_t}{x_\star-\mu_t}
        +\frac{\eta^2}{2}\normm{g_t}{\Sigma_{t+1}}^2                            \\
         & =\frac12\normm{\mu_t-x_\star}{\Sigma_t^{-1}}^2
        +\eta\bigl(q_t(x_\star)-q_t(\mu_t)\bigr)
        +\frac{\eta^2}{2}\normm{g_t}{\Sigma_{t+1}}^2.
    \end{align}
\end{proof}
Now, the ONS identity from \cref{lem:ons-identity} gives us the following telescoping identity for the observed quadratic regret
\begin{align}
    \operatorname{qReg}_{1:\tau -1}(x_{\star})
    =
    \sum_{t=1}^{\tau-1}
    q_t(\mu_t)-q_t(x_\star)
    =
    \frac{R_1^2 - R_\tau^2}{2\eta }
    +\frac{\eta}{2}\sum_{t=1}^{\tau-1}\normm{g_t}{\Sigma_{t+1}}^2.
    \eqlabel{eq:ons-telescope}
\end{align}

Now, starting from the definition of the extension regret and the stopping time $\tau$, we have
\begin{align}
    \operatorname{eReg}_{1:\tau -1}(x_{\star})
     &
    \explan{\leq}{a}
    \overline{\operatorname{qReg}}_{1:\tau -1}(x_{\star})
    +
    2\sum_{t=1}^{\tau -1}
    \tr(\bar{H}_t\Sigma_t)
    +\frac{(\tau -1) n^{-8}}{\eps}
    \\
     &
    \explan{\leq}{b}
    \operatorname{qReg}_{1:\tau -1}(x_{\star})
    +
    \left|
    \sum_{t=1}^{\tau-1} D_t
    \right|
    +
    2\sum_{t=1}^{\tau -1}
    \tr(\bar{H}_t\Sigma_t)
    +1
    \\
     &
    \explan{\leq}{c}
    \frac{R_1^2 - R_\tau^2}{2\eta}
    +
    \frac{\eta}{2}
    \sum_{t=1}^{\tau-1}\normm{g_t}{\Sigma_{t+1}}^2
    +
    \left|
    \sum_{t=1}^{\tau-1} D_t
    \right|
    +
    2\sum_{t=1}^{\tau -1}
    \tr(\bar{H}_t\Sigma_t)
    +1\,,
    \ineqlabel{eq:stopped-regret-after-trace}
\end{align}
where \reason{a} follows from \cref{eq:regret-switch-focus},
\reason{b} follows from the definition of $D_t$ in Step~3 and
$(\tau-1)n^{-8}/\eps\leq1$, and \reason{c} follows from
\cref{eq:ons-telescope}.

The derivation of \cref{eq:stopped-regret-after-trace} is the basis of the regret analysis.
The rest of the proof consists of two main parts: first controlling the terms on the right-hand side of \cref{eq:stopped-regret-after-trace} on a high-probability event, and second showing that $\tau = n+1$ with high probability
so the bound in \cref{eq:stopped-regret-after-trace} gives a high-probability bound on the extension regret up to time $n$.


\proofpart{II}{Bounding the stopped regret}
\proofstepinline{4}
Several later bounds require the fresh randomness of each round to be of typical size.
We collect these requirements into a single event that covers every round up to $n\wedge\tau$.
The first event controls the norm of the standardized Gaussian samples:
\begin{align}
    \mathcal E_Z
    =\left\{\norm{Z_t}^2\leq4d\ell\,\,
    \text{ for all }t\in [1, n \wedge \tau]\right\}\,.
    % \text{ for all }1\leq t\leq n\right\}.
    \eqlabel{eq:gaussian-norm-event}
\end{align}
Next, the sub-Gaussian assumption in \cref{eq:noise} gives the high-probability event
\begin{align}
    \mathcal E_{\mathrm{noise}}
    =\left\{|\xi_t|\leq\sqrt\ell\,\,
    \text{ for all }t\in [1, n \wedge \tau]\right\}\,.
    % \text{ for all }1\leq t\leq n\right\}.
    \eqlabel{eq:noise-cutoff-event}
\end{align}
The displacement $X_t - \mu_t$ is Gaussian with bounded covariance $\Sigma_t\preceq 2\bar{\Sigma}_t \preceq 2 \bar{\Sigma}_1 =2\sigma^2I_d$, and $\pi$ is 1-Lipschitz, so $\pi(X_t - \mu_t)$ can be controlled with high probability by defining the event
\begin{align}
    \mathcal E_{\mathrm{gauge}}
    =\left\{\pi(X_t - \mu_t)\leq3+2\sqrt\ell\,\,
    \text{ for all }t\in [1, n \wedge \tau]\right\}\,.
    \eqlabel{eq:gauge-cutoff-event}
\end{align}
Let $\mathcal E_{\mathrm{global}} = \mathcal E_Z \cap \mathcal E_{\mathrm{noise}} \cap \mathcal E_{\mathrm{gauge}}$.
\begin{lemma}
    \label{lem:global-event}
    It holds that
    \begin{align}
        \Pp(\mathcal E_{\mathrm{global}}^c)
        \leq4ne^{-\ell}\leq\frac\delta4.
        \ineqlabel{eq:global-event-probability}
    \end{align}
\end{lemma}
The proof is in Appendix~\ref{app:concentration}.
In all the steps in this part of the proof, we will assume that the high probability event $\mathcal E_{\mathrm{global}}$ holds.
\proofstepinline{5}
First, we deal with the sum $\frac{\eta}{2}\sum_{t=1}^{\tau - 1}\normm{g_t}{\Sigma_{t+1}}^2$
in \cref{eq:stopped-regret-after-trace}.
We will work with the slightly time-perturbed quantity $\normm{g_t}{\Sigma_{t}}^2$ in place of $\normm{g_t}{\Sigma_{t+1}}^2$ for now.
By the definition of $g_t$ in \cref{eq:sampled-hermite-coefficients}, we have
\begin{align}
    \normm{g_t}{\Sigma_t}^2
    =Y_t^2\norm{Z_t}^2
    \leq
    4d\ell Y_t^2
    ,
    \ineqlabel{eq:energy-feedback-identity}
\end{align}
where we used the bound $\norm{Z_t}^2\leq 4d\ell$ on the event $\mathcal E_{\mathrm{global}}$.
Next, to control the sum $\sum_{t=1}^{\tau-1}Y_t^2$ we use the following version of Freedman's inequality.
\begin{restatable}{lemma}{freedmannonnegative}
    \label{lem:freedman-nonnegative}
    Let $(U_i)_{i=1}^T$ be an adapted sequence.
    Then, for every $b,S,z>0$,
    \begin{align}
        \Pp\left(
        \sum_{i=1}^T U_i^2
        \geq 2S+4bz,
        \quad
        \sum_{i=1}^T\E_i[U_i^2]\leq S,
        \quad
        \max_{1\leq i\leq T}U_i^2\leq b
        \right)
        \leq2e^{-z}.
        \ineqlabel{eq:freedman-nonnegative}
    \end{align}
\end{restatable}
To apply this lemma, with $U_t=Y_t$, we first bound the absolute size of the feedback $|Y_t|^2$.
since $f(A_t)\in[0,1]$, we have
\begin{align}
    |Y_t|=\left|\pi_\eps(X_t)(f(A_t)+\xi_t)+\frac{2v(X_t)}{\eps}\right|
    \leq
    (v(X_t) + 1)(1+|\xi_t|)+\frac{2v(X_t)}{\eps}.
    \ineqlabel{eq:energy-feedback-decomposition}
\end{align}
For the noise $|\xi_t|$, we have the bound $|\xi_t|\leq\sqrt\ell$ on the event $\mathcal E_{\mathrm{global}}$.
% Thus we need to control both the observation noise $|\xi_t|$ and the capped Minkowski functional $v(X_t)$.
To bound $v(X_t)$ for $t < \tau$, write $W_t=X_t-\mu_t$, then the subadditivity of $\pi(\cdot)$ gives
\begin{align}
    v(X_t)\leq v(\mu_t)+\frac{\pi(W_t)}{1-\eps}
    \leq v(\mu_t)+2\pi(W_t)
    \leq v(\mu_t)+2(3 + 2\sqrt\ell)
    \,,
    \ineqlabel{eq:energy-sample-gauge}
\end{align}
where we used the bound $\pi(W_t)\leq 3 + 2\sqrt\ell$ on the event $\mathcal E_{\mathrm{global}}$.
We also need to bound $v(\mu_t)$ in \cref{eq:energy-sample-gauge}.
Projecting the mean onto $\K_\eps$ would make this term zero; our unprojected mean can leave $\K_\eps$.
Instead, we use the stopping time $\tau_{\mathrm{shell}}$ defined in \cref{eq:energy-shell-stopping-time},
which gives
\begin{align}
    v(X_t)
    \leq v(\mu_t) + 2(3 + 2\sqrt\ell)
    \leq 3\ell\,,
\end{align}
where we used $\ell \geq 640$.
Then, for every $t < \tau$, by \cref{eq:energy-feedback-decomposition,eq:energy-sample-gauge} we have
\begin{align}
    |Y_t|
    \leq (3\ell + 1)(1+\sqrt\ell)+\frac{6\ell}{\eps}
    \leq \frac{4\ell^{3/2}}{\eps}
    \ineqlabel{eq:energy-sample-gauge-bound}
\end{align}
where we again used $\ell\geq640$.
Now, by \cref{eq:energy-feedback-identity,eq:energy-sample-gauge-bound}
for $t < \tau$, we have
\begin{align}
    \normm{g_t}{\Sigma_t}^2
    =Y_t^2\norm{Z_t}^2
    \leq\frac{64d\ell^4}{\eps^2}\,.
    \ineqlabel{eq:energy-sample-gauge-bound-2}
\end{align}
To apply \cref{lem:freedman-nonnegative}, it remains to bound the conditional second moment $\sum_{t=1}^{\tau-1}\E_t[Y_t^2]$,
which is bounded by $S$ since $\tau \leq \tau_{\mathrm{moment}}$.
Now, let
\begin{align}
    \mathcal E_Y
    \triangleq
    \left\{
    \sum_{t=1}^{\tau-1}Y_t^2
    \leq
    2S + \frac{4d\ell^5}{\eps^2}
    \right\}\,.
    \eqlabel{eq:interval-feedback-bound}
\end{align}
\begin{lemma}
    \label{lem:feedback-square-sum}
    It holds that
    \begin{align}
        \Pp\left(\mathcal E_{\mathrm{global}}\cap\mathcal E_Y^c\right)
        \leq2e^{-\ell/64}\leq\frac\delta8.
        \ineqlabel{eq:feedback-square-sum-probability}
    \end{align}
\end{lemma}
The proof is an application of \cref{lem:freedman-nonnegative} and is
in Appendix~\ref{app:concentration}.
Therefore, on the event $\mathcal E_\mathrm{global} \cap \mathcal E_Y$, we get
\begin{align}
    \frac{\eta}{2}\sum_{t=1}^{\tau-1}\normm{g_t}{\Sigma_t}^2
    =
    \frac{\eta}{2}\sum_{t=1}^{\tau-1}Y_t^2\|Z_t\|^2
    \leq
    2\eta d\ell
    \sum_{t=1}^{\tau-1}Y_t^2
    \leq
    4\eta d\ell
    \left(
    S + \frac{2d\ell^5}{\eps^2}
    \right)\,.
    \ineqlabel{eq:interval-energy-bound}
\end{align}
Finally, we compare $\normm{g_t}{\Sigma_t}^2$ and $\normm{g_t}{\Sigma_{t+1}}^2$.
For every $t<\tau$, we have
\begin{align}
    \lambda\eta\left\|
    \Sigma_t^{1/2}H_t\Sigma_t^{1/2}
    \right\|_{\op}
     & =\lambda\eta|Y_t|\|Z_tZ_t^\top-I_d\|_{\op}
    \leq\frac{16\lambda\eta d\ell^{5/2}}{\eps}
    \leq\frac12,
\end{align}
where the first inequality follows from
$\|Z_tZ_t^\top-I_d\|_{\op} \leq \max\{1,\|Z_t\|^2\} \leq 4d\ell$,
and the last inequality follows from our parameter choices.
The update $\Sigma_{t+1}^{-1}
    =\Sigma_t^{-1}+\lambda\eta H_t$ therefore implies
$\frac23\Sigma_t
    \preceq\Sigma_{t+1}
    \preceq2\Sigma_t$,
and hence
$
    \frac23\normm{g_t}{\Sigma_t}^2
    \leq\normm{g_t}{\Sigma_{t+1}}^2
    \leq2\normm{g_t}{\Sigma_t}^2
$.
Thus, \cref{eq:interval-energy-bound} gives
\begin{align}
    \frac{\eta}{2}\sum_{t=1}^{\tau-1}\normm{g_t}{\Sigma_{t+1}}^2
    \leq
    8\eta d\ell
    \left(
    S + \frac{2d\ell^5}{\eps^2}
    \right)\,.
    \ineqlabel{eq:interval-energy-bound-final}
\end{align}

\proofstepinline{6}
Next, going back to \cref{eq:stopped-regret-after-trace},
we control the sum $\sum_{t=1}^{\tau}D_t$, which is a martingale difference sequence as $\E_t[D_t]=0$.
We use the following self-normalized martingale inequality \citet[Theorem~3.26]{bercu2015concentration} to control this sum.
\begin{restatable}{lemma}{selfnormalisedmartingale}
    \label{lem:self-normalised-martingale}
    Let $(U_i)_{i=1}^T$ be a martingale-difference sequence. Then, for every $z,v>0$,
    \begin{align}
        \Pp\left(
        \left|\sum_{i=1}^T U_i\right|\geq z\quad\text{and}\quad
        \sum_{i=1}^T \left(U_i^2+\E_i[U_i^2]\right)\leq v
        \right)
        \leq
        2\exp\left(-\frac{z^2}{2v}\right).
        \ineqlabel{eq:self-normalised-martingale}
    \end{align}
\end{restatable}
To apply it, we bound the realized and conditional second moments together.
Since $\bar{q}_t(x_{\star})=\E_t[q_t(x_\star)]$, a straightforward calculation gives
\begin{align}
    \sum_{t=1}^{\tau-1}\left(D_t^2+\E_t[D_t^2]\right)
     & \leq
    2\sum_{t=1}^{\tau-1}
    q_t(x_{\star})^2
    +
    2\sum_{t=1}^{\tau-1}
    \E_t[q_t(x_{\star})^2]\,,
    \ineqlabel{eq:martingale-variance}
\end{align}
where we used $-2\bar{q}_t(x_{\star})q_t(x_{\star}) \leq \bar{q}_t(x_{\star})^2+q_t(x_{\star})^2$ and Jensen's inequality
$\bar{q}_t(x_{\star})^2\leq\E_t[q_t(x_{\star})^2]$.
Now, let $G_t = \langle \Sigma_t^{-1/2}(x_\star-\mu_t)/R_t, Z_t\rangle$
when $R_t>0$, and $G_t = 0$ when $R_t=0$.
% Then we can bound the second moment of $q_t(x_{\circ})$ as
% \begin{align}
%     \E_t\left[q_t(x_{\circ})^2\right]
%      & \leq R_t(x_{\circ})^2\E_t\left[
%         Y_t^2\left(
%         |G_t(x_{\circ})|+\frac12|G_t(x_{\circ})^2-1|
%         \right)^2\right]
%     \\
%      & \quad
%     \leq 4R_t(x_{\circ})^2\E_t\left[
%         Y_t^2\left(G_t(x_{\circ})^2+1\right)^2\right]
%     \leq 8F\E_t\left[
%         Y_t^2\left(G_t(x_{\circ})^2+1\right)^2\right].
% \end{align}
Then just from the definition of $q_t(x_\star)$, we have
\begin{align}
    q_t(x_{\star})^2
     & =
    Y_t^2\left[
        \left\langle \Sigma_t^{-1/2}(x_\star - \mu_t), Z_t \right\rangle
        + \frac{\lambda}{2}\left(
        \left\langle \Sigma_t^{-1/2}(x_\star - \mu_t), Z_t \right\rangle^2
        - \normm{x_\star - \mu_t}{\Sigma_t^{-1}}^2
        \right)
    \right]^2 \\
     & =
    Y_t^2\left[
        R_t G_t
        + \frac{\lambda}{2} R_t^2 \left(G_t^2 - 1\right)
        \right]^2
    \leq
    2R_t^2Y_t^2\left(G_t^2+1\right)^2
    \leq
    4F Y_t^2\left(G_t^2+1\right)^2
    \,,
    \ineqlabel{eq:quadratic-estimate-decomposition}
\end{align}
where we used $(a+b)^2\leq2a^2+2b^2$ and the focus bound $R_t^2\leq 2F \leq \frac{2}{128^2\lambda^2 \ell}$, which also gives $\lambda R_t<1$.
Now, observe that the direction defining $G_t$ is then an
$\F_{t-1}$-measurable unit vector, so as long as $R_t>0$, the conditional distribution of
$G_t$ given $\F_{t-1}$ is $\N(0,1)$.
% Therefore, $|G_t| \leq 2\sqrt{\ell}$ with high probability by standard tail bounds.
This justifies the definition of the event
\begin{align}
    \mathcal E_G
    =\left\{
    |G_t|\leq2\sqrt\ell
    \text{ for all }1\leq t< \tau
    \right\}
    \,,
    \eqlabel{eq:gaussian-score-event}
\end{align}
on which we have $(G_t^2+1)^2\leq25\ell^2$.
Now, on the event $\mathcal E_G$, we can bound the realized part of the RHS of \cref{eq:martingale-variance} as
\begin{align}
    2\sum_{t=1}^{\tau-1}
    q_t(x_\star)^2
    \leq
    2\sum_{t=1}^{\tau-1}
    4F Y_t^2\left(G_t^2+1\right)^2
    \leq
    200 F \ell^2 \sum_{t=1}^{\tau-1}
    Y_t^2
    \,.
    \ineqlabel{eq:quadratic-bound}
\end{align}
For the conditional part of the RHS of \cref{eq:martingale-variance} using the event $\mathcal E_G$ we have
\begin{align}
    2\sum_{t=1}^{\tau-1}
    \E_t[q_t(x_\star)^2]
     & \leq
    200 F \ell^2 \sum_{t=1}^{\tau-1}
    \E_t[Y_t^2]
    +
    8F \sum_{t=1}^{\tau-1}\E_t\left[Y_t^2(G_t^2+1)^2\1_{\{|G_t|>2\sqrt\ell\}}\right]
    \\
     & \leq
    200 F \ell^2 \sum_{t=1}^{\tau-1}
    \E_t[Y_t^2]
    +
    8Fe^{-\ell/16}
    \,,
\end{align}
where the last inequality uses \cref{lem:weighted-gaussian-tail} to bound the summed weighted Gaussian tails.
Thus, starting from \cref{eq:martingale-variance} and
absorbing the tail term into the coefficient gives
\begin{align}
    \sum_{t=1}^{\tau-1}\left(D_t^2+\E_t[D_t^2]\right)
     & \leq
    200 F \ell^2
    \left(
    \sum_{t=1}^{\tau-1}
    \left(Y_t^2
        +
        \E_t[Y_t^2]\right)
    \right)
    + 8Fe^{-\ell/16}
    \\
     &
    \leq300F\ell^2\left(3S+\frac{4d\ell^5}{\eps^2}\right)\,,
\end{align}
where we used the event $\mathcal{E}_Y$ in \cref{eq:interval-feedback-bound} which gives $\sum_{t=1}^{\tau-1}Y_t^2\leq2S+4d\ell^5/\eps^2$.
Now, define the event
\begin{align}
    \mathcal E_{q}
    =\left\{
    \left|\sum_{t=1}^{\tau-1}D_t\right|
    \leq\sqrt{F\ell^3\left(S+\frac{d\ell^5}{\eps^2}\right)}
    \right\},
    \eqlabel{eq:self-normalized-quadratic-event}
\end{align}
The following lemma controls its failure probability.
\begin{lemma}
    \label{lem:quadratic-concentration}
    It holds that
    \begin{align}
        \Pp\left(\mathcal E_{\mathrm{global}}\cap\mathcal E_G\cap\mathcal E_q^c\right)
        \leq4e^{-\ell/64}\leq\frac\delta8.
        \ineqlabel{eq:quadratic-concentration-probability}
    \end{align}
\end{lemma}
The proof is an application of \cref{lem:self-normalised-martingale} and is in Appendix~\ref{app:concentration}.
\proofstepinline{7}
It remains to control the trace sum $\sum_{t=1}^{\tau-1}\tr(\bar H_t\Sigma_t)$ in \cref{eq:stopped-regret-after-trace}.
\begin{lemma}
    \label{lem:stopped-trace}
    It holds that
    \begin{align}
        \eta\sum_{t=1}^{\tau-1}\tr(\bar H_t\Sigma_t)
        \leq
        2\eta\sum_{t=1}^{\tau-1}\tr(\bar H_t\bar\Sigma_t)
        \leq\frac{8d\log(1+n)}{\lambda}.
        \ineqlabel{eq:stopped-trace-budget}
    \end{align}
\end{lemma}
The proof is in \cref{app:trace}.

Fix the two unspecified thresholds as
\begin{align}
    S=16n,
    \qquad
    F=\frac{1}{1000\lambda^2\ell^3}
    =1000d^3M^2\ell^5.
    \eqlabel{eq:global-threshold-choices}
\end{align}
Then, on the event
$\mathcal E_{\mathrm{global}}\cap\mathcal E_Y\cap\mathcal E_G\cap\mathcal E_q$,
from \cref{eq:interval-energy-bound-final} we have
\begin{align}
    \sum_{t=1}^{\tau-1} \normm{g_t}{\Sigma_{t+1}}^2
    \leq
    16d\ell\left(S+\frac{2d\ell^5}{\eps^2}\right)
    \leq288dn\ell\,,
\end{align}
where we used
$d\ell^5/\eps^2=n/(10^4d^3M^2\ell^{11})\leq n$;
from \cref{eq:self-normalized-quadratic-event} we have
\begin{align}
    \left|\sum_{t=1}^{\tau-1}D_t\right|
    \leq\sqrt{F\ell^3\left(S+\frac{d\ell^5}{\eps^2}\right)}
    \leq\sqrt{17Fn\ell^3}
    =\sqrt{17000}\,d^{3/2}M\ell^4\sqrt n\,;
\end{align}
and from the initialization in \cref{alg:direct-hermite},
$\norm{x_\star}\leq2d$, we have
\begin{align}
    \frac{R_1^2}{2\eta}
    =\frac{\norm{x_\star}^2}{2\eta\sigma^2}
    \leq\frac{4d^3M^2}{\eta}
    =\frac12d^2M\sqrt n.
\end{align}

Plugging these bounds into \cref{eq:stopped-regret-after-trace} gives a complete control over the extension regret before the stopping time:
\begin{align}
    \operatorname{eReg}_{1:\tau-1}(x_\star)
     & \leq \frac{1}{2} d^2 M \sqrt n
    + 144\eta d n \ell
    + 131 d^{3/2} M \ell^4 \sqrt n
    + \frac{16d\log(1+n)}{\lambda\eta}
    +1
    \\
     & \leq
    \frac{125}{4}d^2M\ell^5\sqrt n.
    \ineqlabel{eq:global-stopped-regret}
\end{align}
The last inequality uses $\eta=8dM/\sqrt n$,
$\log(1+n)\leq\ell/128$, $\ell\geq640$, and $M\sqrt d\geq1$.

The remaining steps of the proof are devoted to showing that $\tau=n+1$
on a high-probability event contained in
$\mathcal E_{\mathrm{global}}\cap\mathcal E_Y\cap\mathcal E_G\cap\mathcal E_q$.
On that event, \cref{eq:global-stopped-regret} bounds the full-horizon
extension regret.

\proofpart{III}{Ruling out early stopping}
We show that $\tau=n+1$ through a proof by contradiction argument:
\begin{align}
    \text{suppose for contradiction that } \tau \leq n\,.
\end{align}
We will show that under this assumption
\begin{align}
    \tau_{\mathrm{focus}} > \tau,\quad
    \tau_{\mathrm{cov}} > \tau,\quad
    \tau_{\mathrm{shell}} > \tau,\quad
    \tau_{\mathrm{moment}} > \tau,
\end{align}
which contradicts the definition of $\tau$ as the minimum of these stopping times.

\proofstepinline{8}
We first prove that $\tau_{\mathrm{focus}} > \tau$.
On the event
$\mathcal E_{\mathrm{global}}\cap\mathcal E_Y\cap\mathcal E_G\cap\mathcal E_q$,
\cref{eq:stopped-regret-after-trace} and the same estimates used in
\cref{eq:global-stopped-regret} give
\begin{align}
    \operatorname{eReg}_{1:\tau-1}(x_\star)
    +\frac{R_\tau^2}{2\eta}
    \leq\frac{125}{4}d^2M\ell^5\sqrt n.
\end{align}
Since $x_\star$ minimizes $f$ on $\K_\eps$, the extension formula
\cref{eq:extension-lipschitz} gives
$e(x)\geq f(x/\pi_\eps(x))\geq f(x_\star)=e(x_\star)$ for every $x\in\R^d$.
Thus $\operatorname{eReg}_{1:\tau-1}(x_\star)\geq0$, and consequently
\begin{align}
    R_\tau^2
    \leq2\eta\cdot\frac{125}{4}d^2M\ell^5\sqrt n
    =500d^3M^2\ell^5
    =\frac F2<F,
\end{align}
where we used $\eta=8dM/\sqrt n$ and \cref{eq:global-threshold-choices}.
As $\tau\leq n\wedge\tau_{\mathrm{cov}}$ the focus condition is well-defined
and holds at $\tau$, so $\tau_{\mathrm{focus}}>\tau$.

\proofstepinline{9}
Proving that $\tau_{\mathrm{shell}}>\tau$ is more involved.
The proof relies on a local regret analysis over short intervals and turning these local bounds into a pointwise control of the extension value at the current iterate, proving that $v(\mu_\tau)<\ell$ in the end.

Nonnegativity of $f$ and the extension formula \cref{eq:extension-lipschitz} give
\begin{align}
    2v(x)\leq\eps e(x)\qquad\text{for every }x\in\R^d,
\end{align}
so it suffices to upper bound the extension value $e(\mu_\tau)$.
By convexity and Jensen's inequality, we have
\begin{align}
    e(\mu_t) \leq \E_t[e(X_t)]\,,
\end{align}
which is the expected extension loss at time $t$,
which suggests that a control over the expected extension loss $\E_t[e(X_t)]$ will lead to a control over $e(\mu_t)$.

We consider an interval $[s,u]$ of length at most
\begin{align}
    k=\left\lceil\frac{\eps F}{\eta}\right\rceil\,,
\end{align}
so that its length $L=u-s+1$ satisfies $L\leq k$.
A crude bound on the number of such short intervals is $n^2$, which suffices for a union bound over all of them in our subsequent probabilistic analysis.

\proofstepinline{9.1}
In the rest of this step, consider a fixed short local interval
\begin{align}
    [s,u]\quad\text{with }\quad u<\tau \quad\text{and}\quad L=u-s+1\leq k\,.
\end{align}
We will almost repeat the global analysis for the local interval, with the
comparator $\mu_s$ (instead of $x_\star$) which is the starting mean of the local interval and is $\F_{s-1}$-measurable.
Write
\begin{align}
    \tilde{R_t}=\normm{\mu_t-\mu_s}{\Sigma_t^{-1}},
    \qquad
    \tilde{D_t}=\bar q_t(\mu_s)-q_t(\mu_s).
\end{align}
Similar to the global definition, define a local focus stopping time for the interval $[s,u]$ as
\begin{align}
    \tilde{\tau}_{\mathrm{focus}}=\inf\{t \in [s,u]:\tilde{R_t}^2>F_{L}\}
    \wedge (u+1)\,,\qquad
    \text{with }\quad F_L = 16\frac{\eta}{\eps}L\ell.
    \eqlabel{eq:shell-local-parameters}
\end{align}

Let
\begin{align}
    \tilde{\tau} = \tau \wedge \tilde{\tau}_{\mathrm{focus}}\,,
\end{align}
and write, for every $s\leq j<\tilde{\tau}$,
\begin{align}
    \operatorname{eReg}_{s:j}(\mu_s)
     &
    \explan{\leq}{a}
    \overline{\operatorname{qReg}}_{s:j}(\mu_s)
    +
    2\sum_{t=s}^{j}
    \tr(\bar{H}_t\Sigma_t)
    +\frac{(j-s+1)n^{-8}}{\eps}
    \\
     &
    \explan{\leq}{b}
    \operatorname{qReg}_{s:j}(\mu_s)
    +
    \left|
    \sum_{t=s}^{j} \tilde{D}_t
    \right|
    +
    2\sum_{t=s}^{j}
    \tr(\bar{H}_t\Sigma_t)
    +1
    \\
     &
    \explan{\leq}{c}
    \frac{- \tilde{R}_{j+1}^2}{2\eta}
    +
    \frac{\eta}{2}
    \sum_{t=s}^{j}\normm{g_t}{\Sigma_{t+1}}^2
    +
    \left|
    \sum_{t=s}^{j} \tilde{D}_t
    \right|
    +
    2\sum_{t=s}^{j}
    \tr(\bar{H}_t\Sigma_t)
    +1\,,
    \ineqlabel{eq:stopped-regret-after-trace-local}
\end{align}
where \reason{a} follows from \cref{eq:q-regret-switch,cor:bracket-ell}
since $\eta/\eps=2/(25d\ell^8)\leq F$, our choice of $k$ gives
\begin{align}
    F_L\leq32F\ell
    =\frac{32}{1000\lambda^2\ell^2}
    \leq\frac{1}{128^2\lambda^2\ell}\,;
\end{align}
\reason{b} follows from the definition of $\tilde{D}_t$; and
$(j-s+1)n^{-8}/\eps\leq1$; and \reason{c} follows from
\cref{eq:ons-telescope} together with the fact that
\begin{align}
    \tilde{R}_s^2 = \normm{\mu_s-\mu_s}{\Sigma_s^{-1}}^2 = 0.
\end{align}

So far this perfectly matches the counterpart analysis of the global regret in \cref{eq:stopped-regret-after-trace} with the exception of $R_s^2$ being zero for the local comparator.
For the trace term, we use the bound from \cref{lem:uniform-trace} for the local interval, which ensures
\begin{align}
    \sum_{t=s}^{u}\tr(\bar H_t\Sigma_t)\leq\frac{5L}{\eps},
\end{align}
On $\mathcal E_{\mathrm{global}}$,
$\Sigma_{t+1}\preceq2\Sigma_t$ from Step~5 give
\begin{align}
    \frac{\eps\eta}{2}\sum_{t=s}^{u}\normm{g_t}{\Sigma_{t+1}}^2
    \leq
    \eps\eta\sum_{t=s}^{u}\normm{g_t}{\Sigma_{t}}^2
    =
    \eps\eta\sum_{t=s}^{u} Y_t^2 \|Z_t\|^2
    \leq64\frac{\eta}{\eps}Ld\ell^4
    =\frac{128L}{25\ell^4}\leq L\,,
    \ineqlabel{eq:shell-local-energy}
\end{align}
where the second to last inequality uses \cref{eq:energy-sample-gauge-bound-2}
and the last inequality uses $\ell\geq 640$.

Moreover, we need the equivalent of the martingale bound from Step~6 for the local comparator.
The second major difference in the local argument is using the deterministic bound
from \cref{lem:single-round-feedback} which gives
\begin{align}
    \sum_{t=s}^{u}\E_t[Y_t^2]\leq\frac{8L\ell^2}{\eps^2}.
\end{align}
With this bound in hand, we define the counterpart of the event in \cref{eq:self-normalized-quadratic-event} as
\begin{align}
    \mathcal E_{q,\mathrm{loc}} = \left\{ \forall [s,j]\subseteq[s,\tilde{\tau}-1]:
    \left|\sum_{t=s}^{j}\tilde D_t\right|
    \leq\sqrt{F_L\ell^3\left(\frac{8L\ell^2}{\eps^2}
        +\frac{d\ell^5}{\eps^2}\right)}
    \right\}\,.
    \eqlabel{eq:shell-local-martingale-event}
\end{align}
\begin{lemma}[Local quadratic concentration]
    \label{lem:local-quadratic-concentration}
    Fix a deterministic interval $1\leq s\leq u\leq n$ with
    $L=u-s+1\leq k$, and let $\mathcal E_{q,\mathrm{loc}}$ be the
    event in \cref{eq:shell-local-martingale-event}. Then
    \begin{align}
        \Pp\left(\mathcal E_{\mathrm{global}}\cap
        \mathcal E_{q,\mathrm{loc}}^c\right)
        \leq6e^{-\ell/64}\leq\frac{\delta}{8n^2}.
        \ineqlabel{eq:local-quadratic-concentration-probability}
    \end{align}
    Consequently, the probability that $\mathcal E_{\mathrm{global}}$
    holds but one of these local events fails over the at most $n^2$
    short intervals is at most $\delta/8$.
\end{lemma}
The proof is in Appendix~\ref{app:concentration}.
On $\mathcal E_{\mathrm{global}}\cap \mathcal E_{q,\mathrm{loc}}$, we have
\begin{align}
    \eps\left|\sum_{t=s}^{j}\tilde D_t\right|
    \leq\sqrt{\frac{256L^2}{25d\ell^2}+\frac{32L\ell}{25}}
    \leq L+2\sqrt{L\ell}.
    \ineqlabel{eq:shell-local-martingale}
\end{align}
Then, on $\mathcal E_{\mathrm{global}}\cap\mathcal E_{q,\mathrm{loc}}$ and using the above bounds into
\cref{eq:stopped-regret-after-trace-local} we get, for every $s\leq j<\tilde\tau$,
\begin{align}
    \eps\operatorname{eReg}_{s:j}(\mu_s)
    +\frac{\eps}{2\eta}\tilde R_{j+1}^2
     & \leq
    \frac{\eps\eta}{2}\sum_{t=s}^{j}\normm{g_t}{\Sigma_{t+1}}^2
    +\eps\left|\sum_{t=s}^{j}\tilde D_t\right|
    +2\eps\sum_{t=s}^{j}\tr(\bar H_t\Sigma_t)
    +1\nonumber                    \\
     & \quad\leq13L+2\sqrt{L\ell}.
    \ineqlabel{eq:shell-local-potential}
\end{align}

We now show that $\tilde\tau=u+1$, so that the preceding bound applies
to the whole interval $[s,u]$.
Since $s<\tau\leq \tau_{\mathrm{shell}}$, we have $v(\mu_s)\leq\ell$, and the extension formula gives
\begin{align}
    \eps e(\mu_s)
    \leq\eps+(2+\eps)v(\mu_s)\leq3\ell.
\end{align}
Therefore, nonnegativity of $e$ implies that, for every $s\leq j<\tilde\tau$,
\begin{align}
    \eps\operatorname{eReg}_{s:j}(\mu_s)
    \geq-(j-s+1)\eps e(\mu_s)
    \geq-3L\ell.
\end{align}
Substituting this lower bound into \cref{eq:shell-local-potential} gives
\begin{align}
    \tilde R_{j+1}^2
    \leq\frac{2\eta}{\eps}
    \left(3L\ell+13L+2\sqrt{L\ell}\right)
    <16\frac{\eta}{\eps}L\ell=F_L,
\end{align}
where the strict inequality uses $L\geq1$ and $\ell\geq640$.
Since $\tilde R_s=0$, we have $\tilde\tau>s$, so we may take
$j=\tilde\tau-1$ to obtain $\tilde R_{\tilde\tau}^2<F_L$.
If $\tilde\tau\leq u$, then $u<\tau$ implies
$\tilde\tau=\tilde\tau_{\mathrm{focus}}$, whose definition would give
$\tilde R_{\tilde\tau}^2>F_L$, a contradiction.
Hence $\tilde\tau=u+1$, and the same bound for every prefix shows that
\begin{align}
    \max_{s\leq t\leq u+1}\tilde R_t^2\leq F_L.
    \ineqlabel{eq:shell-local-focus}
\end{align}
Taking $j=u$ in \cref{eq:shell-local-potential} and dropping the
nonnegative terminal-distance term gives the desired local regret bound.
Finally, Jensen's inequality gives $e(\mu_t)\leq\E_t[e(X_t)]$, so
\begin{align}
    \eps\sum_{t=s}^{u}\bigl(e(\mu_t)-e(\mu_s)\bigr)
    \leq\eps\operatorname{eReg}_{s:u}(\mu_s)
    \leq13L+2\sqrt{L\ell}.
    \ineqlabel{eq:shell-local-regret}
\end{align}

\proofstepinline{9.2}
We next bound the average extension value over the most recent $k$
rounds, or over all rounds so far if fewer than $k$ rounds have elapsed.
Work on $\mathcal E_{\mathrm{global}}$ and the local events
$\mathcal E_{q,\mathrm{loc}}$ from Step~9.1 for all short interval $[s,u]$ before $\tau$.
Fix $1\leq t<\tau$, and set
\begin{align}
    m=\min\{k,t\},\qquad s=t-m+1.
\end{align}

If $t<k$, then $m=t$ and $s=1$.
Applying \cref{eq:shell-local-regret} on $[1,t]$ with comparator $\mu_1$
gives
\begin{align}
    \frac{\eps}{t}\sum_{j=1}^{t}e(\mu_j)
    \leq\eps e(\mu_1)+13+2\sqrt{\frac{\ell}{t}}
    \leq14+2\sqrt{\frac{\ell}{t}},
\end{align}
where we used $\mu_1=0$ and
$\eps e(\mu_1)=\eps f(0)\leq\eps\leq1/2$.

If $t\geq k$, then $m=k$ and $s=t-k+1$.
Recall that $D_j=\bar q_j(x_\star)-q_j(x_\star)$ is the martingale difference
from Step~3.
We use $x_\star$ as comparator on $[s,t]$.
Let $\mathcal E_{q,\mathrm{local}}$ denote the event
\begin{align}
    \mathcal E_{q,\mathrm{local}}
    =\left\{
    \forall t\in [k, \tau-1]:
    \left|\sum_{j=t-k+1}^{t}D_j\right|
    \leq\sqrt{F\ell^3\left(\frac{8k\ell^2}{\eps^2}
        +\frac{d\ell^5}{\eps^2}\right)}
    \right\}\,,
    \eqlabel{eq:shell-window-martingale-event}
\end{align}
which is the local version of the event $\mathcal E_{q,\mathrm{local}}$
in \cref{eq:self-normalized-quadratic-event} with the local feedback moment budget
$\sum_{j=t-k+1}^{t}\E_j[Y_j^2]\leq \frac{8k\ell^2}{\eps^2}$.
\begin{lemma}
    The event $\mathcal{E}_{q,\mathrm{local}} \cap \mathcal{E}_{\mathrm{global}}$ holds with high probability.
\end{lemma}
Now, on $\mathcal E_{q,\mathrm{local}}$, the choice $k=\lceil\eps F/\eta\rceil$ gives
\begin{align}
    \frac{\eps}{k}\left|\sum_{j=s}^{t}D_j\right|
    \leq\frac{\sqrt{F\ell^3(8k\ell^2+d\ell^5)}}{k}
    \leq\sqrt{\frac{16}{25d\ell^3}+\frac{2}{25k}}
    \leq1,
    \ineqlabel{eq:shell-window-martingale}
\end{align}
where we used $F/k\leq\eta/\eps=2/(25d\ell^8)$.

Repeating the regret-switching and ONS calculation in
\cref{eq:stopped-regret-after-trace-local} on $[s,t]$ with comparator
$x_\star$, and using Jensen's inequality, yields
\begin{align}
    \frac{\eps}{k}\sum_{j=s}^{t}e(\mu_j)
     & \leq\eps e(x_\star)
    +\frac{\eps R_s^2}{2\eta k}
    +\frac{\eps\eta}{2k}\sum_{j=s}^{t}\normm{g_j}{\Sigma_{j+1}}^2
    +\frac{\eps}{k}\left|\sum_{j=s}^{t}D_j\right|
    +\frac{2\eps}{k}\sum_{j=s}^{t}\tr(\bar H_j\Sigma_j)
    +\frac{\eps}{k}\nonumber                                \\
     & \leq\frac12+\frac{\eps F}{2\eta k}+12+\frac{\eps}{k}
    \leq14.
\end{align}
The second inequality uses $R_s^2\leq F$,
$\eps e(x_\star)=\eps f(x_\star)\leq1/2$,
\cref{eq:shell-local-energy} with $L=k$,
\cref{eq:shell-window-martingale}, and \cref{lem:uniform-trace}.
The last inequality uses $\eps F/(\eta k)\leq1$ and $\eps/k\leq1/2$.

Thus, on the intersection of the events above, both cases give
\begin{align}
    \frac{\eps}{m}\sum_{j=t-m+1}^{t}e(\mu_j)
    \leq14+2\sqrt{\frac{\ell}{m}},
    \qquad m=\min\{k,t\},\quad 1\leq t<\tau.
    \ineqlabel{eq:shell-window-average}
\end{align}

\proofstepinline{9.3}
Step~9.2 bounds the average extension value on a recent window.
To control its last mean, we use Step~9.1 to bound how much this average
can increase as we remove the earliest terms of the window.
The following deterministic lemma expresses this implication for an
arbitrary real sequence.
\begin{lemma}[From interval averages to pointwise bounds]
    \label{lem:averages-to-pointwise}
    Let $a_1,\ldots,a_n\in\R$, let $k\geq1$ be an integer, and let $\ell>0$.
    Suppose that, for every $1\leq s\leq u\leq n$ with $u-s+1\leq k$,
    \begin{align}
        \frac{1}{u-s+1}\sum_{j=s}^{u}a_j
        \leq a_s+13+2\sqrt{\frac{\ell}{u-s+1}}.
        \ineqlabel{eq:sequence-local-average}
    \end{align}
    Suppose also that, for every $1\leq u\leq n$,
    \begin{align}
        \frac{1}{\min\{k,u\}}
        \sum_{j=\max\{1,u-k+1\}}^{u}a_j
        \leq14+2\sqrt{\frac{\ell}{\min\{k,u\}}}.
        \ineqlabel{eq:sequence-window-average}
    \end{align}
    Then, for every $1\leq u\leq n$,
    \begin{align}
        a_u\leq27+13\log u+8\sqrt\ell.
        \ineqlabel{eq:sequence-pointwise-bound}
    \end{align}
\end{lemma}
The proof is in \cref{app:averages-to-pointwise}.

We apply the lemma on the same intersection of events as in Step~9.2.
Fix $1\leq t<\tau$ and take $a_j=\eps e(\mu_j)$ for $1\leq j\leq t$.
Dividing \cref{eq:shell-local-regret} by the interval length gives
\cref{eq:sequence-local-average}, and \cref{eq:shell-window-average}
gives \cref{eq:sequence-window-average}.
Thus \cref{lem:averages-to-pointwise} yields
\begin{align}
    \eps e(\mu_t)
    \leq27+13\log t+8\sqrt\ell
    \leq\frac\ell2,
    \qquad 1\leq t<\tau.
    \ineqlabel{eq:shell-pointwise-extension}
\end{align}
For the last inequality, use $\log t\leq\log n\leq\ell/128$,
and $\ell\geq640$.
Finally, the extension bound $2v(\mu_t)\leq\eps e(\mu_t)$ gives
\begin{align}
    v(\mu_t)\leq\frac\ell4,
    \qquad 1\leq t<\tau\,.
    \ineqlabel{eq:shell-pointwise-gauge}
\end{align}

\proofstepinline{9.4}
It remains to control the final update from $\mu_{\tau-1}$ to $\mu_\tau$.
We work on the same intersection of events as in Step~9.3.
If $\tau=1$, then $v(\mu_\tau)=v(0)=0<\ell$.
Otherwise, \cref{eq:shell-pointwise-gauge} gives
$v(\mu_{\tau-1})\leq\ell/4$, so it suffices to show that the final
update cannot use up the remaining distance to the shell threshold.

Since $\tau-1<\tau\leq\tau_{\mathrm{cov}}$, the covariance bounds and
the comparison from Step~5 give
\begin{align}
    \Sigma_{\tau-1}\preceq2\sigma^2I_d\preceq I_d,
    \qquad
    \Sigma_\tau\preceq2\Sigma_{\tau-1},
\end{align}
where $2\sigma^2\leq1$ follows from $M\sqrt d\geq1$.
Consequently,
\begin{align}
    \left\|\Sigma_\tau\Sigma_{\tau-1}^{-1/2}\right\|_{\op}
     & \leq\left\|\Sigma_{\tau-1}^{1/2}\right\|_{\op}
    \left\|\Sigma_{\tau-1}^{-1/2}\Sigma_\tau
    \Sigma_{\tau-1}^{-1/2}\right\|_{\op}
    \leq2.
\end{align}
The mean update in \cref{alg:direct-hermite}, together with
\cref{eq:energy-sample-gauge-bound,eq:gaussian-norm-event}, therefore implies
\begin{align}
    \norm{\mu_\tau-\mu_{\tau-1}}
     & =\eta|Y_{\tau-1}|
    \norm{\Sigma_\tau\Sigma_{\tau-1}^{-1/2}Z_{\tau-1}}
    \leq2\eta|Y_{\tau-1}|\norm{Z_{\tau-1}}
    \leq16\frac{\eta}{\eps}\sqrt d\,\ell^2.
    \ineqlabel{eq:shell-final-update}
\end{align}
Since $v$ is $1/(1-\eps)\leq2$-Lipschitz, this bounds the increase in
the shell value:
\begin{align}
    v(\mu_\tau)
     & \leq v(\mu_{\tau-1})+2\norm{\mu_\tau-\mu_{\tau-1}}
    \leq\frac\ell4+32\frac{\eta}{\eps}\sqrt d\,\ell^2
    =\frac\ell4+\frac{64}{25\sqrt d\,\ell^6}
    <\ell,
\end{align}
where we used $\eta/\eps=2/(25d\ell^8)$ and $\ell\geq640$.
Thus the shell condition holds at $\tau$ in both cases.
It also holds at every earlier round by the definition of $\tau$,
so $\tau_{\mathrm{shell}}>\tau$.

\proofstepinline{10}
We next prove that $\tau_{\mathrm{moment}}>\tau$.
Work on the intersection of the events used for
\cref{eq:global-stopped-regret} and Step~9.
We first use the extension-regret bound to control the conditional
feedback moments through round $\tau-1$, and then add the final round.

For every $t<\tau$, we have $v(\mu_t)\leq\ell$ and
$\Sigma_t\preceq2\sigma^2I_d$.
Applying \cref{lem:feedback} with $u=\tau-1$ therefore gives
\begin{align}
    \sum_{t=1}^{\tau-1}\E_t[Y_t^2]
     & \leq4n+\frac{32\ell}{\eps}
    \operatorname{eReg}_{1:\tau-1}(x_\star)
    +\frac{32ne^{-\ell}}{\eps^2}
    \leq4n+\frac{10n}{\ell^2}
    +\frac{32ne^{-\ell}}{\eps^2}
    \leq6n.
    \ineqlabel{eq:moment-prefix-improvement}
\end{align}
Here \cref{eq:global-stopped-regret} and the value of $\eps$ give
\begin{align}
    \frac{32\ell}{\eps}\operatorname{eReg}_{1:\tau-1}(x_\star)
    \leq\frac{32\ell}{\eps}\cdot\frac{125}{4}d^2M\ell^5\sqrt n
    =\frac{10n}{\ell^2}\leq n.
\end{align}
The remaining tail term is at most $n$, since
$\eps^{-2}=n/(10^4d^4M^2\ell^{16})\leq n$ and
$32ne^{-\ell}\leq1$.

To include round $\tau$, use the shell condition proved in Step~9.4:
$v(\mu_\tau)<\ell$.
Also, $\tau\leq\tau_{\mathrm{cov}}$ ensures
$0\prec\Sigma_\tau\preceq2\sigma^2I_d$.
Thus \cref{lem:single-round-feedback} applies at $\tau$ and yields
\begin{align}
    \E_\tau[Y_\tau^2]
    \leq\frac{8\ell^2}{\eps^2}
    =\frac{n}{1250d^4M^2\ell^{14}}
    \leq n,
\end{align}
where the last inequality uses $M\sqrt d\geq1$ and $\ell\geq640$.
Combining this with \cref{eq:moment-prefix-improvement}, we obtain
\begin{align}
    \sum_{t=1}^{\tau}\E_t[Y_t^2]\leq7n<16n=S.
    \ineqlabel{eq:moment-no-exit}
\end{align}
The moment threshold is therefore not exceeded at $\tau$.
Since $\tau\leq\tau_{\mathrm{moment}}$ by definition,
this proves $\tau_{\mathrm{moment}}>\tau$.

\proofstepinline{11}
We finally prove that $\tau_{\mathrm{cov}}>\tau$.
Intersect the events used in the preceding steps with
$\mathcal E_H$ from \cref{lem:hessian-concentration}.
Step~9 gives $v(\mu_t)\leq\ell$ for every $t\leq\tau$, and
\cref{eq:moment-no-exit} gives
$\sum_{t=1}^{\tau}\E_t[Y_t^2]\leq7n<16n$.
Thus, under our assumption $\tau\leq n$, all the hypotheses of
\cref{lem:hessian-concentration} hold, so
\begin{align}
    -100\ell^{3/2}\sqrt{dn}\,\bar\Sigma_{\tau+1}^{-1}
    \preceq\sum_{t=1}^{\tau}(H_t-\bar H_t)
    \preceq100\ell^{3/2}\sqrt{dn}\,\bar\Sigma_{\tau+1}^{-1}.
\end{align}
The proposed precision update and \cref{eq:predictable-covariance}
start from the same matrix, hence
\begin{align}
    \Sigma_{\tau+1}^{-1}-\bar\Sigma_{\tau+1}^{-1}
    =\lambda\eta\sum_{t=1}^{\tau}(H_t-\bar H_t).
\end{align}
The parameter choices in \cref{thm:main} give
\begin{align}
    100\lambda\eta\ell^{3/2}\sqrt{dn}
    =\frac{4}{5\ell^{5/2}}<\frac14.
\end{align}
Consequently,
\begin{align}
    \frac34\bar\Sigma_{\tau+1}^{-1}
    \preceq\Sigma_{\tau+1}^{-1}
    \preceq\frac54\bar\Sigma_{\tau+1}^{-1}.
    \ineqlabel{eq:covariance-no-exit}
\end{align}
In particular, the proposed matrix is positive definite and satisfies
the covariance condition in \cref{eq:covariance-stopping-time} at
round $\tau$. Since $\tau\leq\tau_{\mathrm{cov}}$, this proves
$\tau_{\mathrm{cov}}>\tau$.

Together with Steps~8--10, all four stopping times are strictly greater
than $\tau$, contradicting its definition as their minimum.
Therefore $\tau=n+1$ on the intersection of the events used above.

\section{Properties of the convex extension}
\label{app:extension}
\begin{lemma}[{\citet[Proposition~3.23]{lattimore2026bandit}}]
    \label{lem:extension}
    The function
    \[
        \bar{e}:
        x
        \longmapsto
        \pi_\eps(x)
        f\!\left(\frac{x}{\pi_\eps(x)}\right)
        +
        \frac{1-\eps}{\eps}v(x)
    \]
    is convex, agrees with $f$ on $\K_\eps$, and
    \begin{align}
        \Lip(\bar e)\leq\frac{2}{\eps(1-\eps)}\,.
    \end{align}
    Moreover,
    $e\!\left(\frac{x}{\pi_\eps(x)}\right)\leq e(x)$
    for all $x\in\R^d$.
\end{lemma}
Now, observe that $e(x)=\bar e(x)+\frac{1+\eps}{\eps}v(x)$, so $e$ is also convex, agrees with $f$ on $\K_\eps$ since $v(x)=0$ for $x\in\K_\eps$.
Moreover, $v$ is $1/(1-\eps)$-Lipschitz, so $e$ is also Lipschitz with
\begin{align}
    \Lip(e)\leq\frac{2}{\eps(1-\eps)}+\frac{1+\eps}{\eps}\cdot\frac{1}{1-\eps}\leq\frac{8}{\eps}\,,
\end{align}
when $\eps\in(0,1/2]$.

\section{Distributional derivatives}
\label{app:distributional-derivatives}
A general convex function $h:\R^d\to\R$ may not be differentiable everywhere and
its Hessian may not exist in the classical sense.
Traditionally, this issue is handled through a mollification process: convolve $h$ with a smooth kernel to obtain a smooth approximation, and then take limits.
However, we will take a different approach that avoids mollification and works directly with the original function $h$.
We employ the theory of distributions, which allows us to define derivatives in a weak sense.
Practically, first we recall the definition of the Schwartz space $\mathcal S(\R^d)$, which consists of smooth functions that decay faster than any polynomial at infinity, along with all their derivatives.
We write $Dh$ and $D^2h$ for the first and second distributional derivatives of $h$, which
are defined as linear functionals on the Schwartz space $\mathcal S(\R^d)$ of smooth rapidly decaying functions.
\begin{align}
    \langle Dh,\varphi\rangle
     & =-\int_{\R^d}h(x)\nabla\varphi(x)\,dx,
     &
    \langle D^2h,\varphi\rangle
     & =\int_{\R^d}h(x)\nabla^2\varphi(x)\,dx,
    \qquad
    \varphi\in\mathcal S(\R^d),
\end{align}
where $\varphi$ is the Gaussian density $\gamma_{\mu,\Sigma}$ for our purposes,
which is a Schwartz function.
The crucial fact here is that if $h$ is convex, then $D^2h$ is a positive-semidefinite matrix-valued \textit{tempered} Radon measure and the Gaussian density $\gamma_{\mu,\Sigma}$ is a Schwartz function, so the pairing $\langle D^2h,\gamma_{\mu,\Sigma}\rangle$ is well-defined and yields a positive-semidefinite matrix.
\alireza{This subsection needs more explanations and citations.}

First we record the standard fact that the distributional Hessian of a convex function is a positive-semidefinite matrix-valued Radon measure.
\begin{lemma}[Hessian measure of a convex function]
    \label{lem:convex-hessian-measure}
    Let $r:\R^d\to\R$ be convex and Lipschitz.
    Then the distributional Hessian $D^2r$ is a positive-semidefinite
    matrix-valued Radon measure.
\end{lemma}
\begin{proof}
    This is standard; see \citet[Section~6.3]{evansgariepy2015measure}.
\end{proof}
% Next lemma gives the Gaussian averaging identity for the Hessian measure, which is the key to the direct Hermite coefficients.
% \begin{lemma}[Gaussian averaging of the Hessian measure]
%     \label{lem:gaussian-hessian-average}
%     Let $r:\R^d\to\R$ be convex and Lipschitz, let
%     $X=\mu+\Sigma^{1/2}Z$, where $Z\sim\N(0,I_d)$ and $\Sigma\succ0$, and
%     let $\gamma_{\mu,\Sigma}$ be the density of $X$.  Then
%     \begin{align}
%         \langle D^2r,\gamma_{\mu,\Sigma}\rangle
%         =
%         \int_{\R^d}
%         \gamma_{\mu,\Sigma}(x)\,D^2r(dx)
%         =
%         \E\left[
%             r(X)\Sigma^{-1/2}(ZZ^\top-I_d)\Sigma^{-1/2}
%             \right].
%         \label{eq:gaussian-hessian-average}
%     \end{align}
% \end{lemma}
% \begin{proof}
%     By \cref{lem:convex-hessian-measure}
%     $D^2r$ is a positive-semidefinite matrix-valued Radon measure.
%     Since $r$ is Lipschitz, $D^2r \in \mathcal S'(\R^d)$ is a tempered distribution,
%     and the Gaussian density $\gamma_{\mu,\Sigma}\in\mathcal S(\R^d)$ is a Schwartz function,
%     which means that $\langle D^2r,\gamma_{\mu,\Sigma}\rangle < \infty$ and
%     \begin{align}
%         \langle D^2r,\gamma_{\mu,\Sigma}\rangle
%         =
%         \langle r,D^2\gamma_{\mu,\Sigma}\rangle
%         =
%         \E\left[
%             r(X)\Sigma^{-1/2}(ZZ^\top-I_d)\Sigma^{-1/2}
%             \right],
%     \end{align}
%     where the first equality is integration by parts applied twice.
% \end{proof}

\section{Regret switching}
\label{app:regret-switch}
\begin{lemma}
    \label{lem:regret-switch}
    Let $X\sim\N(\mu,\Sigma)$, then for every $x\in\R^d$,
    \begin{align}
        \E[e(X)]-e(x)
        =\tr(\bar{H}\Sigma)+\ip{\bar{g}}{\mu-x}-\E[D_e(x\,\|\,X)]\,,
        \eqlabel{eq:regret-switch}
    \end{align}
    where $\bar{g}=\langle De,\gamma_{\mu,\Sigma}\rangle$ and
    $\bar{H}=\langle D^2e,\gamma_{\mu,\Sigma}\rangle$.
\end{lemma}
\begin{proof}
    Splitting $X-x=(X-\mu)+(\mu-x)$ in a rearrangement of \cref{eq:bregman} and taking expectations gives
    \begin{align}
        \E[e(X)]-e(x)
         & =
        \E\left[\ip{\nabla e(X)}{X-\mu}\right]
        +
        \ip{\E[\nabla e(X)]}{\mu-x}
        -
        \E[D_e(x\,\|\,X)]\,,
    \end{align}
    and the proof is completed by showing that
    \[
        \bar{g}=\E[\nabla e(X)]
        \qquad\text{and}\qquad
        \tr(\bar{H}\Sigma)=\E[\ip{\nabla e(X)}{X-\mu}].
    \]

    Since $e$ is Lipschitz, it is differentiable Lebesgue-almost everywhere,
    and its almost-everywhere partial derivatives are bounded by $\Lip(e)$
    and represent its weak derivatives.
    Applying this identity in every coordinate yields $\bar{g}=\E[\nabla e(X)]$,
    which proves the first identity.

    We next identify the trace term. Since
    $\bar{H}=\langle D^2e,\gamma_{\mu,\Sigma}\rangle$,
    distributional integration by parts gives
    \begin{align}
        \tr(\bar{H}\Sigma)
         & =
        \sum_{i,j=1}^d
        \Sigma_{ij}
        \left\langle
        \partial_{ij}e,
        \gamma_{\mu,\Sigma}
        \right\rangle
        =
        -\sum_{i,j=1}^d
        \Sigma_{ij}
        \left\langle
        \partial_i e,
        \partial_j\gamma_{\mu,\Sigma}
        \right\rangle
        \\
         & =
        -\int_{\R^d}
        \left\langle
        \nabla e(x),
        \Sigma\nabla\gamma_{\mu,\Sigma}(x)
        \right\rangle
        \,dx.
    \end{align}
    The Gaussian score identity
    $\Sigma\nabla\gamma_{\mu,\Sigma}(x)
        =
        -(x-\mu)\gamma_{\mu,\Sigma}(x)$
    therefore yields
    \begin{align}
        \tr(\bar{H}\Sigma)
         & =
        \int_{\R^d}
        \left\langle
        \nabla e(x),
        x-\mu
        \right\rangle
        \gamma_{\mu,\Sigma}(x)\,dx
        =
        \E\left[
            \left\langle
            \nabla e(X),
            X-\mu
            \right\rangle
            \right].
    \end{align}
\end{proof}

\section{Curvature and Bregman divergence}
\label{app:bracket}
This appendix gives the general comparison underlying \cref{cor:bracket-ell}.
We state the result with a free tail parameter $\beta$, derive the corollary
for the algorithm parameters, and then prove the general result.

\begin{proposition}[General curvature comparison]
    \label{prop:bracket}
    Let $e:\R^d\to\R$ be convex and $L$-Lipschitz, and let
    $X\sim\N(\mu,\Sigma)$ with $\Sigma\succ0$.
    Write $\bar{H}=\langle D^2e,\gamma_{\mu,\Sigma}\rangle$.
    Then for any $\beta\in(0,1/2]$ and $\lambda\in(0, 1/8]$
    such that $\lambda\left(d+\log\frac2\beta\right) \leq \frac{1}{64}$,
    for all $x\in\R^d$ satisfying
    \begin{align}
        \lambda\normm{\mu-x}{\Sigma^{-1}}
        \sqrt{\log\frac2\beta}                \leq \frac{1}{64},
        \ineqlabel{eq:bracket-focus}
    \end{align}
    it holds that
    \begin{align}
        \frac\lambda2\normm{\mu-x}{\bar{H}}^2
        -\E[D_e(x\,\|\,X)]
        \leq
        \frac1{16}\tr(\bar{H}\Sigma)
        +\frac{\lambda L\beta}{\sqrt{\pi}}
        \norm{\mu-x}\normm{\mu-x}{\Sigma^{-1}}.
        \ineqlabel{eq:bracket}
    \end{align}
\end{proposition}

\subsection{Derivation of the corollary}
The choice $\beta=e^{-\ell}$ makes the tail remainder negligible throughout
the focus region once the covariance is bounded by $2\sigma^2I_d$.

\begin{proof}[Proof of \cref{cor:bracket-ell}]
    Apply \cref{prop:bracket} with $\beta=e^{-\ell}$.
    Since $M\geq d^{-1/2}$, $d\geq1$ and $\ell\geq640$, we have
    $0<\lambda\leq1/8$, $0<e^{-\ell}\leq1/2$, and
    \begin{align}
        \lambda\bigl(d+\log(2e^\ell)\bigr)
        \leq\frac{1+2\ell}{1000\ell^4}
        \leq\frac1{64}.
    \end{align}
    The condition on $x$ implies \cref{eq:bracket-focus} because
    $\log(2e^\ell)\leq2\ell$.
    Moreover, the covariance bound gives
    $\norm{\mu-x}\leq\sqrt2\,\sigma\normm{\mu-x}{\Sigma^{-1}}$.
    Together with \cref{eq:bracket-ell-focus}, this bounds the remainder by
    \begin{align}
        \frac{8\lambda e^{-\ell}}{\eps\sqrt\pi}
        \norm{\mu-x}\normm{\mu-x}{\Sigma^{-1}}
         & \leq\frac{\sigma e^{-\ell}}{512\sqrt{2\pi}\lambda\eps\ell}
        \leq\frac{d\ell^3e^{-\ell}}{\eps}
        \leq\frac{n^{-8}}{\eps}.
    \end{align}
    Here we used the parameter choices, $\ell^3\leq e^{\ell/2}$, and
    $de^{-\ell/2}=d(128nd/\delta)^{-64}\leq n^{-8}$.
\end{proof}

\subsection{Proof of the general comparison}
We first record the measure notation and the curvature representation of
expected Bregman divergence used in the proof.
Let $M$ be a positive-semidefinite matrix-valued Radon measure and let
$\nu=\tr M$ be its trace measure.
Positivity implies that every entry of $M$ is absolutely continuous with
respect to $\nu$.\alireza{Needs proof or citation.}
Hence, by the Radon--Nikodym theorem, there is a
measurable matrix-valued function $Q$ such that\alireza{Needs proof.}
\begin{align}
    M(dz)=Q(z)\nu(dz),
    \qquad
    Q(z)\succeq0,
    \qquad
    \tr Q(z)=1
    \quad\text{for $\nu$-almost every }z.
    \label{eq:matrix-measure-polar}
\end{align}
Therefore, for a nonnegative measurable function $\rho$ and a measurable vector field
$b$, we use the notation
\begin{align}
    \int_B \rho(z)b(z)^\top M(dz)b(z)
    \triangleq
    \int_B
    \rho(z)b(z)^\top Q(z)b(z)\,\nu(dz).
    \eqlabel{eq:matrix-measure-quadratic-form}
\end{align}
In particular, this integral is nonnegative, possibly with value $+\infty$.

\begin{lemma}[Averaged Bregman curvature]
    \label{lem:averaged-bregman-curvature}
    Let $h:\R^d\to\R$ be convex and Lipschitz, let $\bar{x}\in\R^d$, and let
    $Z\sim\N(0,I_d)$.  For $s\in(0,1]$, let $p_s$ denote the density of $Z_s=(1-s)\bar{x}+sZ$.
    Then
    \begin{align}
        \E[D_h(\bar{x}\,\|\,Z)]
        =
        \int_0^1\frac1s
        \int_{\R^d}
        p_s(z)(z-\bar{x})^\top D^2h(dz)(z-\bar{x})\,ds.
        \eqlabel{eq:averaged-bregman-curvature}
    \end{align}
\end{lemma}
\begin{proof}
    Set $\Phi(s)=\E[h(\bar{x}+s(Z-\bar{x}))]$ for $s\geq0$.
    We apply the fundamental theorem of calculus to this averaged
    function. Lipschitz continuity gives $\Phi(s)\to h(\bar{x})$ as
    $s\downarrow0$. For each $s>0$, $h$ is differentiable at $Z_s$
    almost surely, so dominated differentiation gives
    \begin{align}
        \Phi'(s)=\E[\ip{\nabla h(Z_s)}{Z-\bar{x}}],
        \qquad
        |\Phi'(s)|\leq\Lip(h)\E[\norm{Z-\bar{x}}].
    \end{align}

    For $s>0$, we can also differentiate
    $\Phi(s)=\int h(z)p_s(z)\,dz$ twice under the integral:
    $h$ has at most linear growth and the derivatives of $p_s$ have
    Gaussian decay, uniformly on compact intervals of positive $s$.
    Differentiating the Gaussian density gives
    \begin{align}
        \partial_s^2p_s(z)
        =\frac1{s^2}\sum_{i,j=1}^d
        \partial_{ij}\bigl[(z_i-\bar{x}_i)(z_j-\bar{x}_j)p_s(z)\bigr].
    \end{align}
    The functions inside the spatial derivatives are Schwartz
    functions. Distributional integration by parts therefore yields
    \begin{align}
        \Phi''(s)
        =\frac1{s^2}\int_{\R^d}
        p_s(z)(z-\bar{x})^\top D^2h(dz)(z-\bar{x})
        \geq0.
    \end{align}

    For $a\in(0,1)$, integration by parts in $s$ now gives
    \begin{align}
        \int_a^1\frac1s\int_{\R^d}
        p_s(z)(z-\bar{x})^\top D^2h(dz)(z-\bar{x})\,ds
         & =\int_a^1s\Phi''(s)\,ds              \\
         & =\Phi'(1)-a\Phi'(a)-\Phi(1)+\Phi(a).
    \end{align}
    As $a\downarrow0$, the derivative bound implies $a\Phi'(a)\to0$.
    By monotone convergence on the left and continuity of $\Phi$ at
    zero on the right, the full integral equals
    \begin{align}
        \Phi'(1)-\Phi(1)+h(\bar{x})
         & =\E\left[h(\bar{x})-h(Z)
            -\ip{\nabla h(Z)}{\bar{x}-Z}\right]
        =\E[D_h(\bar{x}\,\|\,Z)],
    \end{align}
    as claimed.
\end{proof}


\begin{proof}[Proof of \cref{prop:bracket}]
    Write $\gamma=\gamma_{0,I_d}$ and $\gamma_V=\gamma_{0,V}$ for $V\succ0$.
    First, we switch to the standardized case by defining
    \begin{align}
        \bar{x}=\Sigma^{-\frac12}(x-\mu),
        \qquad
        h(z)=e(\mu+\Sigma^{1/2}z),
        \quad\text{and}
        \quad
        Z=\Sigma^{-\frac12}(X - \mu) \sim \N(0,I_d)\,.
    \end{align}
    Since $\Sigma^{\frac12}\bar{x}=x-\mu$, the chain rule at
    differentiability points of $e$ gives
    $D_h(\bar{x}\,\|\,Z)=D_e(x\,\|\,X)$
    almost surely.

    The plan is to define a set $\mathcal A\subset \R^d$ and control the integrals over $\mathcal A$ and $\mathcal A^c$ separately.
    To this end, write
    \begin{align}
        \begin{aligned}
            \normm{\mu-x}{\bar{H}}^2
             & =
            \bar{x}^\top \Sigma^{\frac12}\bar{H}\Sigma^{\frac12}\bar{x}
            \\
             & =
            \int_{\R^d}\gamma(z)\bar{x}^\top D^2h(dz)\bar{x}
            \\
             & =
            \int_{\mathcal A}\gamma(z)\bar{x}^\top D^2h(dz)\bar{x}
            +
            \int_{\mathcal A^c}\gamma(z)\bar{x}^\top D^2h(dz)\bar{x}
        \end{aligned}
        \eqlabel{eq:standardized-comparator-curvature}
    \end{align}
    In step 1, we upper bound the integral over $\mathcal A$ by the sum of expected Bregman divergence $\E[D_e(x\,\|\,X)]$ and a trace term $\tr(\bar{H}\Sigma)$.
    In step 2, we upper bound the integral over $\mathcal A^c$ by a tail term that is small under the assumptions of the proposition.

    \proofstepinline{1}
    By \cref{lem:averaged-bregman-curvature}, we have
    \begin{align}
        \E[D_e(x\,\|\,X)]
        =
        \E[D_h(\bar{x}\,\|\,Z)]
        =
        \int_0^1\frac1s
        \int_{\R^d}
        p_s(z)(z-\bar{x})^\top D^2h(dz)(z-\bar{x})\,ds\,,
        \eqlabel{eq:bregman-density}
    \end{align}
    where $p_s$ denotes the density of $(1-s)\bar{x}+sZ$.
    Define the set\footnote{
        If $\bar{x}=0$, then the left-hand side of \cref{eq:bracket} is nonpositive, so the inequality holds trivially.
    }
    \begin{align}
        \mathcal A=\left\{z:
        \norm{z}^2\leq2d\log2+4\log\frac2\beta
        \quad
        \text{and}
        \quad
        \left|\ip{\frac{\bar{x}}{\norm{\bar{x}}}}{z}\right|^2
        \leq4\log\frac{2\sqrt2}{\beta}
        \right\}.
        \eqlabel{eq:A}
    \end{align}

    \proofstepinline{1.1}
    Next we show that $\frac12\gamma(z)\leq p_s(z)$ for $s\in[1-2\lambda,1]$ and $z\in\mathcal A$.
    Since $sZ+(1-s)\bar{x}\sim\N((1-s)\bar{x},s^2I_d)$,
    we have
    \begin{align}
        \log\frac{p_s(z)}{\gamma(z)}
         & =-d\log s
        -\frac{1-s^2}{2s^2}\norm{z}^2
        +\frac{1-s}{s^2}\ip{\bar{x}}{z}
        -\frac{(1-s)^2}{2s^2}\norm{\bar{x}}^2\,.
        \eqlabel{eq:density-ratio}
    \end{align}
    For $s\in[1-2\lambda,1]$, the assumption $\lambda\leq1/8$ implies $s\geq\frac34$, and $0\leq1-s\leq2\lambda$.
    Consequently,
    \begin{align}
        \frac{1-s^2}{2s^2}
        \leq
        \frac{1-s}{s^2}
        \leq
        \frac{2\lambda}{(3/4)^2}
        =
        \frac{32}{9}\lambda
        \leq4\lambda,
        \quad\text{and}\quad
        \frac{(1-s)^2}{2s^2}
        \leq
        \frac{(2\lambda)^2}{2(3/4)^2}
        =
        \frac{32}{9}\lambda^2
        \leq4\lambda^2.
    \end{align}
    Moreover, $-d\log s\geq0$. Therefore,
    \cref{eq:density-ratio} yields
    \begin{align}
        \log\frac{p_s(z)}{\gamma(z)}
        \geq
        -4\lambda\norm{z}^2
        -4\lambda\left|\ip{\bar{x}}{z}\right|
        -4\lambda^2\norm{\bar{x}}^2.
        \ineqlabel{eq:density-ratio-lower}
    \end{align}
    Now suppose for $z\in\mathcal A$, then
    together with the assumption $\lambda\left(d+\log\frac2\beta\right)\leq \frac{1}{64}$ we have
    \begin{align}
        \lambda\norm{z}^2
         & \leq
        \lambda\left(
        2d\log2+4\log\left(\frac2\beta\right)
        \right)
        \leq
        4\lambda\left(d+\log\left(\frac2\beta\right)\right)
        \leq
        \frac{1}{16}\,,
        \ineqlabel{ineq:lambda-z-squared-bound}
    \end{align}
    and, again, from $z\in\mathcal A$ we get
    \begin{align}
        \lambda\left|\ip{\bar{x}}{z}\right|
         & \leq
        2\lambda\norm{\bar{x}}
        \sqrt{\log\frac{2\sqrt2}{\beta}}
        \leq
        2\lambda\norm{\bar{x}}
        \sqrt{
            \log\left(
            \frac{2}{\beta}
            \right)
            +\frac12\log2
        }
        \\
         & \leq
        \sqrt5\,
        \lambda\norm{\bar{x}}\sqrt{\log\left(\frac{2}{\beta}\right)}
        \leq
        \frac{\sqrt5}{64}\,,
    \end{align}
    where in the penultimate inequality we used $
        \frac12\log2
        \leq \frac14\log\left(\frac{2}{\beta}\right)$ since $\beta<1/2$,
    and in the last inequality we used \cref{eq:bracket-focus}.
    The last inequality also gives
    $\lambda^2\norm{\bar{x}}^2
        \leq
        \frac{1}{2^{12} \log\left(\frac{2}{\beta}\right)}
        \leq
        \frac{1}{2^{12} \log4}$.
    Substituting these three bounds into
    \cref{eq:density-ratio-lower} gives
    \begin{align}
        \log\frac{p_s(z)}{\gamma(z)}
        \geq
        -4\left(\frac{1}{16}+\frac{\sqrt5}{64}+\frac{1}{2^{12} \log4}\right)
        \geq
        \log\frac12\,.
    \end{align}
    Therefore, we have established
    $p_s(z)
        \geq
        \frac12\gamma(z)$
    for all
    $s\in[1-2\lambda,1]$
    and
    $z\in\mathcal A$.

    \proofstepinline{1.2}
    Since the measure integrand in \cref{eq:bregman-density} is
    nonnegative, we may restrict both integrals to
    $s\in[1-2\lambda,1]$ and $z\in\mathcal A$,
    which together with $-\log(1-2\lambda)\geq2\lambda$ gives
    \begin{align}
        \E[D_h(\bar{x}\,\|\,Z)]
         & \geq
        \int_{1-2\lambda}^1\frac1s
        \int_{\mathcal A}
        p_s(z)(z-\bar{x})^\top D^2h(dz)(z-\bar{x})\,ds
        \\
         & \geq
        \frac12
        \int_{1-2\lambda}^1\frac{ds}{s}
        \int_{\mathcal A}
        \gamma(z)(z-\bar{x})^\top D^2h(dz)(z-\bar{x})
        \\
         & =
        \frac{-\log(1-2\lambda)}{2}
        \int_{\mathcal A}
        \gamma(z)(z-\bar{x})^\top D^2h(dz)(z-\bar{x})
        \\
         & \geq
        \lambda
        \int_{\mathcal A}
        \gamma(z)(z-\bar{x})^\top D^2h(dz)(z-\bar{x}).
        \ineqlabel{eq:bregman-lower}
    \end{align}
    Apply \cref{eq:matrix-measure-polar} to $D^2h$ and write
    $\nu=\tr(D^2h)$, and $D^2h(dz)=Q(z)\nu(dz)$,
    where $Q(z)\succeq0$ and $\tr Q(z)=1$ for $\nu$-almost every $z$.
    For every positive-semidefinite matrix $G$, by the triangle inequality we have
    \begin{align}
        \frac12\bar{x}^\top G\bar{x}
        \leq
        (z-\bar{x})^\top G(z-\bar{x})+z^\top Gz.
    \end{align}
    Applying this inequality to $G=Q(z)$, multiplying by
    $\lambda\gamma(z)$, and integrating over $\mathcal A$ gives
    \begin{align}
        \frac\lambda2
        \int_{\mathcal A}
        \gamma(z)\bar{x}^\top D^2h(dz)\bar{x}
        \leq{} &
        \lambda
        \int_{\mathcal A}
        \gamma(z)(z-\bar{x})^\top D^2h(dz)(z-\bar{x})
        +
        \lambda
        \int_{\mathcal A}
        \gamma(z)z^\top D^2h(dz)z.
    \end{align}
    Moreover, for $\nu$-almost every $z$,
    $z^\top Q(z)z
        \leq
        \norm{z}^2\tr Q(z)
        =
        \norm{z}^2.$
    On $\mathcal A$, \cref{ineq:lambda-z-squared-bound} gives $\lambda\norm{z}^2\leq1/16$.
    Hence,
    \begin{align}
        \lambda
        \int_{\mathcal A}
        \gamma(z)z^\top D^2h(dz)z
        \leq
        \frac1{16}
        \int_{\mathcal A}\gamma(z)\,\nu(dz)
        \leq
        \frac1{16}
        \int_{\R^d}\gamma(z)\tr(D^2h)(dz)
        =
        \frac1{16}\tr(\bar{H}\Sigma),
        \ineqlabel{eq:trace-upper-bound}
    \end{align}
    where the last equality is followed from
    \begin{align}
        \int_{\R^d}\gamma(z)\,\tr(D^2h)(dz)
        %  & =
        % \tr(\langle D^2h, \gamma\rangle)
         & =
        \tr(\langle h, D^2\gamma\rangle)
        =
        \tr\left(\E\left[h(Z)(ZZ^\top-I_d)\right]\right)
        \\
         & =
        \tr(\Sigma^{\frac12}\bar{H}\Sigma^{\frac12})
        =
        \tr(\bar{H}\Sigma).
        \eqlabel{eq:standardized-trace-curvature}
    \end{align}
    Finally, combining \cref{eq:bregman-lower,eq:bregman-lower,eq:trace-upper-bound} we obtain
    \begin{align}
        \frac\lambda2
        \int_{\mathcal A}
        \gamma(z)\bar{x}^\top D^2h(dz)\bar{x}
        \leq
        \E[D_h(\bar{x}\,\|\,Z)]
        +
        \frac1{16}\tr(\bar{H}\Sigma)\,,
        \ineqlabel{eq:bracket-step1}
    \end{align}
    which concludes step 1.

    \proofstepinline{2}
    It remains to bound the curvature measure on $\mathcal A^c$.  Set
    \begin{align}
        u=\frac{\bar{x}}{\norm{\bar{x}}},
        \qquad
        D^2_{\bar{x}}\triangleq\bar{x}^\top D^2h\,\bar{x}.
    \end{align}
    Since $D^2h$ is positive semidefinite, $D^2_{\bar{x}}$ is a nonnegative Radon
    measure.\alireza{we could've just used $D_{\bar{x}}^2$ from the beginning.}
    Direct comparison of Gaussian densities gives
    \begin{align}
        \begin{aligned}
            \gamma(z)\1\left\{\norm z^2>
            2d\log2+4\log\frac2\beta\right\}
             & \leq
            \frac\beta2\gamma_{2I_d}(z)\,\,\,\text{and}\,\,\,
            \\
            \gamma(z)\1\left\{|\ip{u}{z}|^2>
            4\log\frac{2\sqrt2}{\beta}\right\}
             & \leq
            \frac\beta2\gamma_{I_d+uu^\top}(z).
        \end{aligned}
        \ineqlabel{eq:directional-tail}
    \end{align}
    Since $\mathcal A^c$ is contained in the union of the two events above
    and $D^2_{\bar{x}}$ is nonnegative, integrating these inequalities gives
    \begin{align}
        \int_{\mathcal A^c}\gamma(z)\,D^2_{\bar{x}}(dz)
        \leq
        \frac\beta2
        \left(
        \int_{\R^d}\gamma_{2I_d}(z)\,D^2_{\bar{x}}(dz)
        +
        \int_{\R^d}\gamma_{I_d+uu^\top}(z)\,D^2_{\bar{x}}(dz)
        \right).
        \ineqlabel{eq:tail-curvature-domination}
    \end{align}
    We next bound each integral on the right-hand side of
    \cref{eq:tail-curvature-domination}.  Fix $V\in\{2I_d,I_d+uu^\top\}$.
    Since $D^2_{\bar{x}}=\bar{x}^\top D^2h\,\bar{x}$, distributional integration by parts against the Schwartz function
    $\gamma_V$ gives
    \begin{align}
        \int_{\R^d}\gamma_V(z)\,D^2_{\bar{x}}(dz)
         & =
        -\int_{\R^d}
        \ip{\nabla h(z)}{\bar{x}}
        \ip{\nabla\gamma_V(z)}{\bar{x}}\,dz
        \\
         & =
        \int_{\R^d}
        \ip{\nabla h(z)}{\bar{x}}
        \ip{V^{-1}\bar{x}}{z}
        \gamma_V(z)\,dz,
        \eqlabel{eq:wider-gaussian-ibp}
    \end{align}
    where we used
    $\nabla\gamma_V(z)=-V^{-1}z\,\gamma_V(z)$.
    At every differentiability point of $e$, the chain rule gives
    \begin{align}
        \ip{\nabla h(z)}{\bar{x}}
         & =
        \ip{\nabla e(\mu+\Sigma^{1/2}z)}
        {\Sigma^{1/2}\bar{x}}
        =
        \ip{\nabla e(\mu+\Sigma^{1/2}z)}
        {x-\mu}.
    \end{align}
    Since $e$ is $L$-Lipschitz, it follows that
    \begin{align}
        \left|\ip{\nabla h(z)}{\bar{x}}\right|
        \leq
        L\norm{x-\mu}
        \qquad\text{for almost every }z.
    \end{align}
    Therefore,
    \begin{align}
        \int_{\R^d}\gamma_V(z)\,D^2_{\bar{x}}(dz)
         & \leq
        L\norm{x-\mu}
        \int_{\R^d}
        \left|\ip{V^{-1}\bar{x}}{z}\right|
        \gamma_V(z)\,dz
        =
        L\norm{x-\mu}
        \sqrt{\frac{2}{\pi}}
        \sqrt{\bar{x}^\top V^{-1}\bar{x}}.
        \ineqlabel{eq:wider-gaussian-bound}
    \end{align}
    For $V=2I_d$, we have
    $\bar{x}^\top V^{-1}\bar{x}=\frac12\norm{\bar{x}}^2$.
    For $V=I_d+uu^\top$, since $\norm u=1$,
    we have
    $V^{-1}
        =
        I_d-\frac12uu^\top$.
    Moreover, $\bar{x}=\norm{\bar{x}}\,u$, and therefore
    $\bar{x}^\top V^{-1}\bar{x}=
        \bar{x}^\top
        \left(I_d-\frac12uu^\top\right)\bar{x}
        =\frac12\norm{\bar{x}}^2$.
    Thus, for either
    $V\in\{2I_d,I_d+uu^\top\}$,
    we have
    $\sqrt{\bar{x}^\top V^{-1}\bar{x}} = \frac{\norm{\bar{x}}}{\sqrt2}$.
    Substituting this into \cref{eq:wider-gaussian-bound} gives
    \begin{align}
        \int_{\R^d}\gamma_V(z)\,D^2_{\bar{x}}(dz)
        \leq
        \frac{L}{\sqrt\pi}
        \norm{x-\mu}\norm{\bar{x}}.
        \ineqlabel{eq:wider-gaussian}
    \end{align}
    Combining \cref{eq:tail-curvature-domination} and
    \cref{eq:wider-gaussian} now yields
    \begin{align}
        \frac{\lambda}{2}\int_{\mathcal A^c}\gamma(z)\,D^2_{\bar{x}}(dz)
        \leq
        \frac{\beta L \lambda}{2\sqrt\pi}
        \norm{x-\mu}\norm{\bar{x}}.
        \ineqlabel{eq:bracket-step2}
    \end{align}
    Finally, the definition of $\bar{x}$ gives
    $\norm{\bar{x}}^2=\normm{\mu-x}{\Sigma^{-1}}^2$,
    completing the second step.
    Plugging \cref{eq:bracket-step1} and \cref{eq:bracket-step2}
    into \cref{eq:standardized-comparator-curvature}
    completes the proof.
\end{proof}


\section{Trace bounds}
\label{app:trace}

\begin{lemma}[Predictable spectrum]
    \label{lem:predictable-spectrum}
    Suppose that $\eps=\eps_0\leq1/2$. Then
    \begin{align}
        \|\bar{\Sigma}_{t+1}^{-1}\|_{\op}^{\frac12}
        \leq
        \sigma^{-1}+t\,,
        % \lambda\eta\Lip(e)
        \qquad(1\leq t\leq n\wedge\tau_{\mathrm{cov}}),
        \ineqlabel{eq:predictable-spectrum}
    \end{align}
\end{lemma}

\begin{proof}[Proof of \cref{lem:predictable-spectrum}]
    Let $L=\Lip(e)$.
    To prove \cref{eq:predictable-spectrum}, write
    \begin{align}
        \norm{\bar{\Sigma}_{t+1}^{-1}}_{\op}
        =
        \norm{\bar{\Sigma}_{t}^{-1} + \lambda\eta \bar{H}_t}_{\op}
        \leq
        \norm{\bar{\Sigma}_{t}^{-1}}_{\op} + \lambda\eta\norm{\bar{H}_t}_{\op}\,.
    \end{align}
    Next, distributional integration by parts for every $v\in\mathbb{S}_1^d$ gives
    \begin{align}
        \norm{\bar{H}_t}_{\op}
         & =
        \sup_{\|v\|=1} v^\top \bar{H}_t v
        =
        \sup_{\|v\|=1}
        \int_{\R^d}e(x) D^2_v\gamma_{\mu_t,\Sigma_t}(x)\,dx
        \\
         & =
        \sup_{\|v\|=1}
        -\int_{\R^d}\partial_ve(x)\,
        \partial_v\gamma_{\mu_t,\Sigma_t}(x)\,dx
        \\
         & =
        \sup_{\|v\|=1}
        \E_t\left[
            \ip{\nabla e(X_t)}{v}
            \ip{\Sigma_t^{-1}(X_t-\mu_t)}{v}
            \right]
        \leq
        \sup_{\|v\|=1}
        L \|v\|_{\Sigma_t^{-1}}
        =
        L \norm{\Sigma_t^{-1}}_{\op}^{1/2},
        \ineqlabel{eq:H-quadratic-bound}
    \end{align}
    where the inequality is by Cauchy--Schwarz and the Lipschitz property of $e$.
    Plugging this bound into the previous inequality gives
    \begin{align}
        \begin{aligned}
            \norm{\bar{\Sigma}_{t+1}^{-1}}_{\op}
             & \leq
            \norm{\bar{\Sigma}_{t}^{-1}}_{\op} + \lambda\eta L \norm{\Sigma_t^{-1}}_{\op}^{1/2}
            \\
             & \leq
            \norm{\bar{\Sigma}_{t}^{-1}}_{\op} + \sqrt{\frac{3}{2}}\lambda\eta L\norm{\bar{\Sigma}_{t}^{-1}}_{\op}^{1/2}
            \\
             & \leq
            \left(\norm{\bar{\Sigma}_{t}^{-1}}_{\op}^{1/2} + \lambda\eta L\right)^2\,,
        \end{aligned}
        \ineqlabel{eq:predictable-spectrum-bound}
    \end{align}
    where the second inequality uses \cref{eq:covariance-stopping-time}.
    Now, by taking square roots and iterating, we get
    \begin{align}
        \norm{\bar{\Sigma}_{t+1}^{-1}}_{\op}^{1/2}
        \leq
        \norm{\bar{\Sigma}_1^{-1}}_{\op}^{1/2} + t\lambda\eta L
        = \sigma^{-1}+t\lambda\eta L
        \leq
        \sigma^{-1}
        +\frac{64t}{10^5d^{5/2}M\ell^{12}}
        \leq
        \sigma^{-1} + t\,,
    \end{align}
    which proves \cref{eq:predictable-spectrum}.
\end{proof}

\begin{lemma}[Trace sum]
    \label{lem:trace}
    Under the parameter choices of \cref{thm:main}, with
    $\eps=\eps_0\leq1/2$, the following bound holds pathwise for every
    integer $0\leq u\leq n\wedge\tau_{\mathrm{cov}}$:
    \begin{align}
        \eta\sum_{t=1}^u\tr(\bar H_t\Sigma_t)
        \leq2\eta\sum_{t=1}^u\tr(\bar H_t\bar\Sigma_t)
        \leq\frac{8d\log(1+n)}{\lambda}.
        \ineqlabel{eq:trace-budget}
    \end{align}
\end{lemma}
\begin{proof}[Proof of \cref{lem:trace}]
    The case $u=0$ is immediate. For $t\leq u$, the definition of
    $\tau_{\mathrm{cov}}$ and the initialization give
    $\Sigma_t^{-1}\succeq\frac12\bar{\Sigma}_t^{-1}$, and hence
    $\Sigma_t\preceq2\bar{\Sigma}_t$.
    Since $\bar{H}_t\succeq0$, we have
    $\tr(\bar{H}_t\Sigma_t)\leq2\tr(\bar{H}_t\bar{\Sigma}_t)$,
    which proves the first inequality in \cref{eq:trace-budget}.
    To prove the second, let $L=\Lip(e)$.
    By cyclicity of the trace,
    \begin{align}
        \eta\sum_{t=1}^u
        \tr(\bar{H}_t\bar{\Sigma}_t)
        =
        \frac1\lambda\sum_{t=1}^u
        \tr\left(
        \lambda\eta\bar{\Sigma}_t^{1/2}\bar{H}_t
        \bar{\Sigma}_t^{1/2}
        \right).
        \eqlabel{eq:normalised-trace-sum}
    \end{align}
    For every $t\leq u$ and every unit vector $v$, distributional
    integration by parts, Cauchy--Schwarz, and
    \cref{eq:covariance-stopping-time} give
    \begin{align}
        v^\top\bar{\Sigma}_t^{1/2}\bar{H}_t
        \bar{\Sigma}_t^{1/2}v
         & \leq
        L\norm{\bar{\Sigma}_t^{1/2}v}
        \sqrt{v^\top\bar{\Sigma}_t^{1/2}
        \Sigma_t^{-1}\bar{\Sigma}_t^{1/2}v}
        \leq
        \sqrt{\frac32}\,L
        \norm{\bar{\Sigma}_t^{1/2}v}
        \leq
        \sqrt{\frac32}\,L\sigma.
        \ineqlabel{eq:normalised-H}
    \end{align}
    The last inequality uses
    $\bar{\Sigma}_t\preceq\bar{\Sigma}_1=\sigma^2I_d$.
    Therefore, by the parameter choices,
    \begin{align}
        \norm{\eta\lambda\bar{\Sigma}_t^{1/2}\bar{H}_t\bar{\Sigma}_t^{1/2}}_{\op}
        \leq
        \sqrt{\frac32}\lambda\eta L\sigma
        \leq
        \frac1{1000}
        \leq
        1\,.
    \end{align}
    Returning to
    \cref{eq:normalised-trace-sum} and applying
    $s\leq2\log(1+s)$ to eigenvalues of $\lambda\eta\bar{\Sigma}_t^{1/2}\bar{H}_t\bar{\Sigma}_t^{1/2}$
    which are in $[0,1]$, we obtain
    \begin{align}
        \eta\sum_{t=1}^u
        \tr(\bar{H}_t\bar{\Sigma}_t)
         & \leq
        \frac2\lambda\sum_{t=1}^u
        \log\det\left(
        I_d+\lambda\eta\bar{\Sigma}_t^{1/2}\bar{H}_t
        \bar{\Sigma}_t^{1/2}
        \right)                     \\
         & =
        \frac2\lambda\sum_{t=1}^u
        \log\frac{\det(\bar{\Sigma}_{t+1}^{-1})}
        {\det(\bar{\Sigma}_t^{-1})} \\
         & =
        \frac2\lambda
        \log\frac{\det(\bar{\Sigma}_{u+1}^{-1})}
        {\det(\bar{\Sigma}_1^{-1})} \\
         & \leq
        \frac{2d}{\lambda}
        \log\left(
        \sigma^2
        \norm{\bar{\Sigma}_{u+1}^{-1}}_{\op}
        \right)                     \\
         & =
        \frac{4d}{\lambda}
        \log\left(
        \sigma
        \norm{\bar{\Sigma}_{u+1}^{-1}}_{\op}^{1/2}
        \right)                     \\
         & \leq
        \frac{4d}{\lambda}
        \log\left(1+n\right)\,,
    \end{align}
    where the last inequality
    uses \cref{lem:predictable-spectrum},
    $\sigma\leq1$, and $u\leq n$. Multiplying this bound by two
    completes \cref{eq:trace-budget}.
\end{proof}


\begin{proof}[Proof of \cref{lem:stopped-trace}]
    Apply \cref{lem:trace} with $u=\tau-1$, which satisfies
    $0\leq u\leq n\wedge\tau_{\mathrm{cov}}$.
\end{proof}

\begin{lemma}[Uniform trace bound]
    \label{lem:uniform-trace}
    Under the parameter choices of \cref{thm:main}, with
    $\eps=\eps_0\leq1/2$, for every $\mu\in\R^d$ and
    $0\prec\Sigma\preceq2\sigma^2I_d$,
    \begin{align}
        0\leq\eps\tr\left(
        \langle D^2e,\gamma_{\mu,\Sigma}\rangle\Sigma
        \right)\leq5.
    \end{align}
    In particular, $\eps\tr(\bar H_t\Sigma_t)\leq5$ for every
    $1\leq t\leq n\wedge\tau_{\mathrm{cov}}$.
\end{lemma}
\begin{proof}[Proof of \cref{lem:uniform-trace}]
    Write $h=\eps e$ and
    $\K^\circ=\{a:\sup_{z\in\K}\ip{a}{z}\leq1\}$.
    We first show that $\nabla h(x)\in5\K^\circ$ at almost every $x$.
    At a differentiability point in $\operatorname{int}\K_\eps$, we have
    $x+\eps z\in\K$ for every $z\in\K$. Therefore
    \begin{align}
        \ip{\nabla h(x)}{z}
        =\eps\ip{\nabla f(x)}{z}
        \leq f(x+\eps z)-f(x)\leq1.
    \end{align}
    At a differentiability point outside $\K_\eps$, radial homogeneity
    gives $h(tx)+2=t(h(x)+2)$ for $t$ near one, and hence
    $\ip{\nabla h(x)}{x}=h(x)+2$. Applying the supporting-hyperplane
    inequality at $(1-\eps)z$, for $z\in\K$, gives
    \begin{align}
        (1-\eps)\ip{\nabla h(x)}{z}
        \leq h((1-\eps)z)+2\leq\eps+2.
    \end{align}
    Thus $\ip{\nabla h(x)}{z}\leq5$. The boundary of $\K_\eps$ has
    Lebesgue measure zero, so the two cases establish the claimed
    inclusion almost everywhere. In particular,
    $\ip{\nabla h(x)}{w}\leq5\pi(w)$ for every $w\in\R^d$.

    Let $X\sim\N(\mu,\Sigma)$ and $W=X-\mu$. Distributional integration
    by parts gives
    \begin{align}
        \eps\tr\left(\langle D^2e,\gamma_{\mu,\Sigma}\rangle\Sigma\right)
         & =\E\left[\ip{\nabla h(X)}{W}\right]
        \leq5\E[\pi(W)]
        \leq5\sqrt{2d}\,\sigma\E[\pi(U)]
        \leq5.
    \end{align}
    For the penultimate inequality, add an independent centered Gaussian
    with covariance $2\sigma^2I_d-\Sigma$, apply Jensen's inequality to
    the convex gauge, and use the independent radius and direction of
    the resulting isotropic Gaussian. Nonnegativity follows from
    convexity of $e$.
\end{proof}

\section{Feedback control}
\label{app:feedback}
\begin{lemma}[Single-round feedback moment]
    \label{lem:single-round-feedback}
    For every $t\leq n\wedge\tau_{\mathrm{cov}}$ satisfying
    $v(\mu_t)\leq\ell$,
    \begin{align}
        \E_t[Y_t^2]\leq\frac{8\ell^2}{\eps^2}.
        \ineqlabel{eq:single-round-feedback}
    \end{align}
\end{lemma}
\begin{proof}
    Write $W_t=X_t-\mu_t$ and $r_t=\pi_\eps(X_t)=1+v(X_t)$.
    The covariance and shell conditions give
    $\Sigma_t\preceq2\sigma^2I_d$ and $v(\mu_t)\leq\ell$.
    By subadditivity of the gauge and $\eps\leq1/2$,
    \begin{align}
        r_t\leq1+\ell+2\pi(W_t).
    \end{align}
    Conditionally on $\F_{t-1}$, $W_t$ is centered Gaussian with
    covariance $\Sigma_t$. Thus
    \cref{lem:gaussian-gauge-fourth-moment} and Minkowski's inequality give
    \begin{align}
        \E_t[r_t^2]^{1/2}
        \leq1+\ell+2\E_t[\pi(W_t)^2]^{1/2}
        \leq\ell+7.
    \end{align}
    The noise assumption \cref{eq:noise} gives conditional mean zero
    and conditional second moment at most one. Since
    $Y_t=e(X_t)+r_t\xi_t$ and
    $0\leq\eps e(X_t)\leq(2+\eps)r_t\leq(5/2)r_t$,
    \begin{align}
        \eps^2\E_t[Y_t^2]
        =\eps^2\E_t\left[e(X_t)^2+r_t^2\E_t[\xi_t^2\mid X_t]\right]
        \leq7\E_t[r_t^2]
        \leq7(\ell+7)^2
        \leq8\ell^2,
    \end{align}
    where the last inequality uses $\ell\geq640$.
\end{proof}

\begin{lemma}[Gaussian gauge fourth moment]
    \label{lem:gaussian-gauge-fourth-moment}
    Let $\sigma=(M\sqrt{2d})^{-1}$ and $W\sim\N(0,\Sigma)$ with
    $0\preceq\Sigma\preceq2\sigma^2I_d$. Then
    \begin{align}
        \E[\pi(W)^4]^{1/4}\leq3.
    \end{align}
\end{lemma}
\begin{proof}
    Adding independent centered Gaussian noise to make the covariance
    $2\sigma^2I_d$ and applying Jensen's inequality to the convex gauge
    can only increase its expectation. Decomposing the resulting
    isotropic Gaussian into its independent radius and direction gives
    \begin{align}
        \E[\pi(W)]
        \leq\sqrt{2d}\,\sigma\E[\pi(U)]\leq1.
    \end{align}
    Since $\mathbb B_1\subseteq\K$, the gauge $\pi$ is $1$-Lipschitz.
    Also $\Sigma\preceq2\sigma^2I_d\preceq I_d$, so Gaussian
    concentration gives
    \begin{align}
        \Pp\left(|\pi(W)-\E[\pi(W)]|>r\right)
        \leq2e^{-r^2/2}\qquad(r>0).
    \end{align}
    Integrating this tail bound yields
    \begin{align}
        \E\left[|\pi(W)-\E[\pi(W)]|^4\right]
        \leq8\int_0^\infty r^3e^{-r^2/2}\,dr=16.
    \end{align}
    Minkowski's inequality therefore gives
    \begin{align}
        \E[\pi(W)^4]^{1/4}
        \leq\E[\pi(W)]
        +\E\left[|\pi(W)-\E[\pi(W)]|^4\right]^{1/4}
        \leq1+2=3.
    \end{align}
\end{proof}

\begin{lemma}[Conditional feedback cap]
    \label{lem:conditional-feedback-cap}
    Assume the parameter choices of \cref{thm:main} with
    $\eps=\eps_0\leq1/2$. For every $t\leq n\wedge\tau_{\mathrm{cov}}$
    satisfying $v(\mu_t)\leq\ell$,
    \begin{align}
        \Pp_t\left(|\eps Y_t|>4\ell^{3/2}\right)\leq3e^{-\ell}.
        \ineqlabel{eq:conditional-feedback-cap}
    \end{align}
\end{lemma}
\begin{proof}
    Write $W_t=X_t-\mu_t$ and $r_t=\pi_\eps(X_t)=1+v(X_t)$.
    The covariance comparison gives
    $\Sigma_t\preceq2\sigma^2I_d\preceq I_d$, and
    \cref{lem:gaussian-gauge-fourth-moment} gives
    $\E_t[\pi(W_t)]\leq3$. Since $\pi$ is $1$-Lipschitz,
    Gaussian concentration and the noise assumption \cref{eq:noise}
    imply
    \begin{align}
        \Pp_t\left(\pi(W_t)>3+2\sqrt\ell\right)\leq e^{-2\ell},
        \qquad
        \Pp_t\left(|\xi_t|>\sqrt\ell\right)\leq2e^{-\ell}.
    \end{align}
    Outside these two events, subadditivity of the gauge gives
    \begin{align}
        v(X_t)\leq v(\mu_t)+2\pi(W_t)
        \leq\ell+6+4\sqrt\ell\leq2\ell,
        \qquad r_t\leq3\ell.
    \end{align}
    Therefore the extension formula and $f(A_t)\in[0,1]$ give
    \begin{align}
        |\eps Y_t|
        \leq\eps r_t(1+|\xi_t|)+2v(X_t)
        \leq\frac{11}{2}\ell+\frac32\ell^{3/2}
        \leq4\ell^{3/2},
    \end{align}
    using $\eps\leq1/2$ and $\ell\geq640$.
    A union bound proves the claim.
\end{proof}

\begin{lemma}[Weighted Gaussian tail]
    \label{lem:weighted-gaussian-tail}
    Assume the parameter choices of \cref{thm:main} with $\eps=\eps_0\leq1/2$.
    For a fixed comparator $x_{\circ}$ and $1\leq t\leq n\wedge\tau_{\mathrm{cov}}$,
    define
    \begin{align}
        R_t(x_{\circ})=\normm{x_{\circ}-\mu_t}{\Sigma_t^{-1}},
        \qquad
        G_t(x_{\circ})=
        \ip{\Sigma_t^{-1/2}(x_{\circ}-\mu_t)/R_t(x_{\circ})}{Z_t}
    \end{align}
    when $R_t(x_{\circ})>0$, and set $G_t(x_{\circ})=0$ otherwise.
    For every such round with $t<\tau_{\mathrm{shell}}$,
    \begin{align}
        \E_t\left[Y_t^2(G_t(x_{\circ})^2+1)^2
            \1_{\{|G_t(x_{\circ})|>2\sqrt\ell\}}\right]
        \leq\frac{e^{-\ell/4}}{\eps^2}
        \leq e^{-\ell/8}.
    \end{align}
    Consequently, for every interval
    $1\leq s\leq u\leq n\wedge\tau_{\mathrm{cov}}$ with $u<\tau_{\mathrm{shell}}$,
    \begin{align}
        \sum_{t=s}^{u}\E_t\left[Y_t^2(G_t(x_{\circ})^2+1)^2
            \1_{\{|G_t(x_{\circ})|>2\sqrt\ell\}}\right]
        \leq e^{-\ell/16}.
    \end{align}
\end{lemma}
\begin{proof}
    If $R_t(x_{\circ})=0$, then $G_t(x_{\circ})=0$ and the tail term vanishes.
    Otherwise, $G_t(x_{\circ})$ is conditionally standard Gaussian.
    The noise assumption \cref{eq:noise} and
    $0\leq\eps e(X_t)\leq(5/2)\pi_\eps(X_t)$ give
    \begin{align}
        \E_t[Y_t^2\mid X_t]
        =e(X_t)^2+\pi_\eps(X_t)^2\E_t[\xi_t^2\mid X_t]
        \leq\frac7{\eps^2}\pi_\eps(X_t)^2.
    \end{align}
    The covariance and shell conditions, subadditivity of the gauge,
    and \cref{lem:gaussian-gauge-fourth-moment} imply
    \begin{align}
        \E_t[\pi_\eps(X_t)^4]^{1/4}
        \leq1+v(\mu_t)+2\E_t[\pi(X_t-\mu_t)^4]^{1/4}
        \leq\ell+7,
    \end{align}
    so $\E_t[\pi_\eps(X_t)^4]^{1/2}\leq4\ell^2$.
    Gaussian moments and the Gaussian tail bound give
    \begin{align}
        \E_t[(G_t(x_{\circ})^2+1)^8]
         & \leq2^7(15!!+1),
         &
        \Pp_t\left(|G_t(x_{\circ})|>2\sqrt\ell\right)
         & \leq2e^{-2\ell}.
    \end{align}
    Since $G_t(x_{\circ})$ is determined by $X_t$, conditioning first on
    $X_t$ and applying Cauchy--Schwarz twice yields
    \begin{align}
         & \E_t\left[Y_t^2(G_t(x_{\circ})^2+1)^2
            \1_{\{|G_t(x_{\circ})|>2\sqrt\ell\}}\right]
        \nonumber                                         \\
         & \quad\leq\frac7{\eps^2}
        \E_t[\pi_\eps(X_t)^4]^{1/2}
        \E_t[(G_t(x_{\circ})^2+1)^8]^{1/4}
        \Pp_t\left(|G_t(x_{\circ})|>2\sqrt\ell\right)^{1/4}
        \nonumber                                         \\
         & \qquad\leq\frac{5600\ell^2}{\eps^2}e^{-\ell/2}
        \leq\frac{e^{-\ell/4}}{\eps^2},
    \end{align}
    where the last inequality uses $\ell\geq640$.
    The parameter choices and $M\geq d^{-1/2}$ give
    $\eps^{-2}=n/(10^4d^4M^2\ell^{16})\leq n\leq e^{\ell/128}$,
    and hence $\eps^{-2}e^{-\ell/4}\leq e^{-\ell/8}$.
    Summing over the interval proves the second claim, since
    \begin{align}
        (u-s+1)e^{-\ell/8}
        \leq ne^{-\ell/8}
        \leq e^{-15\ell/128}
        \leq e^{-\ell/16}.
    \end{align}
\end{proof}

\begin{lemma}[Feedback moment self-bound]
    \label{lem:feedback}
    Suppose that $\eps=\eps_0\leq1/2$, and let $\tau$ be the common
    stopping time from \cref{sec:stochastic-control}. For every integer
    $0\leq u<\tau$,
    \begin{align}
        \sum_{t=1}^u\E_t[Y_t^2]
        \leq4u+\frac{32\ell}{\eps}\operatorname{eReg}_{1:u}(x_\star)
        +\frac{32ue^{-\ell}}{\eps^2}.
        \ineqlabel{eq:feedback-self-bound}
    \end{align}
\end{lemma}
\begin{proof}
    The case $u=0$ is immediate. For $1\leq t\leq u<\tau$, write
    \begin{align}
        r_t=\pi_\eps(X_t)=1+v(X_t),
        \qquad W_t=X_t-\mu_t.
    \end{align}
    The shell condition and subadditivity of the gauge give
    \begin{align}
        v(X_t)
         & \leq v(\mu_t)+\frac{\pi(W_t)}{1-\eps}
        \leq \ell+2\pi(W_t).
        \ineqlabel{eq:feedback-v-minkowski}
    \end{align}
    The stopped covariance comparison gives
    $\Sigma_t\preceq2\sigma^2I_d\preceq I_d$, and
    \cref{lem:gaussian-gauge-fourth-moment} gives
    $\E_t[\pi(W_t)]\leq3$.
    Since $\pi$ is $1$-Lipschitz, Gaussian concentration gives
    \begin{align}
        \Pp_t\left(\pi(W_t)>\E_t[\pi(W_t)]+s\right)
        \leq\exp\left(-\frac{s^2}{2\|\Sigma_t\|_{\op}}\right)
        \leq e^{-s^2/2}.
        \ineqlabel{eq:feedback-minkowski-tail}
    \end{align}
    In particular,
    \begin{align}
        \Pp_t(2\pi(W_t)>s)\leq e^{-s^2/32}
        \qquad(s\geq12).
        \ineqlabel{eq:feedback-v-tail}
    \end{align}
    Conditioning on $X_t$ first, the noise assumption \cref{eq:noise} gives
    \begin{align}
        \E_t[Y_t^2\mid X_t]
         & =e(X_t)^2+r_t^2\E_t[\xi_t^2\mid X_t]
        \leq e(X_t)^2+r_t^2.
        \ineqlabel{eq:Y2-conditional-simple}
    \end{align}
    Using $e(X_t)\leq1+3v(X_t)/\eps$ and $r_t=1+v(X_t)$, we obtain
    \begin{align}
        \sum_{t=1}^u\E_t[Y_t^2]
        \leq4u+\frac{32}{\eps^2}
        \sum_{t=1}^u\E_t[v(X_t)^2].
        \ineqlabel{eq:feedback-start}
    \end{align}
    Set $a=16\sqrt\ell$. Splitting at $\ell+a$, and then integrating
    \cref{eq:feedback-v-tail}, yields
    \begin{align}
        \E_t[v(X_t)^2]
         & \leq(\ell+a)\E_t[v(X_t)]
        +\E_t\left[v(X_t)^2\1\{v(X_t)>\ell+a\}\right]
        \\
         & \leq2\ell\E_t[v(X_t)]
        +\left(2\ell^2+2a^2+64\right)e^{-a^2/32}
        \\
         & \leq2\ell\E_t[v(X_t)]+e^{-\ell}.
        \ineqlabel{eq:feedback-v2-split}
    \end{align}
    Indeed, the tail event implies $2\pi(W_t)>a$, and on that event
    $v(X_t)^2\leq2\ell^2+2(2\pi(W_t))^2$. Also $a\leq \ell$ and
    $(2\ell^2+512\ell+64)e^{-8\ell}\leq e^{-\ell}$.

    Finally, $A_t\in\K_\eps$ implies $f(A_t)\geq f(x_\star)\geq0$.
    Therefore the definition of the extension gives the pointwise bound
    \begin{align}
        e(X_t)-e(x_\star)
         & =r_tf(A_t)-f(x_\star)+\frac2\eps v(X_t)
        \geq\frac2\eps v(X_t).
        \ineqlabel{eq:extension-violation}
    \end{align}
    Combining \cref{eq:feedback-start,eq:feedback-v2-split,eq:extension-violation}
    proves \cref{eq:feedback-self-bound}.

\end{proof}

\section{Concentration}
\label{app:concentration}
Throughout this section, $\tau$ denotes the common stopping time from
\cref{sec:stochastic-control}. For the proofs, set
\begin{align}
    y_t=\eps Y_t,
    \qquad N=n\eps^2=10^4d^4M^2\ell^{16}.
    \eqlabel{eq:normalised-feedback}
\end{align}
Conditionally on $\F_{t-1}$, the standardized sample
$Z_t=\Sigma_t^{-1/2}(X_t-\mu_t)$ has law $\N(0,I_d)$.
Stopped increments are extended by zero after round $n\wedge\tau$.
Since $\{t\leq\tau\}\in\F_{t-1}$, retaining round $\tau$ preserves
predictability; any additional activity cutoffs used below are also checked
before sampling the round.

\begin{proof}[Proof of \cref{lem:global-event}]
    Fix $1\leq t\leq n$ and work conditionally on $\F_{t-1}$ on
    $\{t\leq\tau\}$. Since $\tau\leq\tau_{\mathrm{cov}}$ and
    the covariance stopping rule tests the update at $t+1$, we have
    \begin{align}
        0\prec\Sigma_t\preceq2\bar\Sigma_t
        \preceq2\sigma^2I_d\preceq I_d.
    \end{align}
    Here $\bar H_i\succeq0$ makes $\bar\Sigma_t\preceq\bar\Sigma_1$,
    and $2\sigma^2=(dM^2)^{-1}\leq1$.

    First, $\norm{Z_t}^2$ is conditionally $\chi_d^2$.
    By \cref{lem:chi-square-upper-tail} with $x=\ell$, and
    $d+2\sqrt{d\ell}+2\ell\leq2d+3\ell\leq4d\ell$ since $\ell\geq640$,
    \begin{align}
        \Pp_t\left(\norm{Z_t}^2>4d\ell\right)\leq e^{-\ell}.
    \end{align}
    Second, conditional Markov's inequality and \cref{eq:noise} give
    \begin{align}
        \Pp_t\left(|\xi_t|>\sqrt\ell\right)
        \leq e^{-\ell}\E_t[\exp(\xi_t^2)]\leq2e^{-\ell}.
    \end{align}
    Third, write $W_t=X_t-\mu_t$. By
    \cref{lem:gaussian-gauge-fourth-moment},
    $\E_t[\pi(W_t)]\leq\E_t[\pi(W_t)^4]^{1/4}\leq3$.
    The map $z\mapsto\pi(\Sigma_t^{1/2}z)$ is $1$-Lipschitz,
    because $\pi$ is $1$-Lipschitz and $\Sigma_t\preceq I_d$.
    Gaussian concentration therefore gives
    \begin{align}
        \Pp_t\left(\pi(W_t)>3+2\sqrt\ell\right)
        \leq
        \Pp_t\left(\pi(W_t)>\E_t[\pi(W_t)]+2\sqrt\ell\right)
        \leq e^{-2\ell}.
    \end{align}

    Since $\{t\leq\tau\}\in\F_{t-1}$, multiplying each conditional
    bound by $\1_{\{t\leq\tau\}}$ and taking expectations bounds
    the probability that the corresponding cutoff fails at a retained
    round. A union bound over the three cutoffs and $1\leq t\leq n$ yields
    \begin{align}
        \Pp(\mathcal E_{\mathrm{global}}^c)
        \leq n\left(3e^{-\ell}+e^{-2\ell}\right)
        \leq4ne^{-\ell}
        =4n\left(\frac{\delta}{128nd}\right)^{128}
        \leq\frac{\delta}{32d}\leq\frac\delta4.
    \end{align}
\end{proof}

\begin{proof}[Proof of \cref{lem:feedback-square-sum}]
    For $1\leq t\leq n$, let
    \begin{align}
        U_t=\begin{cases}
                Y_t & \text{if }t<\tau, \\
                0   & \text{otherwise},
            \end{cases}
        \qquad
        I_t=\1\left\{t\leq\tau,\ \sum_{i=1}^t\E_i[Y_i^2]\leq S\right\},
    \end{align}
    where $I_t=0$ after $\tau$. The sequence $(U_t)_{t=1}^n$ is
    adapted. Since $\{t\leq\tau\}\in\F_{t-1}$ and the conditional
    moment budget through round $t$ is $\F_{t-1}$-measurable,
    $I_t$ is predictable. Also $\tau\leq\tau_{\mathrm{moment}}$ gives
    $\1_{\{t<\tau\}}\leq I_t$, and hence
    \begin{align}
        \sum_{t=1}^n\E_t[U_t^2]
        \leq\sum_{t=1}^n I_t\E_t[Y_t^2]\leq S.
    \end{align}
    The last inequality follows because the cumulative conditional
    second moments are nondecreasing and $I_t$ retains only rounds
    for which their sum through that round is at most $S$.
    Terms with $I_t=0$ are interpreted as zero.

    On $\mathcal E_{\mathrm{global}}$,
    \cref{eq:energy-sample-gauge-bound} gives
    $\max_{1\leq t\leq n}U_t^2\leq b$ with
    $b=16\ell^3/\eps^2$. Take $z=\ell/64$. Since
    \begin{align}
        4bz=\frac{\ell^4}{\eps^2}
        \leq\frac{4d\ell^5}{\eps^2},
    \end{align}
    the event $\mathcal E_{\mathrm{global}}\cap\mathcal E_Y^c$
    is contained in
    \begin{align}
        \left\{
        \sum_{t=1}^nU_t^2\geq2S+4bz,\quad
        \sum_{t=1}^n\E_t[U_t^2]\leq S,\quad
        \max_{1\leq t\leq n}U_t^2\leq b
        \right\}.
    \end{align}
    Applying \cref{lem:freedman-nonnegative} with $T=n$ therefore yields
    \begin{align}
        \Pp\left(\mathcal E_{\mathrm{global}}\cap\mathcal E_Y^c\right)
        \leq2e^{-\ell/64}
        =2\left(\frac{\delta}{128nd}\right)^2
        \leq\frac\delta8.
    \end{align}
\end{proof}

\begin{proof}[Proof of \cref{lem:quadratic-concentration}]
    We first sharpen the second-moment bound on
    $\mathcal E_{\mathrm{global}}\cap\mathcal E_G\cap\mathcal E_Y$.
    For $t<\tau$, the focus condition gives $R_t^2\leq F$ and
    $\lambda R_t\leq1/(128\sqrt\ell)$. Whenever $|G_t|\leq2\sqrt\ell$,
    we have $|G_t^2-1|\leq4\ell$, and hence
    \begin{align}
        \left|R_tG_t+\frac\lambda2R_t^2(G_t^2-1)\right|
        \leq R_t\left(2\sqrt\ell+2\lambda R_t\ell\right)
        \leq\left(2+\frac1{64}\right)R_t\sqrt\ell
        \leq\sqrt{5F\ell}.
    \end{align}
    Thus $q_t(x_\star)^2\leq5F\ell Y_t^2$ on this cutoff.
    For the conditional second moment, use this sharper bound on
    $\{|G_t|\leq2\sqrt\ell\}$ and retain
    \cref{eq:quadratic-estimate-decomposition} on its complement:
    \begin{align}
        \E_t[q_t(x_\star)^2]
        \leq5F\ell\E_t[Y_t^2]
        +4F\E_t\left[Y_t^2(G_t^2+1)^2
            \1_{\{|G_t|>2\sqrt\ell\}}\right].
    \end{align}
    Combining \cref{eq:martingale-variance,lem:weighted-gaussian-tail},
    the definition of $\mathcal E_Y$, and the moment stopping rule gives
    \begin{align}
        \sum_{t=1}^{\tau-1}\left(D_t^2+\E_t[D_t^2]\right)
         & \leq10F\ell\sum_{t=1}^{\tau-1}
        \left(Y_t^2+\E_t[Y_t^2]\right)+8Fe^{-\ell/16}         \\
         & \leq10F\ell\left(3S+\frac{4d\ell^5}{\eps^2}\right)
        +8Fe^{-\ell/16}                                       \\
         & \leq48F\ell\left(S+\frac{d\ell^5}{\eps^2}\right).
        \ineqlabel{eq:quadratic-sharp-variance}
    \end{align}
    The last inequality uses $d\ell^5/\eps^2\geq1$ and $\ell\geq640$.

    To apply the martingale inequality at the random endpoint $\tau-1$,
    define, for $1\leq t\leq n$, the predictable indicators
    \begin{align}
        I_t=\1\left\{t\leq\tau,\ R_t^2\leq F,\ v(\mu_t)\leq\ell,\
        \sum_{i=1}^t\E_i[Y_i^2]\leq S\right\},
        \qquad U_t=I_tD_t,
    \end{align}
    interpreting both quantities as zero after $\tau$.
    All the cutoffs in $I_t$ are $\F_{t-1}$-measurable, so
    $(U_t)_{t=1}^n$ is a martingale-difference sequence and
    $\E_t[U_t^2]=I_t\E_t[D_t^2]$.
    Moreover, $I_t=1$ for every $t<\tau$.
    Apply \cref{lem:self-normalised-martingale} to $(U_t)_{t=1}^m$
    for each deterministic $1\leq m\leq n$, with
    \begin{align}
        z=\sqrt{F\ell^3\left(S+\frac{d\ell^5}{\eps^2}\right)},
        \qquad
        v=48F\ell\left(S+\frac{d\ell^5}{\eps^2}\right).
    \end{align}
    A union bound shows that the probability of any prefix satisfying
    $|\sum_{t=1}^mU_t|>z$ and
    $\sum_{t=1}^m(U_t^2+\E_t[U_t^2])\leq v$ is at most
    $2n\exp(-\ell^2/96)$.
    If $\tau=1$, then $\mathcal E_q$ holds trivially. Otherwise,
    on $\mathcal E_{\mathrm{global}}\cap\mathcal E_G\cap\mathcal E_Y
        \cap\mathcal E_q^c$, take $m=\tau-1$.
    The retained sums coincide with the original sums through $\tau-1$,
    and \cref{eq:quadratic-sharp-variance} bounds their combined
    second moments by $v$. Therefore this event has probability at most
    $2n\exp(-\ell^2/96)$.

    Finally, \cref{lem:feedback-square-sum} bounds the contribution
    from $\mathcal E_Y^c$, giving
    \begin{align}
        \Pp\left(\mathcal E_{\mathrm{global}}\cap\mathcal E_G\cap\mathcal E_q^c\right)
         & \leq2e^{-\ell/64}+2ne^{-\ell^2/96} \\
         & \leq4e^{-\ell/64}\leq\frac\delta8.
    \end{align}
    Here $\log n\leq\ell/128$ and $\ell\geq640$ imply
    $ne^{-\ell^2/96}\leq e^{-\ell/64}$, while
    $4e^{-\ell/64}=4(\delta/(128nd))^2\leq\delta/8$.
\end{proof}

\begin{proof}[Proof of \cref{lem:local-quadratic-concentration}]
    Fix the deterministic interval $[s,u]$. For $s\leq t\leq u$, let
    \begin{align}
        I_t=\1\left\{t\leq\tau,\ \tilde R_t^2\leq F_L,
        \ v(\mu_t)\leq\ell\right\},
        \qquad U_t=I_t\tilde D_t,
    \end{align}
    interpreting all retained quantities as zero when $\tau<s$ or
    $t>\tau$. Since $\mu_s$ is $\F_{s-1}$-measurable, $I_t$ is
    predictable and $(U_t)_{t=s}^u$ is a martingale-difference sequence.
    Moreover, $I_t=1$ for every $s\leq t<\tilde\tau$. If
    $\tilde\tau\leq s$, the local event holds vacuously.

    On each retained round, \cref{lem:single-round-feedback} gives
    $\E_t[Y_t^2]\leq8\ell^2/\eps^2$. Thus
    \begin{align}
        \sum_{t=s}^u\E_t[(I_tY_t)^2]
        \leq\frac{8L\ell^2}{\eps^2}.
    \end{align}
    On $\mathcal E_{\mathrm{global}}$, the calculation in
    \cref{eq:energy-sample-gauge-bound} also gives
    $|I_tY_t|\leq4\ell^{3/2}/\eps$.
    Apply \cref{lem:freedman-nonnegative} to $(I_tY_t)_{t=s}^u$
    with $b=16\ell^3/\eps^2$, moment budget $8L\ell^2/\eps^2$,
    and $z=\ell/64$. Since $4bz=\ell^4/\eps^2\leq4d\ell^5/\eps^2$,
    outside an event whose intersection with $\mathcal E_{\mathrm{global}}$
    has probability at most $2e^{-\ell/64}$, we have
    \begin{align}
        \sum_{t=s}^u I_tY_t^2
        \leq\frac{16L\ell^2}{\eps^2}
        +\frac{4d\ell^5}{\eps^2}.
    \end{align}

    On a retained round with $\tilde R_t>0$, define
    \begin{align}
        \tilde G_t
        =\ip{\Sigma_t^{-1/2}(\mu_s-\mu_t)/\tilde R_t}{Z_t},
    \end{align}
    and set $\tilde G_t=0$ otherwise. Its direction is predictable,
    so a conditional Gaussian tail bound and a union bound show that
    $|\tilde G_t|\leq2\sqrt\ell$ on every retained round except on
    an event of probability at most $2Le^{-2\ell}$.
    Since $F_L\leq1/(128^2\lambda^2\ell)$, the same calculation as
    in the proof of \cref{lem:quadratic-concentration} gives
    \begin{align}
        q_t(\mu_s)^2\leq5F_L\ell Y_t^2
    \end{align}
    on these rounds. The conditional tail estimate in
    \cref{lem:weighted-gaussian-tail} applies to $\mu_s$ as well:
    its proof only requires the comparator to be
    $\F_{t-1}$-measurable. Consequently, on the intersection of the
    preceding events and $\mathcal E_{\mathrm{global}}$,
    \begin{align}
        \sum_{t=s}^u\left(U_t^2+\E_t[U_t^2]\right)
         & \leq10F_L\ell\sum_{t=s}^u
        I_t\left(Y_t^2+\E_t[Y_t^2]\right)+8F_Le^{-\ell/16}
        \\
         & \leq48F_L\ell\left(\frac{8L\ell^2}{\eps^2}
        +\frac{d\ell^5}{\eps^2}\right).
    \end{align}
    This also bounds the combined second moments of every prefix.

    Apply \cref{lem:self-normalised-martingale} to $(U_t)_{t=s}^j$
    for each deterministic $s\leq j\leq u$, with
    \begin{align}
        z=\sqrt{F_L\ell^3\left(\frac{8L\ell^2}{\eps^2}
            +\frac{d\ell^5}{\eps^2}\right)},
        \qquad
        v=48F_L\ell\left(\frac{8L\ell^2}{\eps^2}
        +\frac{d\ell^5}{\eps^2}\right).
    \end{align}
    A union bound over the $L$ prefixes gives probability at most
    $2Le^{-\ell^2/96}$. For $j<\tilde\tau$, their sums coincide
    with $\sum_{t=s}^j\tilde D_t$. Combining the three failure
    probabilities therefore yields
    \begin{align}
        \Pp\left(\mathcal E_{\mathrm{global}}\cap
        \mathcal E_{q,\mathrm{loc}}^c\right)
         & \leq2e^{-\ell/64}+2Le^{-2\ell}+2Le^{-\ell^2/96}
        \\
         & \leq6e^{-\ell/64}
        =6\left(\frac{\delta}{128nd}\right)^2
        \leq\frac{\delta}{8n^2}.
    \end{align}
    Here we used $L\leq n$, $\log n\leq\ell/128$, and
    $\ell\geq640$. A union bound over at most $n^2$ short intervals
    proves the final assertion.
\end{proof}

The following two-sided form of Freedman's inequality follows from
\citet[Theorem~3.10]{bercu2015concentration} applied to a martingale and its
negative.
\begin{lemma}[Freedman's inequality]
    \label{lem:freedman}
    Let $(D_t)_{t=1}^n$ be a square-integrable\\
    martingale-difference sequence such that, for some $b>0$,
    $|D_t|\leq b$ almost surely for every $t$. Then, for every $s,z>0$,
    \begin{align}
        \Pp\left(
        \left|\sum_{t=1}^n D_t\right|
        \geq sz+\frac{bz^2}{3},
        \quad
        \sum_{t=1}^n\E_t[D_t^2]\leq s^2
        \right)
        \leq
        2\exp\left(-\frac{z^2}{2}\right).
    \end{align}
\end{lemma}

\begin{lemma}[Weighted Gaussian score moments]
    \label{lem:gaussian-score-localisation}
    Suppose that $\eps=\eps_0\leq1/2$. On a round $t\leq n\wedge\tau$
    satisfying $v(\mu_t)\leq\ell$, let
    $h_t$ be any $\F_{t-1}$-measurable unit vector and put
    $G_t=\ip{h_t}{Z_t}$. Then
    \begin{align}
        \E_t[y_t^2(G_t^2+1)^2]
        \leq25\ell^2\E_t[y_t^2]+e^{-\ell/4}.
        \ineqlabel{eq:H-score-second-moment}
    \end{align}
\end{lemma}
\begin{proof}
    Write $r_t=\pi_\eps(X_t)$.
    On $\mathcal A_t=\{|G_t|\leq2\sqrt\ell\}$, we have
    $(G_t^2+1)^2\leq25\ell^2$. Thus it suffices to bound the integral on
    $\overline{\mathcal A_t}$. The noise assumption \cref{eq:noise} and
    $0\leq\eps e(X_t)\leq(5/2)r_t$ give
    $\E_t[y_t^2\mid X_t]\leq7r_t^2$. Consequently,
    \begin{align}
        \E_t[y_t^2(G_t^2+1)^2\1_{\overline{\mathcal A_t}}]
         & \leq7\E_t[r_t^2(G_t^2+1)^2\1_{\overline{\mathcal A_t}}]
        \\
         & \leq7\E_t[r_t^4]^{1/2}
        \E_t[(G_t^2+1)^4\1_{\overline{\mathcal A_t}}]^{1/2}.
    \end{align}
    The shell bound and $\Sigma_t\preceq2\sigma^2I_d$ give
    $\E_t[r_t^4]^{1/4}\leq\ell+7$ by
    \cref{lem:gaussian-gauge-fourth-moment} and Minkowski's inequality.
    Thus the first factor is at most $(\ell+7)^2$.
    For the second, Cauchy--Schwarz and
    the standard Gaussian sixteenth moment give
    \begin{align}
        \E_t[(G_t^2+1)^4\1_{\overline{\mathcal A_t}}]^{1/2}
         & \leq\E_t[(G_t^2+1)^8]^{1/4}\Pp_t(\overline{\mathcal A_t})^{1/4}
        \\
         & \leq\bigl(2^8(15!!+1)\bigr)^{1/4}e^{-\ell/2}
        \leq200e^{-\ell/2}.
    \end{align}
    Consequently the omitted integral is at most
    \begin{align}
        1400(\ell+7)^2e^{-\ell/2}\leq e^{-\ell/4},
    \end{align}
    using $\ell\geq640$. No independence between the noise and the
    sampled action is used.
\end{proof}

We will truncate only in the proof, never in the algorithm. Define
\begin{align}
    C_t=\{|y_t|\leq4\ell^{3/2},\ \norm{Z_t}^2\leq3(d+\ell)\}.
    \eqlabel{eq:common-truncation}
\end{align}
On a round $t\leq n\wedge\tau$ satisfying $v(\mu_t)\leq\ell$,
\cref{lem:conditional-feedback-cap} and the
chi-square bound below give $\Pp_t(\overline{C_t})\leq4e^{-\ell}$.
The following chi-square upper-tail inequality follows from
\citet[Lemma~1 and Eq.~(4.3)]{laurent2000adaptive}.
\begin{lemma}[Chi-square upper-tail bound]
    \label{lem:chi-square-upper-tail}
    Let $V\sim\chi_d^2$. Then, for every $x>0$,
    \begin{align}
        \Pp\left(
        V\geq d+2\sqrt{dx}+2x
        \right)
        \leq
        \exp(-x).
        \ineqlabel{eq:chi-square-upper-tail}
    \end{align}
\end{lemma}

We use the following consequence of \cref{lem:freedman}: centered
increments bounded by $b$ with total conditional variance at most $w$
satisfy, for every $z>0$,
\begin{align}
    \Pp\left(\left|\sum_tD_t\right|>\sqrt{2wz}+2bz\right)
    \leq2e^{-z}.
    \ineqlabel{eq:freedman-convenient}
\end{align}
The bounds $b$ and $w$ are deterministic; the Hessian argument below
handles random directional scales by truncating at deterministic
dyadic levels.

\begin{lemma}[Stopped Hessian concentration]
    \label{lem:hessian-concentration}
    Suppose that $\eps=\eps_0\leq1/2$. There is an event $\mathcal E_H$
    satisfying
    \begin{align}
        \Pp(\overline{\mathcal E}_H)
        \leq ne^{-\ell}+2e^{-5d\ell/128}
    \end{align}
    such that, on $\mathcal E_{\mathrm{global}}\cap\mathcal E_H$,
    if $\tau\leq n$ and
    \begin{align}
        v(\mu_t)\leq\ell\quad(1\leq t\leq\tau),
        \qquad
        \sum_{t=1}^{\tau}\E_t[Y_t^2]\leq16n,
    \end{align}
    then
    \begin{align}
        -100\ell^{3/2}\sqrt{dn}\,\bar{\Sigma}_{\tau+1}^{-1}
        \preceq\sum_{t=1}^{\tau}(H_t-\bar{H}_t)
        \preceq100\ell^{3/2}\sqrt{dn}\,\bar{\Sigma}_{\tau+1}^{-1}.
        \ineqlabel{eq:H-concentration}
    \end{align}
\end{lemma}
\begin{proof}
    To obtain a deterministic conditional-variance bound, retain only
    rounds satisfying the predictable cutoffs
    \begin{align}
        I_t=\1\left\{
        t\leq\tau,\quad v(\mu_t)\leq\ell,\quad
        \sum_{i=1}^t\E_i[Y_i^2]\leq16n
        \right\},\qquad 1\leq t\leq n.
        \eqlabel{eq:H-predictable-cutoffs}
    \end{align}
    All three conditions are $\F_{t-1}$-measurable. Quantities multiplied
    by $I_t$ are interpreted as zero when $I_t=0$, so no sample after
    $n\wedge\tau$ is needed. Monotonicity of the cumulative conditional
    second moments gives, on every path,
    \begin{align}
        \sum_{t=1}^n I_t\E_t[y_t^2]\leq16N.
        \ineqlabel{eq:H-retained-moment-budget}
    \end{align}
    Under the hypotheses of the lemma, $I_t=1$ for every $t\leq\tau$.
    It therefore suffices to control
    \begin{align}
        T=\eps\sum_{t=1}^n I_t(H_t-\bar{H}_t),
        \qquad A=\bar{\Sigma}_{(n\wedge\tau)+1}^{-1}.
    \end{align}

    We also require $\norm{Z_t}^2\leq3(d+\ell)$ for all
    $t\leq n\wedge\tau$. By \cref{lem:chi-square-upper-tail}, this
    extra Gaussian event fails with probability at most $ne^{-\ell}$.
    On $\mathcal E_{\mathrm{global}}$, the shell cutoff and the
    calculation in \cref{eq:energy-sample-gauge-bound} give
    $|y_t|\leq4\ell^{3/2}$ whenever $I_t=1$, including at $t=\tau$.
    Thus $C_t$ holds on every retained round if this extra Gaussian
    event also holds.

    Let $\mathcal U$ be a deterministic
    $(512dn^2)^{-1}$-net of the Euclidean unit sphere. The volumetric
    covering bound gives
    \begin{align}
        \log|\mathcal U|
        \leq d\log(1+1024dn^2)\leq\frac{d\ell}{64}.
        \ineqlabel{eq:H-net-size}
    \end{align}
    For $u\in\mathcal U$ and $t\leq n\wedge\tau$, define
    \begin{align}
        a_{t,u}=u^\top\Sigma_t^{-1}u,
        \qquad
        G_{t,u}=
        \ip{\Sigma_t^{-1/2}u/\sqrt{a_{t,u}}}{Z_t}.
        \eqlabel{eq:H-standardisation}
    \end{align}
    Then
    \begin{align}
        \eps u^\top H_tu
        =a_{t,u}y_t(G_{t,u}^2-1).
        \eqlabel{eq:H-direction}
    \end{align}
    By covariance comparison and \cref{lem:predictable-spectrum},
    \begin{align}
        1\leq a_{t,u}\leq12dn^2,
        \qquad
        \kappa(A)\leq3dn^2.
        \ineqlabel{eq:concentration-spectrum}
    \end{align}
    These comparisons hold also at $t=\tau$, since the covariance
    stopping rule tests $\Sigma_{t+1}^{-1}$. To verify the bounds, use
    $\Sigma_t^{-1}\succeq\sigma^{-2}I_d/2\succeq I_d$,
    $\Sigma_t^{-1}\preceq3\bar{\Sigma}_t^{-1}/2$,
    $M\leq1$ (because $\K$ contains the unit ball), and
    \begin{align}
        \|\bar{\Sigma}_{t+1}^{-1}\|_{\op}
        \leq(\sigma^{-1}+n)^2
        \leq2(2d+n^2)\leq6dn^2,
        \qquad A\succeq2I_d.
    \end{align}

    Set $J=\lceil\log_2(12dn^2)\rceil$ and $b_j=2^j$ for
    $0\leq j\leq J$. For a fixed $(u,j)$, truncate the score in
    \cref{eq:H-direction} on $C_t$, center it, and multiply by
    the predictable indicator $I_t\1\{a_{t,u}\leq b_j\}$.
    Since $\norm{Z_t}^2+1\leq4d\ell$ on $C_t$, these centered increments
    have absolute value at most $32b_jd\ell^{5/2}$.
    By \cref{lem:gaussian-score-localisation} and
    \cref{eq:H-retained-moment-budget}, their total conditional
    variance is at most
    \begin{align}
        b_j^2\sum_{t=1}^n I_t
        \left(25\ell^2\E_t[y_t^2]+e^{-\ell/4}\right)
        \leq401b_j^2N\ell^2\leq426b_j^2N\ell^2.
        \ineqlabel{eq:H-directional-predictable}
    \end{align}
    Apply \cref{eq:freedman-convenient} with $z=d\ell/16$. The resulting
    bound is
    \begin{align}
        b_j\left(
        \sqrt{\frac{426}{8}}\sqrt{dN}\,\ell^{3/2}
        +4d^2\ell^{7/2}\right)
        \leq8b_j\sqrt{dN}\,\ell^{3/2}.
        \ineqlabel{eq:H-peeled-bound}
    \end{align}
    For the last inequality, divide the second term by
    $b_j\sqrt{dN}\ell^{3/2}$ and use
    \begin{align}
        \frac{4d^{3/2}\ell^2}{\sqrt N}
        =\frac{1}{25\sqrt d\,M\ell^6}
        \leq\frac{1}{25\ell^6}.
    \end{align}

    We must also restore the untruncated conditional means. On a
    retained round, \cref{lem:single-round-feedback} gives
    $\E_t[y_t^2]\leq8\ell^2$.
    Cauchy--Schwarz, \cref{lem:gaussian-score-localisation}, and
    $\Pp_t(\overline{C_t})\leq4e^{-\ell}$ then give the bias bound
    \begin{align}
        \left|\E_t[a_{t,u}y_t(G_{t,u}^2-1)\1_{\overline{C_t}}]\right|
        \leq40a_{t,u}\ell^2 e^{-\ell/2}.
    \end{align}
    Since $a_{t,u}\leq(3/2)u^\top Au$ on every retained round, the sum
    of these biases is at most
    \begin{align}
        60n\ell^2e^{-\ell/2}u^\top Au\leq n^{-8}u^\top Au.
    \end{align}
    Here $\eps\leq1/2$ implies
    $n\geq4\cdot10^4d^4M^2\ell^{16}\geq\ell$, and the last inequality
    follows from $\ell=128\log(128nd/\delta)$.

    If no round is retained, $T=0$ and the claim is immediate.
    Otherwise, on the dyadic level matching
    $a_u^\star=\max_{t:I_t=1}a_{t,u}$,
    no directional summand is removed and
    $b_j\leq2a_u^\star\leq3u^\top Au$. Hence, on
    $\mathcal E_{\mathrm{global}}$, the extra Gaussian event, and all
    the Freedman events used in \cref{eq:H-peeled-bound}, we obtain
    \begin{align}
        |u^\top Tu|
        \leq30\sqrt{dN}\,\ell^{3/2}u^\top Au
        \qquad(u\in\mathcal U).
        \ineqlabel{eq:H-net}
    \end{align}
    The dyadic level is selected only after applying the inequalities
    simultaneously to all deterministic levels.

    To pass from the net to all directions, write
    $r=(512dn^2)^{-1}$ and $c=30\sqrt{dN}\ell^{3/2}$.
    For any unit vector $w$, choose $u\in\mathcal U$ with
    $\norm{w-u}\leq r$. First \cref{eq:H-net} implies
    $\|T\|_{\op}\leq2c\|A\|_{\op}$, since $2r\leq1/2$.
    Applying the approximation once more gives
    \begin{align}
        |w^\top Tw|
         & \leq c\,w^\top Aw+6cr\|A\|_{\op}
        \leq c(1+6r\kappa(A))w^\top Aw
        \leq2c\,w^\top Aw.
    \end{align}
    Thus $-100\sqrt{dN}\ell^{3/2}A\preceq T
        \preceq100\sqrt{dN}\ell^{3/2}A$. Under the hypotheses of the
    lemma, the retained sum equals the sum through $\tau$ and
    $A=\bar{\Sigma}_{\tau+1}^{-1}$. Dividing by $\eps$ proves
    \cref{eq:H-concentration}.

    Define $\mathcal E_H$ as the intersection of the extra Gaussian
    event and all the Freedman events for $(u,j)\in\mathcal U\times
        \{0,\ldots,J\}$. By \cref{eq:H-net-size} and
    $\log(J+1)\leq\ell/128$, its failure probability is at most
    \begin{align}
        ne^{-\ell}+2|\mathcal U|(J+1)e^{-d\ell/16}
        \leq ne^{-\ell}+2e^{-5d\ell/128}.
    \end{align}
\end{proof}

\section{From interval averages to pointwise bounds}
\label{app:averages-to-pointwise}

\begin{proof}[Proof of \cref{lem:averages-to-pointwise}]
    Fix $1\leq u\leq n$. We start with the average over the most recent
    window and remove its earliest terms one at a time until only $a_u$
    remains. To track this change, write
    \begin{align}
        A_r=\frac1r\sum_{j=u-r+1}^{u}a_j,
        \qquad 1\leq r\leq\min\{k,u\}.
    \end{align}
    The window assumption \cref{eq:sequence-window-average} bounds
    $A_{\min\{k,u\}}$, while $A_1=a_u$ is the quantity we want to control.

    Removing $a_{u-r}$ from a window of length $r+1$ changes its average by
    \begin{align}
        A_r-A_{r+1}
        =\frac{1}{r(r+1)}
        \sum_{j=u-r+1}^{u}(a_j-a_{u-r}).
    \end{align}
    The sum is unchanged if we include its zero summand at $j=u-r$.
    We may therefore apply \cref{eq:sequence-local-average} to
    $[u-r,u]$, whose length is $r+1\leq k$, to obtain
    \begin{align}
        A_r-A_{r+1}
        \leq\frac{13(r+1)+2\sqrt{(r+1)\ell}}{r(r+1)}
        =\frac{13}{r}+\frac{2\sqrt\ell}{r\sqrt{r+1}}.
    \end{align}
    Telescoping these increases from the full window to its last term gives
    \begin{align}
        a_u
         & =A_{\min\{k,u\}}
        +\sum_{r=1}^{\min\{k,u\}-1}(A_r-A_{r+1})\nonumber \\
         & \leq14+2\sqrt{\frac{\ell}{\min\{k,u\}}}
        +13\sum_{r=1}^{\min\{k,u\}-1}\frac1r
        +2\sqrt\ell\sum_{r=1}^{\min\{k,u\}-1}\frac{1}{r\sqrt{r+1}}.
    \end{align}
    The two sums satisfy
    \begin{align}
        \sum_{r=1}^{\min\{k,u\}-1}\frac1r
         & \leq1+\log u,\nonumber                   \\
        \sum_{r=1}^{\min\{k,u\}-1}\frac{1}{r\sqrt{r+1}}
         & \leq\sum_{r=1}^{\infty}\frac{1}{r^{3/2}}
        \leq1+\int_1^\infty x^{-3/2}\,dx=3.
    \end{align}
    These bounds also hold when $\min\{k,u\}=1$, with empty sums equal to
    zero. Substituting them and using $\min\{k,u\}\geq1$ gives
    \begin{align}
        a_u
        \leq14+2\sqrt\ell+13(1+\log u)+6\sqrt\ell
        =27+13\log u+8\sqrt\ell.
    \end{align}
\end{proof}

\end{document}
